% concatenated from DisjunctiveSOS.tex, definitions.tex
% Dual-format toggle: preprint (current style) or Springer journal (Mathematical Programming).
\newif\ifpreprint
\preprinttrue   % <-- comment out this line for the journal (Springer) version

\ifpreprint
%%%%%%%%%%%%%%%%%%%%%%%%% PREPRINT %%%%%%%%%%%%%%%%%%%%%%%%%
\documentclass[12pt]{article}

% Prevent words to overflow over the margin
% \sloppy  % Alternative command
\emergencystretch 3em

\usepackage{fullpage}

% Hyperreferences (the journal version uses the hyperref loaded by sn-jnl)
\usepackage{hyperref}
\hypersetup{colorlinks=true, citecolor=cyan, linkcolor=cyan, urlcolor=cyan}

% Bibliography: biblatex/biber, IEEE style
\usepackage[
backend=biber,
style=ieee,
sortcites=false,
sorting=nyt
% mcitethresholds=1
% compress
]{biblatex}
\addbibresource{bibliography.bib}

% force the order of volume-page-year
\newtoggle{bbx:datemissing}

\renewbibmacro*{date}{\toggletrue{bbx:datemissing}%
}

\renewbibmacro*{issue+date}{%
  \toggletrue{bbx:datemissing}%
  \iffieldundef{issue}
    {}
    {\printtext[parens]{%
       \printfield{issue}}}%
  \newunit}

\newbibmacro*{date:print}{%
  \togglefalse{bbx:datemissing}%
  \printdate}

\renewbibmacro*{chapter+pages}{%
  \printfield{chapter}%
  \setunit{\bibpagespunct}%
  \printfield{pages}%
  \newunit
  \usebibmacro{date:print}%
  \newunit}

\renewbibmacro*{note+pages}{%
  \printfield{note}%
  \setunit{\bibpagespunct}%
  \printfield{pages}%
  \newunit
  \usebibmacro{date:print}%
  \newunit}

\renewbibmacro*{addendum+pubstate}{%
  \iftoggle{bbx:datemissing}
    {\usebibmacro{date:print}}
    {}%
  \printfield{addendum}%
  \newunit\newblock
  \printfield{pubstate}}

\else
%%%%%%%%%%%%%%% JOURNAL (Springer Mathematical Programming) %%%%%%%%%%%%%%%
\documentclass[pdflatex,sn-mathphys-num]{sn-jnl}
% sn-jnl loads hyperref, geometry, fontenc, and natbib, and sets \bibliographystyle internally.
\raggedbottom
\fi

%%%%%%%%%%%%%%%%%%%%%%%%% COMMON (both versions) %%%%%%%%%%%%%%%%%%%%%%%%%
% Images and plots
\usepackage{graphicx}
\usepackage{pgfplots}
% \usepackage{subfigure}
\usepackage{subcaption}
% Math
\usepackage{amsmath,amssymb,amsthm,enumitem,mathtools}
\renewcommand{\qedsymbol}{$\blacksquare$}
\renewcommand\arraystretch{1.2}
\usepackage{multirow}
\usepackage{algorithm}% http://ctan.org/pkg/algorithms
\usepackage[noend]{algpseudocode}

\renewcommand{\algorithmicrequire}{\textbf{Input:}}
\renewcommand{\algorithmicensure}{\textbf{Output:}}
\algnewcommand\algorithmicinitialization{\textbf{Initialization:}}
\algnewcommand\Initialization{\item[\algorithmicinitialization]}

% Tables
\usepackage{csvsimple}	% reading CSV files in tables
\usepackage{booktabs}   % Nicer Tables
\usepackage{adjustbox}  % To adjust table length
\usepackage{calc}

% Todos: no-ops by default; colored notes shown only in the preprint version
\newcommand{\bnote}[1]{}
\newcommand{\anote}[1]{}
\newcommand{\ynote}[1]{}
\newcommand{\snote}[1]{}
\ifpreprint
\renewcommand{\bnote}[1]{\textcolor{cyan}{\textbf{[B: #1]}}}
\renewcommand{\anote}[1]{\textcolor{magenta}{\textbf{[A: #1]}}}
\renewcommand{\ynote}[1]{\textcolor{black}{\textbf{[Y: #1]}}}
\renewcommand{\snote}[1]{\textcolor{brown}{\textbf{[S: #1]}}}
\fi

% Acronyms
\usepackage[xindy, acronyms]{glossaries}   % Glossaries and Acronyms
\makeglossaries

% Definitions

% Commands
\newcommand{\eg}{{\it e.g.}}
\newcommand{\ie}{{\it i.e.}}

% Sets
\newcommand{\ones}{\mathbf 1}
\newcommand{\reals}{\mathbb{R}}
\newcommand{\integers}{\mathbb{Z}}
\newcommand{\pintegers}{\mathbb{N}}
\newcommand{\symm}{{\mbox{\bf S}}}  % symmetric matrices
\newcommand{\NPhard}{\mbox{$\mathcal{NP}$-hard}}  % symmetric matrices
\newcommand{\prob}{{\mathbf P}}
\newcommand{\distrib}{{\mathcal P}}
\newcommand{\identity}{I}
\newcommand{\tpose}{T}
\newcommand{\nullspace}{{\mathcal N}}
\newcommand{\range}{{\mathcal R}}
\newcommand{\Rank}{\mathop{\bf Rank}}
\newcommand{\Tr}{\mathop{\bf Tr}}
\newcommand{\diag}{\mathop{\bf diag}}
\newcommand{\lambdamax}{{\lambda_{\rm max}}}
\newcommand{\lambdamin}{\lambda_{\rm min}}
\newcommand{\Expect}{\mathop{\bf E{}}}
\newcommand{\Prob}{\mathop{\bf Prob}}
\newcommand{\Co}{{\mathop {\bf Co}}} % convex hull
\newcommand{\dist}{\mathop{\bf dist{}}}
\newcommand{\argmin}{\mathop{\rm argmin}}
\newcommand{\argmax}{\mathop{\rm argmax}}
\newcommand{\epi}{\mathop{\bf epi}} % epigraph
\newcommand{\Vol}{\mathop{\bf vol}}
\newcommand{\dom}{\mathop{\bf dom}} % domain
\newcommand{\intr}{\mathop{\bf int}}

\newcommand{\sign}{\mathop{\bf sign}}
\newcommand{\round}{\mathop{\bf round}}
\newcommand{\define}{\coloneqq}

% new commands for this paper
\newcommand{\NN}{P}
\newcommand{\SOS}{\mathbf{\Sigma}}
\newcommand{\HNN}{\widehat{P}}
\newcommand{\HSOS}{\widehat{\mathbf{\Sigma}}}
\newcommand{\DiSOS}{\mathbf{\Sigma}}
\newcommand{\DiSOSL}{\text{NC}}
\newcommand{\defeq}{\coloneq} 

\newcommand{\supp}[1]{\operatornamewithlimits{supp}\left(#1\right)}
\newcommand{\vecp}{\operatornamewithlimits{vec}}
\newcommand{\maxi}{\operatornamewithlimits{maximize}}
\newcommand{\mini}{\operatornamewithlimits{minimize}}
\newcommand{\st}{\text{s.t.}}
\newcommand{\cone}{\operatornamewithlimits{cone}}
\newcommand{\conv}{\operatornamewithlimits{conv}}
\newcommand{\size}{\operatornamewithlimits{size}}
\newcommand{\cond}{\operatornamewithlimits{cond}}
\newcommand{\vol}{\operatornamewithlimits{vol}}
\newcommand{\Null}{\operatornamewithlimits{Null}}

% Trees
\newcommand{\tree}{\mathcal{T}}
\newcommand{\mcQ}{\mathcal{Q}}
\newcommand{\mcP}{\mathcal{P}}
\newcommand{\mcS}{\mathcal{S}}
\newcommand{\mcA}{\mathcal{A}}
\newcommand{\mcC}{\mathcal{C}}
\newcommand{\mcL}{\mathcal{L}}
\newcommand{\mcK}{\mathcal{K}}
\newcommand{\mcN}{\mathcal{N}}
\newcommand{\mbS}{\mathbb{S}}
\newcommand{\mcT}{\mathcal{T}}
\newcommand{\mcD}{\mathcal{D}}
\newcommand{\mcV}{\mathcal{V}}
\newcommand{\mcPN}{\mathcal{PN}}
\newcommand{\mbF}{\mathbb{F}}
\newcommand{\DD}{\text{DD}}
\newcommand{\SDD}{\text{SDD}}

% Strategy
\newcommand{\strategy}{\mathcal{S}}
\newcommand{\intvars}{\mathcal{I}}
\newcommand{\tightconstraints}{\mathcal{T}}

% Machine learning
\newcommand{\loss}{\mathcal{L}}

% Sections (for citation)
\newcommand{\Sec}{Section\;}

% Create theorems and other environments
\newtheorem{theorem}{Theorem}%[section]  % Restart counter every section
\newtheorem{lemma}[theorem]{Lemma}  % Restart counter every section
% \newtheorem{corollary}{Corollary}[section]
\newtheorem{corollary}{Corollary}[theorem]  % Restart counter every theorem
\newtheorem{proposition}{Proposition}%[section]
\newtheorem{assumption}{Assumption}%[section]
\newtheorem{definition}{Definition}%[section]
\newtheorem{example}{Example}%[section]
\newtheorem{exercise}{Exercise}%[section]
\theoremstyle{remark}
% \newtheorem*{remark}{Remark}
\newtheorem{remark}{Remark}%[section]
\newtheorem{problem}{Problem}[section]

% Acronyms
\newacronym{LO}{LO}{linear optimization}
\newacronym{QO}{QO}{quadratic optimization}
\newacronym{MIQO}{MIQO}{mixed-integer quadratic optimization}
\newacronym{MIO}{MIO}{mixed-integer optimization}
\newacronym{MILO}{MILO}{mixed-integer linear optimization}
\newacronym{MINLO}{MINLO}{mixed-integer nonlinear optimization}
\newacronym{sBB}{sBB}{spacial branch and bound}
\newacronym{NLO}{NLO}{nonlinear optimization}
\newacronym{PWA}{PWA}{piecewise affine}
\newacronym{SVM}{SVM}{support vector machines}
\newacronym{ReLU}{ReLU}{rectified linear unit}
\newacronym{CPU}{CPU}{central processing unit}
\newacronym{GPU}{GPU}{graphics processing unit}
\newacronym{MPC}{MPC}{model predictive control}
\newacronym{ADMM}{ADMM}{alternating direction method of multipliers}
\newacronym{ADP}{ADP}{approximate dynamic programming}
\newacronym{FPGA}{FPGA}{field-programmable gate array}

% breakable algorithm
\makeatletter
\newenvironment{breakablealgorithm}
  {% \begin{breakablealgorithm}
   \begin{center}
     \refstepcounter{algorithm}% New algorithm
     \hrule height.8pt depth0pt \kern2pt% \@fs@pre for \@fs@ruled
     \renewcommand{\caption}[2][\relax]{% Make a new \caption
       {\raggedright\textbf{\fname@algorithm~\thealgorithm} ##2\par}%
       \ifx\relax##1\relax % #1 is \relax
         \addcontentsline{loa}{algorithm}{\protect\numberline{\thealgorithm}##2}%
       \else % #1 is not \relax
         \addcontentsline{loa}{algorithm}{\protect\numberline{\thealgorithm}##1}%
       \fi
       \kern2pt\hrule\kern2pt
     }
  }{% \end{breakablealgorithm}
     \kern2pt\hrule\relax% \@fs@post for \@fs@ruled
   \end{center}
  }
\makeatother

% insert two subfigures without captions
\newcommand{\includesubfig}[3][0.4]{
    \subfigure[#3]{
      \centering
      \includegraphics[width=#1\textwidth]{#2}  
    }
}




% Abstract text shared by both versions
\newcommand{\abstracttext}{We introduce the concept of \emph{disjunctive sum of squares} for certifying nonnegativity of polynomials. Unlike the popular sum of squares approach where nonnegativity is certified by a single algebraic identity, the disjunctive sum of squares approach certifies nonnegativity with multiple algebraic identities which can be found in parallel. Our main result is a disjunctive Positivstellensatz proving that we can keep the degree of each algebraic identity as low as the degree of the polynomial whose nonnegativity is in question. Based on this result, we construct a semidefinite programming based converging hierarchy of lower bounds for the problem of minimizing a polynomial over a compact basic semialgebraic set, where the size of the largest semidefinite constraint is fixed throughout the hierarchy. We further prove a second disjunctive Positivstellensatz which leads to an optimization-free hierarchy for polynomial optimization. We specialize this result to the problem of proving copositivity of matrices. Finally, we describe how the disjunctive sum of squares approach can be combined with a branch-and-bound algorithm and we present numerical experiments on polynomial, copositive, and combinatorial optimization problems.}

\begin{document}

\ifpreprint
%%%%%%%%%%%%%%%%%%%%%%%%% PREPRINT FRONT MATTER %%%%%%%%%%%%%%%%%%%%%%%%%
\title{ Disjunctive Sum of Squares }
% \author{Amir Ali Ahmadi, Sanjeeb Dash, Yixuan Hua, Bartolomeo Stellato}
\author{Amir Ali Ahmadi\thanks{Princeton University, Operations Research and Financial Engineering.
AAA was partially supported by the Sloan Fellowship, the Princeton AI Lab Seed Grant, the Princeton SEAS Innovation Grant, and a Research Gift in Mathematical Optimization.
YH was partially supported by the Princeton AI Lab Seed Grant, the Princeton SEAS Innovation Grant, and the IBM PhD Fellowship.},
 Sanjeeb Dash\thanks{IBM TJ Watson Research Center.},
Yixuan Hua\footnotemark[1],
Bartolomeo Stellato\footnotemark[1]}
\date{}
\maketitle

\begin{abstract}
\abstracttext
\end{abstract}
\else
%%%%%%%%%%%%%%%%%%%%%%%%% JOURNAL FRONT MATTER %%%%%%%%%%%%%%%%%%%%%%%%%
\title[Disjunctive Sum of Squares]{Disjunctive Sum of Squares}
\author*[1]{\fnm{Amir Ali} \sur{Ahmadi}}\email{aaa@princeton.edu}
\author[2]{\fnm{Sanjeeb} \sur{Dash}}\email{sanjeebd@us.ibm.com}
\author[1]{\fnm{Yixuan} \sur{Hua}}\email{yh7422@princeton.edu}
\author[1]{\fnm{Bartolomeo} \sur{Stellato}}\email{bstellato@princeton.edu}
\affil[1]{\orgdiv{Department of Operations Research and Financial Engineering}, \orgname{Princeton University}, \orgaddress{\city{Princeton}, \state{NJ}, \country{USA}}}
\affil[2]{\orgdiv{Mathematics of Computation Department}, \orgname{IBM TJ Watson Research Center}, \orgaddress{\city{Yorktown Heights}, \state{NY}, \country{USA}}}
\abstract{\abstracttext}
% TODO: complete keywords for the journal submission
\keywords{TODO, TODO, TODO}
% TODO: complete the MSC classification codes
\pacs[MSC Classification]{TODO}
\maketitle
\fi

% Table of contents
%\newpage \tableofcontents \newpage


\section{Introduction}
A polynomial $p:\mathbb{R}^n\rightarrow\mathbb{R}$ is said to be \emph{nonnegative} if $p(x)\ge 0$ for all $x\in\reals^n$. 
%We denote the set of nonnegative polynomials in $n$ variables and degree $d$ as $\NN_{n,d}$.
The problem of optimizing over the intersection of the cone of nonnegative polynomials of a given degree with an affine subspace %Linear optimization over the intersection of an affine subspace with $\NN_{n,d}$ 
is fundamental to optimization and algebraic geometry, with broad applications in polynomial optimization~\cite{lasserre2001global}, control systems and robotics~\cite{henrion2005positive,tedrake2010lqr}, probability theory~\cite{bertsimas2005optimal}, and statistics~\cite{curmei2023shape}, to name a few; see the textbooks~\cite{lasserre2009moments,blekherman2012semidefinite,lasserre2015introduction}, or survey papers~\cite{laurent2009sums,hall2020applications} for many other applications and further context.  

%Closely related to nonnegative polynomials are sum of squares polynomials. 
A polynomial $p:\mathbb{R}^n\rightarrow\mathbb{R}$ is said to be a \emph{sum of squares} (sos) if it can be expressed as
$p(x)=\sum_i q_i^2(x)$ for some polynomials $q_i:\mathbb{R}^n\rightarrow\mathbb{R}$. 
% We denote the set of sos polynomials in $n$ variables and degree $d$ as $\SOS_{n,d}$. 
A sum of squares decomposition of a polynomial provides a single algebraic identity that certifies its nonnegativity.
%Since every sos polynomial is nonnegative, the sos representation provides a single algebraic identity to certify nonnegativity of polynomials.
Moreover, as is well known by now, such a decomposition (assuming it exists) can be obtained by solving a semidefinite program (SDP)~\cite{Parrilo2000,lasserre2001global,parrilo2003semidefinite}.

%However, not every nonnegative polynomial is sos. In a seminal paper~\cite{Hilbert1888}, 

Throughout the paper, we denote the set of nonnegative polynomials and sos polynomials in $n$ variables and degree at most $d$ by $\NN_{n,d}$ and $\SOS_{n,d}$, respectively. 
It has been known since the seminal work of Hilbert~\cite{Hilbert1888} that not every nonnegative polynomial is a sum of squares. In fact, Hilbert established that $\NN_{n,d}$ and $\SOS_{n,d}$ coincide only in three specific cases: when $n=1$, $d=2$, or $(n,d)=(2,4)$. The first example of a nonnegative polynomial which is not sos was constructed by Motzkin many years later~\cite{Motzkin1967}:
\begin{equation}\label{eq:motzkin_nonhomo}
M(x_1,x_2)=x_1^4x_2^2+x_1^2x_2^4-3x_1^2x_2^2+1.
\end{equation}

%Can one still prove the nonnegativity of a non-sos polynomial with a single algebraic identity?

Despite the existence of such examples, it turns out that a proof of nonnegativity based on a single algebraic identity is always possible though at the expense of increasing the degree of the sos polynomials that appear in the identity. More specifically, it follows from Artin's solution to Hilbert's 17th problem~\cite{Artin1927} that for every nonnegative polynomial $p$, there exists a nonzero sos polynomial $q$, such that the product $pq$ is sos (and hence nonnegativity of $p$ is algebraically certified by a representation as a ratio of two sos polynomials). For example, for the Motzkin polynomial, we have the following degree-8 identity:
\begin{equation}\label{eq:Motzkin.with.multiplier}
(x_1^2+x_2^2)M(x_1,x_2)=x_1^2(1-x_2^2)^2+x_2^2(1-x_1^2)^2+x_1^2x_2^2(x_1^2+x_2^2-2)^2.    
\end{equation}
By utilizing Artin's theorem, one can design a complete SDP-based hierarchy for proving nonnegativity of a polynomial $p$. In level $j$ of this hierarchy, one searches for a nonzero sos polynomial~$q$ of degree~$2j$ that makes the product $pq$ also sos.\footnote{Unlike the approach we later propose in this paper, one cannot use this SDP hierarchy to \emph{impose} a nonnegativity constraint on a polynomial $p$. This is because of the bilinear matrix inequality that arises from the sos constraint on the product of the two unknown polynomials $p$ and $q$.}
However, the degree of the polynomial $pq$ can be much higher than the degree of~$p$ (see, e.g.,~\cite{sos_on_hypercube}) and is upper bounded currently only by a tower of five exponentials in the number of variables and degree of $p$~\cite{lombardi2020elementary}. Since the size of the SDP underlying an sos %a sum of squares 
constraint grows rapidly with the degree of the 
sos polynomial, %sum of squares polynomial
the SDP hierarchy based on Artin's theorem can quickly become impractical in certain applications. This challenge also arises in the context of other sum of squares based SDP hierarchies (e.g., in~\cite{Parrilo2000,lasserre2001global,parrilo2003semidefinite}) which rely on related seminal theorems of algebraic geometry such as the
% Positivstellens\"atze of Stengle *, Schm\"udgen *, or Putinar *.
Positivstellens\"atze of Stengle~\cite{stengle1974nullstellensatz}, Schm\"udgen~\cite{Schmudgen1991}, or Putinar~\cite{putinar1993positive}.

With this context in mind, the central question that motivates our research is the following: \emph{``Can one lower the complexity of a sum of squares based proof of polynomial nonnegativity by constructing several algebraic identities instead of a single one?''} Our primary notion of complexity here is the degree of the sos polynomials that appear in the proof of nonnegativity. However, our results also simplify the structure of these sos polynomials making the search for them amenable to cheaper classes of conic programs, such as linear programs. As an example, consider the following proof of nonnegativity of the Motzkin polynomial using two degree-6 algebraic identities:
\begin{equation}\label{eq:Motzkin.disos.proof}
\begin{aligned}
M(x_1,x_2) &= (1-x_1x_2)^2 + (x_1^2x_2-x_1x_2^2)^2 + 2x_1x_2(1-x_1x_2)^2\\
    & = (1+x_1x_2)^2 + (x_1^2x_2+x_1x_2^2)^2 - 2x_1x_2(1+x_1x_2)^2. 
    \end{aligned}    
\end{equation}
%The sos method, based on the idea of finding a multiplier guaranteed by Artin's theorem, yields a complete hierarchy to find an algebraic identity certifying nonegativity. On each level of the hierarchy, the optimization problem is an SDP. However, the multiplier arising from the hierarchy can have a high degree~\cite{lombardi2014elementary}, making the SDP computationally expensive.
%Can one prove nonnegativity with low-degree identities to avoid solving large SDPs? We propose a disjunctive sum of squares (disos) method to give a positive answer to this question. We use multiple algebraic identities with exactly the same degree as $p$, to certify its nonnegativity. For example, the Motzkin polynomial can be written as the following two identities, which we collectively call a \emph{disjunctive certificate}. 
% \begin{equation}\label{eq:Motzkin.disos}
% \begin{aligned}
% M(x_1,x_2) &= (1-x_1x_2)^2 + (x_1^2x_2-x_1x_2^2)^2 + 2x_1x_2(1-x_1x_2)^2\\
%     & = (1+x_1x_2)^2 + (x_1^2x_2+x_1x_2^2)^2 - 2x_1x_2(1+x_1x_2)^2. 
%     \end{aligned}    
% \end{equation}
Indeed, for every point $(x_1,x_2)\in\reals^2$, either $x_1x_2\ge 0$ or $x_1x_2<0$. When $x_1x_2\ge 0$, nonnegativity of $M(x_1,x_2)$ follows from the first identity. Similarly, when $x_1x_2<0$ or equivalently $-x_1x_2>0$, nonnegativity of $M(x_1,x_2)$ follows from the second identity. We call a proof of this type a \emph{disjunctive sum of squares} proof of nonnegativity. Comparing the two algebraic proofs in (\ref{eq:Motzkin.with.multiplier}) and (\ref{eq:Motzkin.disos.proof}), we see that the former involves an sos polynomial of degree 8 while the latter involves two sos polynomials of degree 6, which is the same as the degree of the Motzkin polynomial. From a computational standpoint, this means that the size of the semidefinite constraint underlying the second proof is smaller.\footnote{We recall that, imposing an sos constraint on a polynomial in $n$ variables and degree $2d$ results in a positive semidefinite constraint on a matrix of size ${{n+d}\choose{d}}\times {{n+d}\choose{d}}$ (see, e.g.,~\cite{Parrilo2000, parrilo2003semidefinite}).}
% Furthermore, reducing the constraint size may lead to a greater reduction in the overall complexity, even if more constraints are introduced~\cite{vandenberghe1996semidefinite}. } 
Moreover, the sos polynomials appearing in (\ref{eq:Motzkin.disos.proof}) have a special structure (known as ``diagonal dominance''; see~\cite{aaa2018dsos}), that enables the search for this particular disjunctive sum of squares proof to be carried out by solving a linear program. This is, in general, much cheaper than solving a semidefinite program.


%This implies that the search for this particular disjunctive sum of squares proof can be carried out by solving a linear program which in general can be much cheaper than solving a semidefinite program.

Our main contribution in this paper is to formalize the concept of disjunctive sum of squares proofs of nonnegativity and show that \emph{low-degree} proofs of this type always exist under standard assumptions. We outline our contributions in some more detail next.

% Our main contribution in this paper is to formalize the concept of disjunctive sum of squares proofs of nonnegativity and show that low-degree proofs of this type exist for a large class of polynomials. We outline our contributions in some more detail next.


%The disjunctive certificate can be found by solving multiple small SDPs. Moreover, in some cases, e.g., the Motzkin polynomial, the sos terms have simple structures and can be found by solving linear programs~\cite{aaa2018dsos}.

\subsection{Organization and Main Contributions}
The remainder of this paper is organized as follows. In Section~\ref{sec:disos}, we define the notion of a disjunctive sum of squares proof of nonnegativity of a polynomial. 
% We also describe an SDP-based approach to searching for such a proof.
% We also show that there are three different disjunctive proof of nonnegativity
We also show that a natural alternating maximization approach based on SDP can find {low-degree disjunctive sum of squares proofs of nonnegativity of several classical non-sos polynomials}.


In Section~\ref{sec:disos_psatz}, we prove a ``disjunctive Positivstellensatz'', which shows that the degree of the sos polynomials appearing in a disjunctive sum of squares proof can be bounded above by the degree of the polynomial whose nonnegativity is being certified. Sections~\ref{sec:local_psatz} and \ref{sec:global_psatz} are devoted to the proof of this statement. Based on this result, Section~\ref{sec:disos_grid} introduces an SDP-based hierarchy, with semidefinite constraints of \emph{fixed size}, which yields 
a converging sequence of 
lower bounds on the minimum value of a homogeneous polynomial over the unit sphere. Our analysis also yields a convergence rate for this hierarchy. 

% arbitrarily tight lower bounds on the minimum value of a homogeneous polynomial over the unit sphere. 

%-a polynomial whose top homogeneous component is positive definite (true)

% Suppose p(x)>0 for all x in R^n,
% then
% ph(x,y)=y^d*p(x/y) is positive definite (as a form in x,y)
% iff
% the thc(p) is positive definite (as a form in x).

% Suppose p(x)>0 for all x in R^n
% ==>thc(p) is nonnegative. 

% slightly stronger assumption:
% thc(p) is positive definite (as a form in x).
% form=homegneous polynomial
% f(x) homogeneous
% f(x) positive means f(x)>0 for all x nonzero 
% iff
% f(x)>0 for all x with ||x||=1.

% f(alpha*x)=alpha^d*f(x) 




%- of a homogeneous polynomial on the unit sphere.


%-? a polynomial which is coercive (false)
%-?a polynomial which attains its infimum (false)





% p(x,y)= x^2+(1-x*y)^2  >0 for all x,y nonnegative
% x-->0, y=1/x
% p(x,y)-->0
% ph(x,y,z)=x^2*z^2+(z^2-x*y)^2
% ph(1,0,0)=0 this homogeneous polynomial is NOT positive definite.

% q(x,y)= x^2+y^2+(1-x*y)^2 
% >=2xy + (1-xy)^2= 1+ (xy)^2 >=1
% (x,y)=(0,0)
% qh(x,y,z)=x^2*z^2+y^2*z^2+(z^2-x*y)^2
% qh(1,0,0)=0

% homog nonnegative



% We show that, for any fixed dimension and degree, there exists a family of partitioning polynomials such that any positive definite form can be certified with a disjunctive sum of squares proof based on certain polynomials in the family. Building on this result, we design a uniform gridding algorithm in Section~\ref{sec:disos_grid}, which provides a complete SDP-based hierarchy to certify global nonnegativity of polynomials. At each level of the hierarchy, the algorithm solves multiple small, fixed-size SDPs, where the number of SDPs equals the number of subregions created by the partition.

In Section~\ref{sec:disos_new}, we prove a second disjunctive Positivstellensatz, which shows that a proof of nonnegativity of a polynomial can be given by simply checking nonnegativity of the coefficients of the same polynomial under several linear changes of coordinates. These linear transformations are explicit and derived from disjunctions based on simplicial partitions of the Euclidean space. An upper bound on the number of linear transformations is provided.



% present an optimization-free approach to certify nonnegativity of polynomials. This approach uses linear polynomials to partition the Euclidean space and a different proof of nonnegativity. We show that one can always certify nonnegativity of a polynomial by examining the coefficients of its transformed version under some suitable linear mappings.

% Additionally, we establish a disjunctive Positivstellensatz tailored to this setting, analogous to the result introduced in Section~\ref{sec:disos_psatz}.

In Section~\ref{sec:constrained_opt}, we demonstrate how our disjunctive sum of squares approach can be applied to constrained polynomial optimization problems. 
In Section~\ref{sec:hierarchy}, 
using a reduction from~\cite{aaa2019}, we construct a semidefinite programming based converging hierarchy of lower bounds for the general problem of minimizing a polynomial over a compact basic semialgebraic set, where the size of the largest semidefinite constraint is fixed throughout the hierarchy.
% we create a converging hierarchy of optimization problems that produces arbitrarily tight lower bound for general polynomial optimization problems with a compact feasible set. 
In Section~\ref{sec:copositive}, 
we focus on the case where the constraint set is the unit simplex, and we provide a specialized disjunctive Positivstellensatz for certifying copositivity of matrices.

% we refine our approach to certify copositivity of matrices through a modified hierarchy.
% To address efficiency issues of using reductions in the hierarchy,
% we adapt our approach to directly certify copositivity of matrices in Section~\ref{sec:copositive}.

% This is based on a reduction from constrained optimization problems to unconstrained ones developed in previous work. While the reduction is theoretically sound, it is not always efficient in practice. To address this,
% in Section~\ref{sec:copositive}, we adapt the disjunctive sum of squares method to directly certify copositivity of matrices without such reduction.

In Section~\ref{sec:bnb}, we show how one can more practically search for disjunctive sum of squares proofs by dynamically partitioning the Euclidean space through a 
spatial branch-and-bound framework. As examples, 
% more practical disjunctions of the Euclidean space can be obtained dynamically by 
% the spatial branch-and-bound framework.
%
we propose branch-and-bound algorithms for minimizing a homogeneous polynomial over the unit sphere and for certifying copositivity of matrices. In Section~\ref{sec:num}, we present some numerical experiments. 
% illustrating the algorithms' behavior across a range of optimization tasks. 
% including certifying nonnegativity of non-sos polynomials, solving standard quadratic programmings via the copositive cone, and computing the clique number of graphs. 
Finally, Section~\ref{sec:discussion} concludes our paper and highlights potential future directions.



\section{SOS with Respect to an Algebraic Disjunction}\label{sec:disos}
In this section, we introduce the concept of sum of squares with respect to an algebraic disjunction, which is central to our framework. 
% which provides an algebraic certificate of nonnegativity of polynomials. 
Throughout the paper, we denote the degree of a polynomial $p$ by $\deg p$.
%sufficient condition for nonnegativity of polynomials and leads to a disjunctive sum of squares proof. 

% As odd-degree polynomials cannot be nonnegative on the entire Euclidean space, we focus only on even-degree polynomials. Further, it suffices to focus on even-degree forms (i.e. homogeneous polynomials). Indeed, for a polynomial $p$ in $n$ variables and degree $d$, its homogenized version  $p_h(x,x_{n+1})\defeq x_{n+1}^dp(x/x_{n+1})$ is a homogeneous polynomial in $n+1$ variables and degree $d$ which is nonnegative if and only if $p$ is nonnegative. 


% For example, the homogeneous Motzkin polynomial is
% \begin{equation}\label{eq:homo_Motzkin}
%     M_h(x_1,x_2,x_3)=x_1^4x_2^2+x_1^2x_2^4-3x_1^2x_2^2x_3^2+x_3^6.
% \end{equation}
% Note that homogenization preserves nonnegativity, i.e., $p$ being nonnegative is equivalent to $p_h$ being nonnegative. 
% Thus, unless otherwise indicated, all polynomials 
% are assumed to be homogeneous for the rest of the paper. 

% Throughout the paper, we denote the degree of a polynomial $p$ by $\deg p$, and the set of forms, nonnegative forms, and sum of squares forms in $n$ variables and degree $d$ by $H_{n,d},\HNN_{n,d},\HSOS_{n,d}$,  respectively. 
% % $H_{n,d},\HNN_{n,d},\HSOS_{n,d}$ as the set of homogeneous polynomials (i.e., forms), nonnegative polynomials, and sum of squares polynomials in $n$ variables and degree $d$, respectively. 
% Observe that 
% $\HSOS_{n,d}\subseteq \HNN_{n,d}$, though the inclusion is strict in general. Definition~\ref{def:disos} below introduces families of polynomials of interest to us that lie in between $\HSOS_{n,d}$ and $\HNN_{n,d}$.

%and these two sets do not coincide in general. Therefore, we aim to construct a larger inner approximation of $\HNN_{n,d}$.

%Motivated by the algebraic identities in~\eqref{eq:Motzkin.disos.proof}, proving nonnegativity of polynomials over specific semialgebraic sets may be more tractable. To this end, we consider a collection of basic closed semialgebraic subsets of $\mathbb{R}^n$ such that their union is $\mathbb{R}^n$. We refer to the polynomials defining these semialgebraic subsets as an algebraic disjunction of $\mathbb{R}^n$.

\begin{definition}[Algebraic Disjunction]\label{def:disjunction}
Let $r,n_1,\ldots,n_r$ be positive integers. We say that a collection of polynomials $\mcD=\{\{q_{k,j}\}_{j=1,\ldots, n_k},\quad k=1,\ldots,r\}$ forms an \emph{algebraic disjunction} of $\reals^n$ if 
\[\mathbb{R}^n=\bigcup\limits_{k=1}^r\Omega_{k},\ \text{where }\Omega_k=\{x\in\mathbb{R}^n\mid q_{k,j}(x)\ge 0,\quad j=1,\ldots,n_k\}.\]
\end{definition}
% We next introduce the concept of a polynomial being a sum of squares with respect to an algebraic disjunction. A polynomial with this property means that a polynomial has an algebraic proof of nonnegativity over each basic closed semialgebraic set in the disjunction.
%L $\deg p$ as the degree of a polynomial $p$.
\begin{definition}\label{def:disos}
Let $p$ be a polynomial in $n$ variables and degree at most $d$, $\bar{d}\ge d$ be an integer, and $\mcD=\{\{q_{k,j}\}_{j=1,\ldots, n_k},\quad k=1,\ldots,r\}$ be an algebraic disjunction of $\reals^n$.
We say that $p$ is \emph{sos of order $\bar{d}$ with respect to }$\mcD$ \emph{(or $p$ is sos$^{\bar{d}}_\mathcal{D}$)} if there exist sos polynomials $s_{k,j}$ for $k=1,\ldots, r$ and $j=0,1,\ldots,n_k$, such that 
\begin{equation}\label{def:sos_ad}
    p(x)=s_{k,0}(x)+\sum_{j=1}^{n_k} s_{k,j}(x)q_{k,j}(x),\quad  k=1,\ldots, r,
\end{equation}
where $\deg s_{k,0}\le \bar{d}$ and $\deg(s_{k,j}q_{k,j})\le \bar{d}$ for $k=1,\ldots,r$ and  $j=1,\ldots,n_k$. 
We denote the set of sos$^{\bar{d}}_\mcD$ polynomials in $n$ variables and degree at most $d$ as $\DiSOS_{n,d}^{\bar{d}}(\mcD)$.
\end{definition}
% \begin{remark}
% Observe that the set $\DiSOS_{n,d}^{d_t}(\mcD)$ is a convex cone. Moreover, given polynomials $q_{k,j}$ for $j=1,\ldots,n_k$, the search of $s_{k,j}$ can be done by solving an SDP problem. In implementation, we do not include $s_{k,0}$ as a decision variable; instead, we impose the constraint that $p(x)-\sum_{j=1}^{n_k} s_{k,j}(x)q_{k,j}(x)$ is sos.
% \end{remark}


Observe that for any algebraic disjunction $\mcD$ and integer $\bar{d}\ge d$, the set $\DiSOS^{\bar{d}}_{n,d}(\mcD)$ is a convex cone and lies between $\SOS_{n,d}$ and $\NN_{n,d}$. 
Indeed, by setting $s_{k,j}=0$ for $k=1,\ldots,r$ and $j=1,\ldots,n_k$, we note that $\SOS_{n,d}\subseteq\DiSOS^{\bar{d}}_{n,d}(\mcD)$. This inclusion can be strict even when $\bar{d}=d$; e.g., the Motzkin polynomial in~\eqref{eq:motzkin_nonhomo} belongs to $\DiSOS_{2,6}^{6}(\{ \{+x_1x_2\},\{-x_1x_2\}\})$, as shown by the identities in~\eqref{eq:Motzkin.disos.proof}. 
To see that $\DiSOS^{\bar{d}}_{n,d}(\mcD)\subseteq \NN_{n,d}$, note that for any point $x\in\Omega_k$, we have $q_{k,j}(x)\ge 0$ for all $j=1,\ldots,n_k$, and therefore, in view of~\eqref{def:sos_ad},~$p(x)\ge 0$. 
% Applying identity~\eqref{def:sos_ad}, this implies that $p(x)\ge 0$ for all $x\in\Omega_k$. 
Consequently, $p$ is nonnegative over each subregion $\Omega_k$, and thus nonnegative over~$\reals^n$. We call the $r$ algebraic identities in~\eqref{def:sos_ad} a disjunctive sos proof of nonnegativity of~$p$. % Moreover, given an algebraic disjunction $\mcD$ and integer $\bar{d}\ge d$, 

From a computational standpoint, for any fixed $\bar{d}$, linear optimization over the intersection of an affine subspace with the cone $\DiSOS_{n,d}^{\bar{d}}(\mcD)$ can be carried out by an SDP of size polynomial in $n$. We also note that for implementation purposes, there is no need to include $s_{k,0}$ as a decision variable; instead, one can impose the constraint that~${p(x)-\sum_{j=1}^{n_k} s_{k,j}(x)q_{k,j}(x)\in\SOS_{n,\bar{d}}}$ for $k=1,\ldots,r$. 
% Moreover, this constraint can be imposed as $r$ independent SDPs, one for each $k\in\{1,\ldots,r\}$, that can be solved in parallel.

% For any polynomial $p\in \DiSOS_{n,d}^{d_t}(\mcD)$, the order-$d_t$ disjunctive sos proof of nonnegativity of $p$ shows the implication that
% \[q_{k,j}(x)\ge 0,\quad j=1,\ldots,n_k\ \Rightarrow\ p(x)\ge 0,\quad k=1,\ldots, r.\]
% Consequently, $p$ is nonnegative over each subregion $\Omega_k$, and thus nonnegative over $\reals^n$. This implies that $\DiSOS_{n,d}^{d_t}(\mcD)$ is an inner approximation of $\HNN_{n,d}$.

Unlike the well-known hierarchies for proving nonnegativity of a polynomial (e.g., based on the Positivstellens\"atze of Artin~\cite{Artin1927}, Reznick~\cite{reznick1995uniform}, Stengle~\cite{stengle1974nullstellensatz, Parrilo2000}, or Putinar~\cite{putinar1993positive, lasserre2001global}) that increase the degrees of the sos polynomials involved in the proof at each level of the hierarchy, our approach keeps the degrees of these sos polynomials fixed (and low). 
% Unlike much of the literature on sum of squares optimization that provides hierarchies of sufficient conditions for nonnegativity of a polynomial by increasing the degrees of the sos polynomials involved in the proof (e.g., the hierarchies based on Positivstellens\"atze of Artin~\cite{Artin1927}, Reznick~\cite{reznick1995uniform}, Stengle~\cite{stengle1974nullstellensatz, Parrilo2000}, or Putinar~\cite{putinar1993positive, lasserre2001global}),  
%the Lasserre's hierarchy~\cite{lasserre2001global}), 
% our approach in this paper is to keep the degrees of the sos polynomials involved fixed (and low), and 
We instead increase the number of subsets in the algebraic disjunction $\mcD$ (see Theorem~\ref{cor:strictly_pos} and Theorem~\ref{cor:strictly_pos_new}). 
% Unlike the Lasserre's hierarchy, which involves increasing the order $d_t$ to improve the representability of the algebraic identity, we focus on keeping the order $d_t$ the same as the degree of polynomial whose nonnegativity is being certified. Note that the set $\DiSOS_{n,d}^d(\mcD)$ is at least as large as $\HSOS_{n,d}$ for any algebraic disjunction $\mcD$, and may contain non-sos polynomials. if $\mcD$ is carefully chosen. 
% Indeed, even when $\bar{d}=d$, the cone $\DiSOS_{n,d}^{\bar{d}}(\mcD)$ can include non-sos polynomials for some choices of $\mcD$. For example, the Motzkin polynomial~\eqref{eq:motzkin_nonhomo} belongs to $\DiSOS_{2,6}^{6}(\{ \{+x_1x_2\},\{-x_1x_2\}\})$, as seen by the identities in~\eqref{eq:Motzkin.disos.proof}.
% For example, the homogeneous Motzkin polynomial 
% \begin{equation}\label{eq:homo_Motzkin}
%     M_h(x_1,x_2,x_3)=x_1^4x_2^2+x_1^2x_2^4-3x_1^2x_2^2x_3^2+x_3^6,
% \end{equation}
% satisfies $M_h\in \DiSOS_{3,6}^{6}(\{ \{+x_1x_2\},\{-x_1x_2\}\})$, as seen by the homogenized version of~\eqref{eq:Motzkin.disos.proof}:
% satisfies $M_h\in \DiSOS_{6,3}(\{ \{+x_1x_2\},\{-x_1x_2\}\})$ as seen by the following algebraic identities:
% \begin{equation}\label{eq:Motzkin_disos}
%  \begin{aligned}
%  M_h 
%    &= (x_3^3-x_1x_2x_3)^2 + (x_1^2x_2-x_1x_2^2)^2 + 2x_1x_2(x_3^2-x_1x_2)^2\\
%     & = (x_3^3+x_1x_2x_3)^2 + (x_1^2x_2+x_1x_2^2)^2 - 2x_1x_2(x_3^2+x_1x_2)^2.
% \end{aligned} 
% \end{equation}

% Clearly,~\eqref{eq:Motzkin_disos} is the homogenized version of~\eqref{eq:Motzkin.disos.proof}. 
% Additionally, the above observation justifies our focus on homogeneous polynomials. Given a polynomial $p$ in $n$ variables and degree $d$, $p\in\DiSOS_{n,d}^d(\mcD)$ is equivalent to its homogeneization $p_h\in\DiSOS_{n+1,d}^d(\widehat{\mcD})$, where the polynomials in $\widehat{\mcD}$ are the homogenization of the polynomials in $\mcD$. Indeed, if $p$ is sos$^d_\mcD$, it can be expressed as in equation~\eqref{def:sos_ad}, then homogenizing all polynomials in~\eqref{def:sos_ad} implies that $p_h$ is sos$^d_{\widehat{\mcD}}$.
% On the other hand, if $p_h$ is sos$^d_{\widehat{\mcD}}$, 
% setting $x_{n+1}=1$ in the disjunctive sos proof of nonnegativity of $p_h$ implies that $p$ is sos$^d_{\mcD}$.
% it can be expressed as 
% \[    p_h(x)=\hat{s}_{k,0}(x)+\sum_{j=1}^{n_k} \hat{s}_{k,j}(x)\hat{q}_{k,j}(x),\quad  k=1,\ldots, r,\]
% where $\hat{s}_{k,0}$ is sos and $\deg \hat{s}_{k,0}\le d$, $\hat{s}_{k,j}$ is sos and $\deg(\hat{s}_{k,j}\hat{q}_{k,j})\le d$ for $j=1,\ldots,n_k$.
% Setting the last element of $x$ to 1 in the above equation implies that $p$ is sos$^d_{\mcD}$.
% Given a polynomial $p\in H_{n,d}$, we consider its minimum over the unit sphere $\mbS^{n-1}$. The minimum, denoted as $p^*$, equals the optimal value of the following optimization problem.
% \begin{equation}\label{eq:poly_min}
% 	\begin{array}{cl}
% 		\max\limits_{\gamma\in{\scriptsize\reals}} & \gamma\\
% 		\st & p(x)-\gamma\|x\|_2^d\ge 0,\quad\forall x\in\reals^n.
% 	\end{array}
% \end{equation}
% Traditional sos method replaces the nonnegativity constraint with the sos constraint:
% \begin{equation}
% 	\begin{array}{cl}	\max\limits_{\gamma\in{\scriptsize\reals}} & \gamma\\
% 		\st & p(x)-\gamma\|x\|_2^d\in\HSOS_{n,d}.
% 	\end{array}
% \end{equation}
% This optimization problem yields a lower bound on $p^*$ since $\HSOS_{n,d}\subseteq \HNN_{n,d}$. Though effective in many cases, however, it will fail for polynomials in $\HNN_{n,d}\backslash\HSOS_{n,d}$. We, instead, replace the nonnegativity constraint with the sos$_\mcD$ constraint:
% \begin{equation}\label{eq:poly_min_disos_relax}
% 	\begin{array}{cl}
% 		\max\limits_{\gamma\in{\scriptsize\reals}} & \gamma\\
% 		\st & p(x)-\gamma\|x\|_2^d\in\DiSOS_{n,d}(\mcD).
% 	\end{array}
% \end{equation}
% Given an algebraic disjunction $\mcD$, this optimization problem is convex and can be solved with multiple small, fixed-size SDPs. 

% Our goal is to certify nonnegativity of polynomials. As observed earlier, it suffices to focus on homogeneous polynomials. Moreover, given a polynomial $p\in H_{n,d}$ and an algebraic disjunction $\mcD$, one can check whether $p$ is sos$^d_\mcD$ by solving an SDP. Thus, 
% The key remaining challenge is 


A natural question is how to search for disjunctive proofs of nonnegativity of a polynomial~$p$; i.e., for the identities in~\eqref{def:sos_ad}. Note that as written, the joint search for the polynomials $q_{k,j}$ that define the algebraic disjunction $\mathcal{D}$ and the sos multipliers $s_{k,j}$ cannot be formulated as a convex optimization problem. An obvious workaround to this issue is to alternate the search for these decision variables. The details of this approach are described in the rest of this section together with some examples. In Section~\ref{sec:disos_psatz} and Section~\ref{sec:disos_new}, we present a more theoretically sound approach to finding disjunctive proofs of nonnegativity. We show that by appropriately fixing a family of algebraic disjunctions, we can always construct low-degree proofs of nonnegativity using convex optimization.


%A more principled approach is to fix the algebraic disjunctions, which would make the search for the sos multipliers a semidefinite program. However, the challenge is to do this in a way that guarantees existence of proofs of polynomial nonnegativity. This is our main contribution and it is presented in Section~\ref{sec:disos_psatz} and Section~\ref{sec:disos_new}.


% study whether proofs of nonnegativity can be guaranteed

% In the remainder of this section, we propose an ``alternating maximization 


% An important question is how 
% to identify a suitable algebraic disjunction that yields a low-degree disjunctive sos proof (assuming one exists). We propose two algorithms for constructing such disjunctions. The first algorithm, described below, is an alternating maximization heuristic. This method jointly searches for $q_{k,j}$ and its corresponding multiplier $s_{k,j}$ in~\eqref{def:sos_ad}. Although the joint search leads to a non-convex problem, we can alternate between the two components by solving small, fixed-sized SDPs for each subproblem.
% This iterative process generates an improving sequence of lower bounds of $p^*$. 
% The second algorithm, presented in Section~\ref{sec:disos_grid}, is a spherical-caps covering method. It employs a fixed choice of algebraic disjunctions based on Theorem~\ref{cor:strictly_pos}. This approach provides a complete SDP-based proof system for certifying nonnegativity of polynomials.


\paragraph{An alternating maximization algorithm.}

% In the joint search for $q_{k,j}$ and its corresponding multiplier $s_{k,j}$, the algebraic disjunction must remain valid at each step. Without prior knowledge of the polynomials forming the disjunction, 
% a valid algebraic disjunction of $\reals^n$ can be constructed by selecting $\ell$ polynomials $q_1,\ldots, q_{\ell}$ and enumerating each side of them.

% One way to guarantee this is to take $\ell$ polynomials $h_1, \ldots, h_{\ell}$ for some positive integer $\ell$,  let $r = 2^{\ell}$, let $\varepsilon \in \{-1,1\}^{r \times \ell}$ with all rows of $\varepsilon$ distinct, and 
% let $q_{kj} = \varepsilon_{kj}h_j$.

In the joint search for the polynomials~$q_{k,j}$ and the sos multipliers~$s_{k,j}$ in~\eqref{def:sos_ad}, a challenge is to ensure that the polynomials~$q_{k,j}$ form a valid algebraic disjunction (see Definition~\ref{def:disjunction}). 
% keeping
% the algebraic disjunction defined by $q_{k,j}$ valid at each step. 
A simple way to guarantee this is to take $\ell$ polynomials $h_1,\ldots, h_{\ell}$, let $r=2^{\ell}$, $\varepsilon \in \{-1,1\}^{r \times \ell}$ with all rows of $\varepsilon$ distinct, and let 
\[q_{k,j}=\varepsilon_{k,j}h_j,\quad k=1,\ldots,r, \quad j=1,\ldots,\ell.\]
% where $r=2^{\ell}$, and $\varepsilon_1,\ldots,\varepsilon_{r}\in\{-1,1\}^{\ell}$ denote all binary sign vectors of length $\ell$. 
With this construction, it is easy to see that the union of the corresponding $r$ subregions $\Omega_1,\ldots,\Omega_r$ as defined in Definition~\ref{def:disjunction} is $\reals^n$.
% always form a partition of~$\reals^n$,
% \[\Omega_{k}=\{x\mid \varepsilon_{k,j}h_j(x)\ge 0,\quad j=1,\ldots, \ell\},\quad k=1,\ldots,r.\]
% This observation is the key idea behind the alternating maximization algorithm. 
Building on this, given a polynomial $p$ of degree $d$ and an integer $\bar{d}\ge d$, the problem boils down to finding polynomials~$h_j$ and sos multipliers~$s_{k,j}$ such that
\begin{equation}\label{eq:am_joint_feas}
		p(x) =s_{k,0}(x)+\sum\limits_{j=1}^{\ell} \varepsilon_{k,j}h_j(x)s_{k,j}(x),\quad  k=1,\ldots,r,    
\end{equation}
% where $s_{k,0}\in\SOS_{n,d},s_{k,j}\in\SOS_{n,d-d'}$ for $k=1,\ldots,r,j = 1,\ldots, \ell$, and $\deg h_j\le d'$ for $j=1,\ldots,\ell$. Here,~$d'\in\{2,\ldots,d\}$ is fixed and even. 
where $\deg s_{k,0}\le \bar{d}$, 
$\deg(h_js_{k,j})\le \bar{d}$ for $k=1,\ldots,r$ and $j=1,\ldots,\ell$. 
% To track progress across iterations, 
To approach this problem using convex optimization, we alternate between the search for the polynomials~$h_j$ and the sos multipliers~$s_{k,j}$. 
% The alternating maximization algorithm proceeds by alternating between the polynomials $h_j$ and sos multipliers $s_{k,j}$.
To track progress across iterations, we always maximize a lower bound on~$p$ that can be certified with identities of the form~\eqref{eq:am_joint_feas}. 
More concretely, we alternate between the following two steps, which can both be reformulated as a semidefinite program:
% We use the largest lower bound on~$p$ that we can certify with identities of the form~\eqref{eq:am_joint_feas} 
% in the current round, as our goal is to certify nonnegativity of $p$. 
% Thus, we apply an alternating optimization scheme to the following problem:
% In order to iteratively refine $q_i$ until becomes feasible, we introduce a slack variable $\gamma$ to measure the quality of the current algebraic disjunction and solve
% \begin{equation}\label{eq:am_joint}
% 	\begin{array}{cl}
% \max\limits_{\gamma,h_j,s_{k,j}} & \gamma\\
% 		\st & p(x)-\gamma =s_{k,0}(x)+\sum\limits_{j=1}^{\ell} \varepsilon_{k,j}h_j(x)s_{k,j}(x),\quad  k=1,\ldots,r\\
%   & \deg s_{k,0}\le \bar{d},\ 
% \deg(h_js_{k,j})\le \bar{d},\quad  k=1,\ldots,r,\ j = 1,\ldots, \ell,\\
% & s_{k,j} \text{ is sos},\quad  k=1,\ldots,r,\ j = 0,\ldots, \ell.
% 	\end{array}
% \end{equation}
% Problem~\eqref{eq:am_joint} is non-convex, because given an optimal solution to~\eqref{eq:am_joint}, denoted by
% $\{s^*_{k,0}\}_{k=1,\ldots,r}$, $\{(q_i^*,s^*_{k,i})\}_{k=1,\ldots,r,i=1,\ldots,\ell}$, we have that $\{s^*_{k,0}\}_{k=1,\ldots,r}, \{(\alpha q_i^*,\alpha^{-1}s^*_{k,i})\}_{k=1,\ldots,r,i=1,\ldots,\ell}$ is also optimal for any constant $\alpha>0$. 
% In practice, however, 
% Note that if the polynomials $h_j$ are fixed, optimizing over the multipliers $s_{k,j}$ is a convex problem that can be solved via SDP, and conversely, if the $s_{k,j}$ are fixed, optimizing over the $h_j$ is also convex. In practice, we begin by generating initial polynomials $h_j$ for $i=1,\ldots, \ell$ (e.g., by randomly generating their coefficients), and then iteratively update $h_j$ and $s_{k,j}$ in an alternating fashion as follows:
\begin{enumerate}
    \item 
    % Let $h_j^*$ for $j=1,\ldots,\ell$ be an optimal solution returned by the solver from the previous step (or the initial polynomials if this is the first round). Let $\gamma^*$ and $s_{k,j}^*$ for $k=1,\ldots,r$ and $j=0,\ldots,\ell$ denote an optimal solution to the following problem returned by the solver:
    Given polynomials~$h_j^*$ for $j=1,\ldots,\ell$, solve the SDP
    \begin{equation*}
	\begin{array}{cl}
\max\limits_{\gamma,s_{k,j}} & \gamma\\
		\st & p(x)-\gamma =s_{k,0}(x)+\sum\limits_{j=1}^{\ell} \varepsilon_{k,j}h_j^*(x)s_{k,j}(x),\quad  k=1,\ldots,r\\
         & \deg s_{k,0}\le \bar{d},\ 
\deg(h_j^*s_{k,j})\le \bar{d},\quad  k=1,\ldots,r,\ j = 1,\ldots, \ell,\\
& s_{k,j} \text{ is sos},\quad  k=1,\ldots,r,\ j = 0,\ldots, \ell.
  % & s_{k,0}\in\SOS_{n,d},s_{k,j}\in\SOS_{n,d-d'},\quad  k=1,\ldots,r,\ j = 1,\ldots, \ell.\\
	\end{array}
\end{equation*}
     % If $\gamma^*\ge 0$, return the disjunctive sos proof of nonnegativity of $p$.
    % \item Let $s_{k,j}^*$ for $k=1,\ldots,r$ and $j=1,\ldots,\ell$ be an optimal solution returned by the solver from the previous step. Let $\hat{\gamma}^*$, $h_j^*$ for $j=1,\ldots,\ell$ and $s_{k,0}^*$ for $k=1,\ldots,r$ denote an optimal solution to the following problem returned by the solver:
    \item Given sos polynomials~$s_{k,j}^*$ for $k=1,\ldots,r$ and $j=1,\ldots,\ell$, solve the SDP\footnote{{In fact, both Step~1 and Step~2 can be decomposed into $r$ independent SDPs and solved in parallel; see Remark~\ref{remark:uniform_grid} in Section~\ref{sec:disos_grid}.}} 
     \begin{equation*}
	\begin{array}{cl}
\max\limits_{\gamma,h_j,s_{k,0}} & \gamma\\
		\st & p(x)-\gamma =s_{k,0}(x)+\sum\limits_{j=1}^{\ell} \varepsilon_{k,j}h_j(x)s_{k,j}^*(x),\quad  k=1,\ldots,r\\
         & \deg s_{k,0}\le \bar{d},\ \deg(h_js_{k,j}^*)\le \bar{d},\quad  k=1,\ldots,r,\ j = 1,\ldots, \ell,\\
& s_{k,0} \text{ is sos},\quad  k=1,\ldots,r.
  % & s_{k,0}\in\SOS_{n,d}\quad  k=1,\ldots,r,\\
  % & \deg h_j\le d'\quad  j = 1,\ldots, \ell.\\
	\end{array}
\end{equation*}   
    % If $\hat{\gamma}^*\ge 0$, return the disjunctive sos proof of nonnegativity of $p$.
\end{enumerate}
To run Step 1, we let the polynomials $h_j^*$ for ${j}=1,\ldots, \ell$ be arbitrary (e.g., chosen at random) if we are in the first iteration, or an optimal solution to Step 2 if we are in subsequent odd iterations. Similarly, to run Step 2, we let $s_{k,j}^*$ for $k=1,\ldots,r$ and $j=1,\ldots,\ell$ be an optimal solution to Step 1.



% we fix $h_j^*$ for $i=1,\ldots, \ell$, potentially chosen at random,

% We initialize the loop by generating polynomials $h_j^*$ for $i=1,\ldots, \ell$ (e.g., by randomly generating their coefficients).


Applying the alternating maximization algorithm to the Motzkin polynomial in~\eqref{eq:motzkin_nonhomo} yields two additional disjunctive sos proofs of nonnegativity. In this example, we set $\ell=1$ and $\bar{d}=6$. To initialize the algorithm, we generate $h_1$ randomly, once of degree 1 and once of degree 4. The alternating maximization algorithm converges in fewer than 20 iterations and produces (after some minor post-processing to polish the coefficients to rational numbers) the following two algebraic identities, respectively: 
% in~\eqref{eq:Motzkin.disos.proof2} and~\eqref{eq:Motzkin.disos.proof3}. 
% leads respectively to the identities in~\eqref{eq:Motzkin.disos.proof2} and~\eqref{eq:Motzkin.disos.proof3}. 
% With this setup, the algorithm converges in fewer than 20 iterations and produces the algebraic identities shown below. 
% The coefficients are then manually adjusted to rational values.
\begin{equation}\label{eq:Motzkin.disos.proof2}
 \begin{aligned}
    M(x_1,x_2) &= (1-x_1x_2^2)^2 + (x_1^2x_2-x_1x_2)^2 + 2x_1(x_1x_2-x_2)^2\\
    & = (1+x_1x_2^2)^2 + (x_1^2x_2+x_1x_2)^2 - 2x_1(x_1x_2+x_2)^2,\\
\end{aligned} 
\end{equation}
\begin{equation}\label{eq:Motzkin.disos.proof3}
 \begin{aligned}
M(x_1,x_2) &= \frac{1}{2}(x_1^2x_2+x_2^3-2x_2)^2 + (x_2^2-1)^2  +\frac{1}{2}(x_1^4-x_2^4-2x_1^2+2x_2^2)x_2^2\\
  &=\frac{1}{2}(x_1^3+x_1x_2^2-2x_1)^2 + (x_1^2-1)^2 - \frac{1}{2}(x_1^4-x_2^4-2x_1^2+2x_2^2)x_1^2.
\end{aligned} 
\end{equation}
% \begin{equation}
%  \begin{aligned}
%     M_h &= (x_3^3-x_1x_2^2)^2 + (x_1^2x_2-x_1x_2x_3)^2 + 2x_3x_1(x_1x_2-x_2x_3)^2\\
%     & = (x_3^3+x_1x_2^2)^2 + (x_1^2x_2+x_1x_2x_3)^2 - 2x_3x_1(x_1x_2+x_2x_3)^2,\\
%     \\
% M_h &= \frac{1}{2}(x_1^2x_2+x_2^3-2x_2x_3^2)^2 + (x_2^2x_3-x_3^3)^2  +\frac{1}{2}(x_1^4-x_2^4-2x_1^2x_3^2+2x_2^2x_3^2)x_2^2\\
%   &=\frac{1}{2}(x_1^3+x_1x_2^2-2x_1x_3^2)^2 + (x_1^2x_3-x_3^3)^2 - \frac{1}{2}(x_1^4-x_2^4-2x_1^2x_3^2+2x_2^2x_3^2)x_1^2.
% \end{aligned} 
% \end{equation}
Figure~\ref{fig:Motzkin} demonstrates the subregions associated with the three degree-6 disjunctive sos proofs of nonnegativity of the Motzkin polynomial in~\eqref{eq:Motzkin.disos.proof},~\eqref{eq:Motzkin.disos.proof2}, and~\eqref{eq:Motzkin.disos.proof3}.


%We produce additional disjunctive sos certificates with rational coefficients for classical nonnegative non-sos polynomials, drawn from~\cite{Reznick2000}.

Similarly, we are able to find disjunctive proofs of nonnegativity of several polynomials that are known to \emph{not} be a sum of squares. We present three such certificates below.
\begin{itemize}[noitemsep]
\item Choi-Lam-1~\cite{ChoiLam1976}. Let $ CL_1(x_{1},x_{2},x_{3},x_{4}) = -4x_{1} x_{2} x_{3} x_{4} + x_{1}^{2} x_{2}^{2} + x_{1}^{2} x_{3}^{2} + x_{2}^{2} x_{3}^{2} + x_{4}^{4}$.
Then $CL_1 = s_1 + h\,s_2 = s_3 - h\,s_4$, with $h = x_{1} x_{2}$,
\[
  \begin{array}{rcl@{\qquad}rcl}
    s_1 & = & (x_{1} x_{3} - x_{2} x_{3})^2 +\, (x_{1} x_{2} - x_{4}^{2})^2, & s_2 & = & 2\,(-x_{3} + x_{4})^2, \\
    s_3 & = & (x_{1} x_{3} + x_{2} x_{3})^2 +\, (x_{1} x_{2} + x_{4}^{2})^2, & s_4 & = & 2\,(x_{3} + x_{4})^2.
  \end{array}
\]
\item Choi-Lam-2~\cite{ChoiLam1977}.
Let $CL_2(x_{1},x_{2},x_{3}) = -3x_{1}^{2} x_{2}^{2} x_{3}^{2} + x_{1}^{2} x_{3}^{4} + x_{1}^{4} x_{2}^{2} + x_{2}^{4} x_{3}^{2}$. 
Then $CL_2~=~ s_1 + h\,s_2 = s_3 - h\,s_4$, with $h = x_{1} x_{2}$,
\[
  \begin{array}{rcl@{\qquad}rcl}
    s_1 & = & (-x_{1} x_{3}^{2} + x_{1}^{2} x_{2})^2 +\, (-x_{1} x_{2} x_{3} + x_{2}^{2} x_{3})^2, & s_2 & = & 2\,(-x_{1} x_{3} + x_{2} x_{3})^2, \\
    s_3 & = & (x_{1} x_{3}^{2} + x_{1}^{2} x_{2})^2 +\, (x_{1} x_{2} x_{3} + x_{2}^{2} x_{3})^2, & s_4 & = & 2\,(x_{1} x_{3} + x_{2} x_{3})^2.
  \end{array}
\]
\item Stengle-1~\cite{Stengle1979}. Let $ S_1(x_{1},x_{2},x_{3}) = -2x_{1} x_{2}^{2} x_{3}^{3} + x_{1}^{2} x_{3}^{4} - 2x_{1}^{3} x_{2}^{2} x_{3} + x_{1}^{3} x_{3}^{3} + 2x_{1}^{4} x_{3}^{2} + x_{1}^{6} + x_{2}^{4} x_{3}^{2}$.
Then $S_1 = s_1 + h\,s_2 = s_3 - h\,s_4$, with $h = x_{1} x_{3}$,
\[
  \begin{array}{rcl@{\qquad}rcl}
    s_1 & = & (-x_{1} x_{3}^{2} - x_{1}^{3} + x_{2}^{2} x_{3})^2, & s_2 & = & (x_{1} x_{3})^2, \\
    s_3 & = & \multicolumn{4}{l}{(x_{1} x_{3}^{2} + \tfrac{1}{2}\,x_{1}^{2} x_{3} + \tfrac{7}{10}\,x_{1}^{3} - \tfrac{6}{7}\,x_{2}^{2} x_{3})^2 +\, \tfrac{1}{14}\,(x_{1} x_{2} x_{3})^2} \\
     &  & \multicolumn{4}{l}{+\, \tfrac{7}{20}\,(x_{1}^{2} x_{3} - x_{1}^{3} + \tfrac{5}{7}\,x_{2}^{2} x_{3})^2 +\, \tfrac{4}{25}\,(x_{1}^{3} - \tfrac{5}{7}\,x_{2}^{2} x_{3})^2 +\, \tfrac{1}{196}\,(x_{2}^{2} x_{3})^2,} \\
    s_4 & = & \tfrac{2}{7}\,(-\tfrac{1}{2}\,x_{1} x_{2} + x_{2} x_{3})^2. & & &
  \end{array}
\]
\end{itemize}



\begin{figure}[ht]
\centering
\includegraphics[width=0.33\textwidth]{figs/Motzkin1.pdf}\hfill
\includegraphics[width=0.33\textwidth]{figs/Motzkin2.pdf}
\includegraphics[width=0.33\textwidth]{figs/Motzkin3.pdf}
    \caption{Three different partitions of $\reals^2$ corresponding to three different degree-6 disjunctive sos proofs of nonnegativity of the Motzkin polynomial (see~\eqref{eq:motzkin_nonhomo}). Left: subregions correspond to the identity in~\eqref{eq:Motzkin.disos.proof} and are defined by $x_1x_2\ge 0$ and $-x_1x_2 \ge 0$. 
    Middle: subregions correspond to the identity in~\eqref{eq:Motzkin.disos.proof2} and are defined by $x_1\ge 0$ and $-x_1\ge 0$. 
    Right: subregions correspond to the identity in~\eqref{eq:Motzkin.disos.proof3} and are defined by $x_1^4-x_2^4-2x_1^2+2x_2^2 \ge 0$ and $-(x_1^4-x_2^4-2x_1^2+2x_2^2) \ge 0$.
    }
\label{fig:Motzkin}
\end{figure}

\section{A Disjunctive Positivstellensatz}
\label{sec:disos_psatz}
%We now present our main theorem on the existence of low-degree disjunctive sum of squares certificate of nonnegativity for forms (Theorem~\ref{cor:strictly_pos} below). 

In this section, we introduce our main result in Theorem~\ref{cor:strictly_pos} which shows that one can a priori fix a family of low-degree algebraic disjunctions with respect to which there always exists a low-degree disjunctive sos proof of nonnegativity. 


Let us establish some basic notation related to \emph{forms}, i.e., homogeneous polynomials, which will be useful in the remainder of this paper.
We denote the set of forms, nonnegative forms, and sum of squares forms in~$n$ variables and degree~$d$ by $H_{n,d},\HNN_{n,d},$ and $\HSOS_{n,d}$ respectively. 
Observe that~$\HNN_{n,d}=\NN_{n,d}\cap H_{n,d}$ and $\HSOS_{n,d}=\SOS_{n,d}\cap H_{n,d}$. 
Additionally, we say that a form~$p$ in~$n$ variables is \emph{positive definite} if $p(x)>0$ for all $x\in\reals^n$ and $x\neq0$.
% (Theorem~\ref{cor:strictly_pos} proven in Section~\ref{sec:local_psatz}). 
%Recall $H_{n,d}$ denote the set of homogeneous polynomials (i.e. forms) in $n$ variables 
We are now ready to present our main result.
\begin{theorem}[Disjunctive Positivstellensatz]\label{cor:strictly_pos}
For any fixed dimension~$n$ and even degree~$d$, there exists an explicit family of algebraic disjunctions
$\mcD^m$, indexed by $m$, such that for any positive
 definite form $p\in H_{n,d}$, we have $p\in \DiSOS_{n,d}^{d}(\mcD^m)$ for some~$m$. Furthermore, each subset in $\mcD^m$ contains exactly one form of degree at most~$d$.\footnote{Our proof in fact provides several constructions: the forms in the algebraic disjunction may be chosen to have any even degree between~2 and~$d$.}
\end{theorem}
%We refer to this result as a disjunctive Positivstellensatz.
% More specifically, we show that one can a priori fix a family of low-degree algebraic disjunctions with respect to which existence of a disjunctive sos proof of nonnegativity is always possible. 

We prove Theorem~\ref{cor:strictly_pos} as follows.
In Section~\ref{sec:local_psatz}, we introduce local Positivstellens\"atze, which provide a single algebraic identity proving nonnegativity of a polynomial in the neighborhood of a given point where it is known to be positive.
In Section~\ref{sec:global_psatz}, we combine these neighborhoods to cover the entire Euclidean space and use the corresponding algebraic identities to construct a disjunctive sos proof of (global) nonnegativity.
In Section~\ref{sec:disos_grid}, we show how this theorem leads to a hierarchy of SDPs---with semidefinite constraints of fixed size---that produces arbitrarily tight lower bounds on the minimum value of a homogeneous polynomial over the unit sphere. 


% \subsection{Reduction to the Homogeneous Case}
% \label{sec:reduction}
% As odd-degree polynomials cannot be nonnegative on the entire Euclidean space, we focus only on even-degree polynomials. Moreover, it suffices to consider even-degree forms. For a polynomial~$p$ in~$n$ variables and degree~$d$, we denote its homogenization by
% $$p_h(x,x_{n+1})\defeq x_{n+1}^dp(x/x_{n+1}).$$
% % It is well known that homogenization preserves nonnegativity and sos property, i.e., 
% It is straightforward to check that homogenization preserves both nonnegativity and the sos property; that is, $p\in\NN_{n,d}$ (resp., $p\in\SOS_{n,d}$) if and only if $p_h\in\HNN_{n+1,d}$ (resp., $p_h\in\HSOS_{n+1,d}$). The following lemma further shows that the disjunctive sos proof of nonnegativity of~$p$ can be directly converted into one for~$p_h$, and vice versa. This justifies our focus on homogeneous polynomials.

% \begin{lemma}
% Let $\mcD=\{\{q_{k,j}\}_{j=1,\ldots, n_k},\quad k=1,\ldots,r\}$ be an algebraic disjunction of $\reals^n$ and $\widehat{\mcD}=\{\{\hat{q}_{k,j}\}_{j=1,\ldots, n_k},\quad k=1,\ldots,r\}$, where~$\hat{q}_{k,j}$ denotes the  homogenization of~$q_{k,j}$.
% Then a polynomial~$p$ in~$n$ variables and degree~$d$ is sos$_{\mcD}^d$ if and only if~$p_h$ is sos$_{\widehat{\mcD}}^d$. 
% \end{lemma}

% \begin{proof}
% If $p$ is sos$^d_\mcD$, then it can be written in the form of the identities in~\eqref{def:sos_ad}. Homogenizing all polynomials in~\eqref{def:sos_ad} immediately yields that $p_h$ is sos$^d_{\widehat{\mcD}}$.
% Conversely, suppose $p_h$ is sos$^d_{\widehat{\mcD}}$. Then it can be expressed as 
% \[    p_h(x)=\hat{s}_{k,0}(x)+\sum_{j=1}^{n_k} \hat{s}_{k,j}(x)\hat{q}_{k,j}(x),\quad  k=1,\ldots, r,\]
% where $\hat{s}_{k,0}$ is sos with $\deg \hat{s}_{k,0}\le d$, and $\hat{s}_{k,j}$ is sos with $\deg(\hat{s}_{k,j}\hat{q}_{k,j})\le d$ for $k=1,\ldots,r$ and $j=1,\ldots,n_k$. Setting $x_{n+1}=1$ in the above identities shows that $p$ is sos$^d_{\mcD}$.

% \end{proof}


%an SDP-based converging hierarchy of lower bounds for the problem of minimizing a homogeneous polynomial over the unit sphere, where the size of the largest semidefinite constraint is fixed throughout the hierarchy.


%to a hierarchy of arbitrarily tight lower bounds on the the minimum value of a polynomial on the 

% Specifically, we show that the
% there exists an explicit family of low-degree algebraic disjunctions, such that any pd form is sos with respect to a member in the family. 
% The algorithmic implications of the theorem are presented in Section~\ref{sec:disos_grid}.
% The key idea in proving the theorem is to show the existence of a low-degree algebraic proof locally. Specifically, if $p$ is positive at a point $x_0$, there is an algebraic identity showing that a small distance between any point and $x_0$ implies the nonnegativity of $p$. This yields a subregion centered at $x_0$ over which $p$ has an algebraic proof of nonnegativity. Thus, for a pd form, we can cover the unit sphere with a finite number of subregions to obtain a disjunctive certificate.

% Moreover, we show that there exists an explicit family of algebraic disjunctions, such that any pd form is sos with respect to a certain algebraic disjunction in the family. This leads to the uniform gridding algorithm in Section~\ref{sec:disos_grid}.

\subsection{Local Positivstellens\"atze}\label{sec:local_psatz}
For a polynomial~$p$ in~$n$ variables and degree~$d$, we denote its homogenization by
$$p_h(x,x_{n+1})\defeq x_{n+1}^dp(x/x_{n+1}).$$
It is straightforward to check that homogenization preserves both nonnegativity and the sos property; that is, $p\in\NN_{n,d}$ (resp., $p\in\SOS_{n,d}$) if and only if $p_h\in\HNN_{n+1,d}$ (resp., $p_h\in\HSOS_{n+1,d}$). Hence, given a general (not necessarily homogeneous) polynomial~$p$, we can construct a disjunctive sos proof of nonnegativity of~$p$ by first obtaining one for $p_h$ and then setting the last variable to 1 in the identities. This justifies our focus on homogeneous polynomials. Moreover, since odd-degree polynomials cannot be nonnegative on the entire Euclidean space, it suffices to consider even-degree forms. 

We now turn to establishing local certificates of nonnegativity. Suppose a polynomial~$p$ takes {a} positive value at some point $x^*\in\mathbb{R}^n$. By continuity, $p$ remains nonnegative within a sufficiently small neighborhood of $x^*$. This observation motivates the development of a local Positivstellensatz: an algebraic identity demonstrating that if the angle between a point $x$ and $x^*$ is sufficiently small, then $p(x)\ge 0$. 

% To prepare for the theorem, we introduce some notation. 
We introduce some notation for the theorems that follow. We denote the unit sphere in~$\reals^n$ by~$\mbS^{n-1}$, and the $n\times n$ identity matrix by $I_n$. 
For a polynomial~$p$, we denote its largest coefficient in absolute value (resp.\ the sum of the absolute values of its coefficients) in the standard monomial basis by $\|p\|_{\infty}$ (resp.\ by $\|p\|_1$).

% let $\|p\|_{\infty}$ denote the largest coefficient of~$p$ in absolute value in the standard monomial basis, and let $\|p\|_1$ denote the sum of the absolute values of its coefficients. 

% For a polynomial~$p$, let $\|p\|_{\infty}$ denote the largest coefficient of~$p$ in absolute value in the standard monomial basis, and let $\|p\|_1$ denote the sum of the absolute values of its coefficients. 

\begin{theorem}\label{thm:degd_homo}
For a polynomial $p\in H_{n,d}$ with $d$ even, if $p(x^*)>0$ for some $x^*\in \mbS^{n-1}$, then for $t>0$ sufficiently large, $p(x)$ can be written as
\begin{equation}\label{eq:degd_homo}
    p(x)=\left(\frac{p(x^*)}{2}(x^Tx^*)^d-th_d(x;x^*)\right)+\sigma(x),
\end{equation}
where $\sigma(x)\in\HSOS_{n,d}$ and 
\begin{equation}\label{eq:degd_homo_h}
h_d(x;x^*)\defeq\sum_{k=1}^{d/2}\left(\|x\|_2^2-(x^Tx^*)^2\right)^k(x^Tx^*)^{d-2k}.
\end{equation}
%TEMP:
%$$c(x)=\frac{p(x^*)}{2}(x^Tx^*)^d-t\left( \sum_{k=1}^{d/2}\left(\|x\|_2^2-(x^Tx^*)^2\right)^k(x^Tx^*)^{d-2k} \right)$$
%
In particular, the above statement holds for any
\begin{equation}\label{eq:degd_homo_t}
t\ge T(p;x^*)\defeq\sup_{U\in\reals^{n\times n}:U^TU=I_n}\left( \frac{n}{2p(x^*)}\|p(Ux)\|_{\infty}^2 + \|p(Ux)\|_1\right).  
\end{equation}
% Here, $\|p\|_{\infty}$ denotes the largest coefficient of~$p$ in absolute value, and $\|p\|_1$ denotes the sum of the absolute values of the coefficients of~$p$. 
\end{theorem}

\begin{remark}
The identity in~\eqref{eq:degd_homo} shows the implication that
\begin{equation}\label{eq:imp_homo}
\frac{p(x^*)}{2}(x^Tx^*)^d-th_d(x;x^*)\ge 0\quad\Rightarrow\quad p(x)
\ge 0.
\end{equation}
To gain geometric intuition, consider the region defined by
\begin{equation}\label{eq:spherical_cone}
A(x^Tx^*)^d-th_d(x;x^*)\ge 0,
\end{equation}
where $A$ and $t$ are positive scalars. Let $\theta$ denote the angle between $x$ and $x^*$. Since   $x^Tx^*=\|x\|_2\cos\theta$, the inequality above is equivalent to 
$\sum_{k=1}^{d/2}(\tan\theta)^{2k}\le A/t$.
Therefore, the subregion corresponds to $|\theta|\le\arctan z$, where $z$ is the unique positive real root of 
$\sum_{k=1}^{d/2} z^{2k}= A/t$. Consequently, identity~\eqref{eq:degd_homo} shows that a sufficiently small angle between $x$ and $x^*$ implies nonnegativity of $p(x)$.
\textcolor{black}{Figure~\ref{fig:spherical_cap} illustrates the intersection of the region defined by~\eqref{eq:spherical_cone} and $\mbS^{n-1}$, which is a spherical cap.}
\end{remark}

\begin{figure}[ht]
    \centering
    \includegraphics[width=0.3\textwidth]{figs/spherical_caps.pdf}
    \caption{\textcolor{black}{A spherical cap centered at $x^*$: the local identity~\eqref{eq:degd_homo} certifies $p(x)\ge 0$ for every $x\in \mbS^{n-1}$ within angle $\theta$ of $x^*$.}}
    \label{fig:spherical_cap}
\end{figure}

% To prove Theorem~\ref{thm:degd_homo}, we initially aim to establish the implication
% \begin{equation}\label{eq:imp_nonhomo}
%  \|x-x^*\|_2\le \delta\ \Rightarrow\ p(x)\ge 0.
% \end{equation}
% However, the non-homogeneous nature of $\delta- \|x-x^*\|_2$ makes it difficult to construct an algebraic identity involving only homogeneous polynomials. To address this, we first allow non-homogeneous polynomials in the algebraic identities. Under this less restrictive setting, Lemma~\ref{thm:degd_nonhomo} shows that such an algebraic identity can indeed be constructed for a general (not necessarily homogeneous) polynomial. We then extend this result to the homogeneous case. This extension reveals a key insight: in the homogeneous setting, the vector inner product naturally controls the distance between points. In particular, the inner product encodes the angle between vectors on $\mbS^{n-1}$ while preserving homogeneity in the algebraic identities. Finally, this observation leads to the derivation of $h_d(x;x^*)$ in Theorem~\ref{thm:degd_homo}, whose form may not appear immediately intuitive at first glance.

To prove Theorem~\ref{thm:degd_homo}, we first establish the following lemma. 
% to prove Theorem~\ref{thm:degd_homo}. 
% It shows that one can construct a non-homogeneous algebraic identity for a general (not necessarily homogeneous) polynomial. This result will then be extended to the homogeneous case. 
Let us recall the multi-index notation that will be used in its proof. For an $n$-tuple of nonnegative integers $\alpha\defeq(\alpha_1,\ldots,\alpha_n)$, and a vector $x\in\reals^n$, we write $x^{\alpha}$ to denote $x_1^{\alpha_1}\ldots x_n^{\alpha_n}$ and $|\alpha|$ to denote $\sum_{i=1}^n \alpha_i$. For two multi-indices $\alpha$ and $\beta$, we denote by $\alpha+\beta$ their componentwise sum.  


% we recall the standard multi-index notation. An $n$-dimensional multi-index $\alpha$ is an $n$-tuple of nonnegative integers, written as $\alpha=(\alpha_1,\ldots,\alpha_n)$. We define  $|\alpha|\defeq\sum_{i=1}^n \alpha_i$. For any $x\in\reals^n$, we write  $x^{\alpha}=x_1^{\alpha_1}\ldots x_n^{\alpha_n}$. For two multi-indices $\alpha$ and $\beta$, we denote by $\alpha+\beta$ as their componentwise sum.

\begin{lemma}\label{thm:degd_nonhomo}
Consider a (not necessarily homogeneous) polynomial~$p$ in $n$ variables and degree at most $d$, where $d$ is even. If $p(x^*)>0$ for some $x^*\in\reals^n$, then for $t>0$ sufficiently large, $p(x)$ can be written as
\[p(x)=\left(\frac{p(x^*)}{2}-tg_d(x-x^*)\right)+\sigma(x),\]
where $\sigma(x)\in\SOS_{n,d}$ and 
\begin{equation}\label{eq:degd_nonhomo_g}
  g_d(x)\defeq\sum_{k=1}^{d/2}\|x\|_2^{2k}.  
\end{equation}
In particular, the above statement holds for any 
% $t$ such that
\[t\ge \tilde{T}(p;x^*)\defeq\dfrac{n}{2p(x^*)}\|q\|^2_{\infty} + \|q\|_1,\]
where $q(x)=p(x+x^*)-p(x^*)$. 
% For any multi-index $\alpha$, $q_{\alpha}$ denotes the coefficient of the term $x^{\alpha}$ in $q(x)$, and $\|q\|_1$ denotes the sum of the absolute value of the coefficients of $q$.
\end{lemma}

\begin{proof}
Without loss of generality, we assume $x^*=0$; otherwise, applying the statement of the lemma to the polynomial $p(x+x^*)$ and recalling that the existence of an sos decomposition is invariant under an affine change of variables, yields the desired result. 
% Therefore, we assume $q(x)=p(x)-p(0)$. 
Let $$r_t(x)\defeq p(x)+tg_d(x)-\frac{p(0)}2=q(x)+tg_d(x)+\frac{p(0)}2.$$ 
Thus, our task is to show that $r_t(x)$ is sos for $t\ge \tilde{T}(p;0)$. 
To do this, we work with the following representation of the summands: 
% we work with the following representation of the polynomials that sum up to $r_t(x)$
% a gram matrix, symmetric matrices $Q$ and $D$ such that $Q+tD$ is positive semidefinite for $t\ge \tilde{T}(p;0)$, and the following two identities hold:
\[\begin{aligned}
    q(x)+\frac{p(0)}{2}&=z(x)^TQz(x)=\sum_{\alpha:|\alpha|\le d/2}\sum_{\beta:|\beta|\le d/2} Q_{\alpha,\beta}x^{\alpha+\beta},\\
    g_d(x)&=z(x)^TDz(x)=\sum_{\alpha:|\alpha|\le d/2}\sum_{\beta:|\beta|\le d/2} D_{\alpha,\beta}x^{\alpha+\beta}.
\end{aligned}\]
Here, $z(x)\defeq(1,x_1,\ldots,x_n,\ldots,x_1^{d/2},\ldots,x_n^{d/2})$ is the vector of monomials of degree at most $d/2$, and $Q$ and $D$ are symmetric matrices. We use the multi-index notation 
$Q_{\alpha,\beta}$ (resp. $D_{\alpha,\beta}$) to refer to the appropriate entry of the matrix $Q$ (resp. $D$). 

% Let $y=x-x^*$ and $A=p(x^*)$. 
% Then $q(y)=p(y+x^*)-p(x^*)=p(x)-A$, and 
% $$p(x)+tg_d(x-x^*)-\frac{A}{2}=q(y)+tg_d(y)+\frac{A}{2}.$$ Hence $\sigma$ can be expressed as a polynomial in $y$, which we denote by $\sigma(y)$ after the change of variables. Since $\sigma(x)$ is sos in $x$ if and only if $\sigma(y)$ is sos in $y$, it suffices to show that $\sigma(y)$ is sos in $y$ when $t$ is sufficiently large.

% Using multi-index notation, let $q(x)=\sum_{\alpha:|\alpha|\le d} q_{\alpha}x^{\alpha}$. 
% We can write $q(x)$ as 
% \[q(x)=z(x)^TQz(x)=\sum_{\alpha:|\alpha|\le d/2}\sum_{\beta:|\beta|\le d/2} Q_{\alpha,\beta}x^{\alpha+\beta},\]
% where $z(x)=(1,x_1,\ldots,x_n,\ldots,x_1^{d/2},\ldots,x_n^{d/2})$ is the vector of monomials of degree at most $d/2$, and $Q$ is a symmetric matrix. 
% Expanding the quadratic form gives
% where we use multi-indices as subscripts for the entries of $Q$.
% Here, we use the multi-index notation 
% $Q_{\alpha,\beta}$ to refer to the appropriate entry of the matrix $Q$. 
% Similarly, $g_d(x)$ can be expressed as $g_d(x)=z(x)^TDz(x)$ for some symmetric matrix $D$. Hence 
% $r_t(x)=z(x)^T(Q+tD)z(x)+p(0)/2$. Although the matrices $Q$ and $D$ are not uniquely determined, it suffices to show that there exist $Q$ that satisfy this identity, $Q+tD$ is positive semidefinite for $t\ge \tilde{T}(p;0)$.
Since $g_d(0)=0$, we have $D_{0,0}=0$. Observe that $g_d(x)$ is a positive linear combination of all monomials of the form $x^{2\alpha}$, where the multi-index $\alpha\neq 0$ and $|\alpha|\le d/2$. Thus, we can take $D$ to be diagonal with $D_{\alpha,\alpha}\ge 1$ for every nonzero multi-index $\alpha$. For the matrix $Q$, note that $q(0)=0$, so  $Q_{0,0}=p(0)/2$. Furthermore, the entries of $Q$ can be rearranged such that $Q_{0,\alpha}=Q_{\alpha,0}=0$ for all $\alpha$ with $|\alpha|\in\{2,\ldots, d/2\}$. This is possible because each monomial $x^{\alpha}$ with $|\alpha|\ge 2$ can be written as the product $x^{\beta}x^{\gamma}$ for some nonzero multi-indices $\beta,\gamma$ with $\beta+\gamma=\alpha$. With this rearrangement, we can write
\begin{equation}\label{eq:degd_nonhomo_rtx}
    \begin{aligned}
r_t(x)&=z(x)^T(Q+tD)z(x)\\
&=\frac{p(0)}{2n}\sum_{i=1}^n\left(1+\frac{nq_i}{p(0)}x_i\right)^2+z(x)^T(\hat{Q}+tD)z(x),
\end{aligned}
\end{equation}
where $q_i$ denotes the coefficient of the term $x_i$ in $q(x)$, and $\hat{Q}$ is defined entrywise as 
\[\hat{Q}_{\alpha,\beta}=
\begin{cases}
    0  & \text{if } \alpha= 0 \text{ or } \beta=0, \\
    Q_{\alpha,\beta}-\dfrac{nq_{i}^2}{2p(0)} & \text{if } (\alpha,\beta)=(e_i,e_i),\quad i=1,\ldots,n,\\ 
    Q_{\alpha,\beta} & \text{otherwise.}
\end{cases}\]
Here, $e_i$ is the $i$-th standard basis vector of $\reals^n$ (so that $x^{e_i}=x_i$ in multi-index notation). Finally, one can verify that when $t\ge \tilde{T}(p;0)$, the matrix $\hat{Q}+tD$ is diagonally dominant\footnote{Recall that a symmetric matrix $A=(a_{ij})$ is diagonally dominant if $a_{ii}\ge \sum_{j\neq i}|a_{ij}|$ for all $i$. It follows from Gershgorin's circle theorem~\cite{gershgorin1931uber} that diagonally {dominant} matrices are positive semidefinite.} and thus positive semidefinite. Therefore, when $t\ge\tilde{T}(p;0)$, the polynomial $z(x)^T(\hat{Q}+tD)z(x)$ is sos. In view of~\eqref{eq:degd_nonhomo_rtx}, when $t\ge\tilde{T}(p;0)$, $r_t(x)$ is also sos.
% In particular, when $x^*=0$, the coefficients of $q(x)$ are identical to the coefficients of $p(x)$, except for the constant term. Since $q_0=0$, we have $\|p\|_1\ge \|q\|_1$. Therefore, the statement of the lemma holds.

\end{proof}
We are now ready to prove Theorem~\ref{thm:degd_homo}.

% Now we extend Lemma~\ref{thm:degd_nonhomo} to the homogeneous case. The key idea is to first analyze the special case $x^*=e_1$, the first standard basis vector of $\reals^n$, and then use a rotation argument to generalize the result to an arbitrary point on the unit sphere.



% \begin{lemma}\label{lem:degd_homo_e1}
% For a polynomial $p\in H_{n,d}$ with $d$ even, if $p(e_1)>0$, then for $t>0$ sufficiently large, $p(x)$ can be written as
% \begin{equation*}
%     p(x)=\left(\frac{p(e_1)}{2}x_1^d-th_d(x;e_1)\right)+\sigma(x),
% \end{equation*}
% where $\sigma(x)\in\HSOS_{n,d}$ and 
% \[h_d(x;e_1)=\sum_{k=1}^{d/2}\left(\|x\|_2^2-x_1^2\right)^kx_1^{d-2k}.\]
% In particular, the above statement holds for any $t$ such that
% \[t\ge \hat{T}_d(p;e_1)\defeq \frac{n}{2p(e_1)}\|p\|_{\infty}^2 + \|p\|_1.\]
% \end{lemma}

\begin{proof}[Proof of Theorem~\ref{thm:degd_homo}]
We first consider the special case where $x^*=e_1$. Let $z\defeq(x_2,\ldots,x_n)$. We dehomogenize~$p$ by defining $\tilde{p}(z)\defeq p(1,z)$. Then $\tilde{p}(z)$ is a polynomial with $\tilde{p}(0)=p(e_1)>0$. 
By Lemma~\ref{thm:degd_nonhomo}, whenever $t\ge \tilde{T}(\tilde{p};0)$, $\tilde{p}(z)$ admits the representation
\[\tilde{p}(z)=\left(\frac{p(e_1)}{2}-tg_d(z)\right)+\tilde{\sigma}(z),\]
where $\tilde{\sigma}(z)\in\SOS_{n-1,d}$ and $g_d$ is defined as in~\eqref{eq:degd_nonhomo_g}. 
% \[g_d(z)=\sum\limits_{k=1}^{d/2}\|z\|_2^{2k}.\]
Next, homogenizing $\tilde{p}(z)$ yields the identity:
\[\begin{aligned}
p(x)&=p(x_1,z)=x_1^d \tilde{p}(z/x_1)\\
&=\left(\frac{p(e_1)}{2}x_1^d-t\sum\limits_{k=1}^{d/2}\|z\|_2^{2k}x_1^{d-2k}\right)+x_1^d \tilde{\sigma}(z/x_1)\\
&=\left(\frac{p(e_1)}{2}x_1^d-t\sum\limits_{k=1}^{d/2}(\|x\|_2^2-x_1^2)^{k}x_1^{d-2k}\right)+\sigma(x)\\
&=\left(\frac{p(e_1)}{2}x_1^d-th_d(x;e_1)\right)+\sigma(x),
\end{aligned}\]
where $\sigma(x)=x_1^d \tilde{\sigma}(z/x_1)\in\HSOS_{n,d}$, since homogenization preserves the sos property. Note that
\[\begin{aligned}
T(p;e_1)&\ge \dfrac{n}{2p(e_1)}\|p\|_{\infty}^2+\|p\|_1\\
&= \dfrac{n}{2\tilde{p}(0)}\|\tilde{p}\|^2_{\infty} + \|\tilde{p}\|_1\\
&=\tilde{T}(\tilde{p};0),
\end{aligned}\]
% Hence, whenever $t\ge T(p;e_1)$, we have 
% \[\begin{aligned}
% t&\ge \dfrac{n}{2p(e_1)}\|p\|_{\infty}^2+\|p\|_1\\
% &= \dfrac{n}{2\tilde{p}(0)}\|\tilde{p}\|^2_{\infty} + \|\tilde{p}\|_1\\
% &=\tilde{T}(\tilde{p};0),
% \end{aligned}\]
where the second line follows from the fact that the coefficients of~$p$ remain the same after dehomogenization. 
Therefore, in view of the fact that $h_d(x;e_1)$ is sos, the claim of Theorem~\ref{thm:degd_homo} follows in the case where $x^*=e_1$.

% we derive $\hat{T}_d(p;e_1)$ by upper bounding $\tilde{T}_d(\hat{p};e_1)$. 
% Note that the coefficients of $p(x)$ and $\hat{p}(z)$ are identical, we have
% \[
% T(\hat{p})=\dfrac{n}{2A}\max\limits_{1\le i\le n}\hat{p}_i^2 + \sum\limits_{\alpha\neq 0}|\hat{p}_{\alpha}|
% \le \dfrac{n}{2A}\|p\|_{\infty}^2+\|p\|_1= \tilde{T}(p;e_1),
% \]
% \[\begin{aligned}
% \tilde{T}_d(\hat{p};e_1)&=\dfrac{n}{2A}\max\limits_{1\le i\le n}\hat{p}_i^2 + \|\hat{p}\|_1\\
% &\le \dfrac{n}{2A}\|p\|_{\infty}^2+\|p\|_1= \hat{T}_d(p;e_1),
% \end{aligned}\]
% which completes the proof.
% \end{proof}

% \begin{proof}
To establish the claim for a general point $x^*\in \mbS^{n-1}$, observe that for any such point there exists a matrix $U\in\reals^{n\times n}$ with $U^TU=I$, such that~$Ue_1=x^*$. Define $p_U(x)\defeq p(Ux)$. Then $p_U$ is a polynomial with $p_U(e_1)=p(x^*)>0$. Applying the statement of the theorem to~$p_U$, whenever $t\ge T(p_U;e_1)$, we have
\[
p_U(x)=\left(\frac{p_U(e_1)}{2}(x^Te_1)^d-t\sum\limits_{k=1}^{d/2}(\|x\|_2^2-(x^Te_1)^2)^{k}(x^Te_1)^{d-2k}\right)+\hat{\sigma}(x),\]
where $\hat{\sigma}(x)\in\HSOS_{n,d}$.
Substituting $x$ with $U^Tx$ gives
\[\begin{aligned}
p(x)&=p_U(U^Tx)\\
&=\left(\frac{p(x^*)}{2}(x^TUe_1)^d-t\sum\limits_{k=1}^{d/2}(\|U^Tx\|_2^2-(x^TUe_1)^2)^{k}(x^TU  e_1)^{d-2k}\right)+\hat{\sigma}(U^Tx)\\
&=\left(\frac{p(x^*)}{2}(x^Tx^*)^d-t\sum\limits_{k=1}^{d/2}(\|x\|_2^2-(x^Tx^*)^2)^{k}(x^Tx^*)^{d-2k}\right)+\hat{\sigma}(U^Tx)\\
&=\left(\frac{p(x^*)}{2}(x^Tx^*)^d-th_d(x;x^*)\right)+\hat{\sigma}(U^Tx),
\end{aligned}
\]
where in the second line we use the fact that the 2-norm is unitarily invariant. 
Since ${\hat{\sigma}(x)\in\HSOS_{n,d}}$ and the sos property is preserved under a linear change of variables, we have ${\hat{\sigma}(U^Tx)\in\HSOS_{n,d}}$. Finally, note that $T(p_U;e_1)=T(p;x^*)$, since the set of orthogonal matrices is invariant under multiplication by any of its elements.
\end{proof}

% \begin{remark}
% A closer examination of the proof of Theorem~\ref{thm:degd_homo} reveals that the polynomial $\sigma(x)$ in identity~\eqref{eq:degd_homo} is not only sos; it decomposes as the sum of an explicit sos polynomial and a ``diagonally-dominant-sum-of-squares'' (dsos) polynomial with respect to a specific basis~\cite{ahmadi2017sum}. 
% Indeed, equation~\eqref{eq:degd_nonhomo_rtx} in the proof of Lemma~\ref{thm:degd_nonhomo} exhibits $\hat{\sigma}(x)$ as the sum of an explicit sos part and a dsos part. Under the change of variables $x\mapsto U^Tx$, the explicit sos component remains explicit sos, while the dsos component becomes dsos with respect to the rotated monomial basis induced by $U^T$.
% % This implies that the search for a proof of nonnegativity of a polynomial in the form of identity~\eqref{eq:degd_homo} can be achieved by a second order cone program.
% \end{remark}

\begin{remark}
A closer examination of the proof of Theorem~\ref{thm:degd_homo} reveals that the polynomial $\sigma(x)$ in identity~\eqref{eq:degd_homo} is not only sos, but in fact ``scaled-diagonally-dominant-sum-of-squares'' (sdsos) with respect to a specific basis~\cite{ahmadi2017sum}. 
Indeed, from equation~\eqref{eq:degd_nonhomo_rtx} in the proof of Lemma~\ref{thm:degd_nonhomo}, we see that the polynomial $\hat{\sigma}(x)$ is sdsos. Consequently, $\sigma(x)=\hat{\sigma}(U^Tx)$ is sdsos with respect to a rotated monomial basis induced by the linear transformation $U^T$. This implies that the search for an algebraic proof of polynomial nonnegativity in the form of identity~\eqref{eq:degd_homo} can be achieved by a second-order cone program (SOCP). This observation connects to Algorithm~\ref{alg:uniform_grid} in Section~\ref{sec:disos_grid}, which can potentially be implemented as an SOCP-based method.
\end{remark}

% Theorem~\ref{thm:degd_homo} shows the existence of an algebraic proof of nonnegativity of a polynomial~$p$ in a neighborhood of point where $p$ is positive. 
% over a subregion defined by a single degree-$d$ polynomial. 

We now present another local Positivstellensatz, analogous to Theorem~\ref{thm:degd_homo}, which provides an algebraic certificate of nonnegativity for a polynomial~$p$ in a neighborhood of a point $x^*$ where $p$ is positive. The neighborhood is defined by a single polynomial of even degree $d' \in \{2,\ldots,d-2\}$. This formulation offers additional flexibility in the choice of the sos polynomial multiplied by the polynomial defining the neighborhood.

% We now present another local Positivstellensatz, analogous to Theorem~\ref{thm:degd_homo}, that gives an algebraic proof of nonnegativity of a polynomial~$p$ in a region around a point $x^*$ where $p$ is positive, with the region defined by a single polynomial of even degree $d'\in\{2,\ldots,d-2\}$. This provides additional flexibility in the search for the sos polynomial multiplied with the defining polynomial of that region.

% Theorem~\ref{thm:degd_homo_d'} further shows that an analogous proof exists under the same assumption when the subregion is defined by a single polynomial of lower degree. 

\begin{theorem}\label{thm:degd_homo_d'}
Let $d$ be an even integer, $d'$ be an even integer in $\{2,\ldots,d-2\}$, and $p\in H_{n,d}$. 
% For a polynomial $p\in H_{n,d}$ with $d$ even and any even degree $d'\in\{2,\ldots,d-2\}$,
If $p(x^*)>0$ for some $x^*\in \mbS^{n-1}$, then for $t>0$ sufficiently large, $p(x)$ can be written as
\begin{equation*}
    p(x)=\left(\frac{p(x^*)}{2}(x^Tx^*)^{d'}-th_{d'}(x;x^*)\right)\left((x^Tx^*)^{d-d'}+h_{d-d'}(x;x^*)\right)+\sigma(x),
\end{equation*}
where $\sigma(x)\in\HSOS_{n,d}$, and $h_{d'}(x;x^*)$, $h_{d-d'}(x;x^*)$ are as in~\eqref{eq:degd_homo_h}. In particular, the above statement holds for any
\[t\ge T'(p;x^*)\defeq T(p;x^*) + \frac{p(x^*)}{2},\]
where $T(p;x^*)$ is defined as in~\eqref{eq:degd_homo_t}. 
\end{theorem}

% To prove Theorem~\ref{thm:degd_homo_d'}, it suffices to show to following lemma. 

% \begin{lemma}\label{lem:degd_homo_d'}
% Consider a polynomial $p\in H_{n,d}$ with $d$ even. Suppose there exists $A>0$ and $t>0$, such that $p(x)$ can be written as
% \[p(x)=\left(A(x^Tx^*)^d-th_d(x;x^*)\right)+\sigma(x),\]
% where $\sigma(x)\in\HSOS_{n,d}$ and $h_d(x;x^*)$ is defined in Theorem~\ref{thm:degd_homo}.
% Then, for any even degree $d'\in\{2,\ldots,d-2\}$, $p(x)$ can be written as
% \begin{equation}\label{eq:degd_homo_d'}
%     p(x)=\left(A(x^Tx^*)^{d'}-\hat{t}h_{d'}(x;x^*)\right)\left((x^Tx^*)^{d'}+h_{d-d'}(x;x^*)\right)+\hat{\sigma}(x),
% \end{equation}
% where $\hat{t}=t+A$, and $\hat{\sigma}(x)\in\HSOS_{n,d}$.
% \end{lemma}

% Theorem~\ref{thm:degd_homo_d'} is closely related to Theorem~\ref{thm:degd_homo}. Lemma~\ref{thm:degd_nonhomo_d'} below illustrates this connection in the non-homogeneous setting.

% The proof of Theorem~\ref{thm:degd_homo_d'} employs the statement of Theorem~\ref{thm:degd_homo}. To this end, we first establish the following lemma.

% We prove Theorem~\ref{thm:degd_homo_d'} using a similar approach as the proof of Theorem~\ref{thm:degd_homo}. 

%The proof of Theorem~\ref{thm:degd_homo_d'} builds on the proof of Theorem~\ref{thm:degd_homo}. To this end, we first establish the following lemma, based on Lemma~\ref{thm:degd_nonhomo}.

%The proof of Theorem~\ref{thm:degd_homo_d'} follows in the same fashion as the proof of Theorem~\ref{thm:degd_homo}. The main step 

The main step in the proof of Theorem~\ref{thm:degd_homo_d'} is the following lemma, which builds on Lemma~\ref{thm:degd_nonhomo}.

%builds on the proof of Theorem~\ref{thm:degd_homo}. To this end, we first establish the following lemma, based on Lemma~\ref{thm:degd_nonhomo}.



\begin{lemma}\label{thm:degd_nonhomo_d'}
Consider a (not necessarily homogeneous) polynomial~$p$ in $n$ variables and degree at most $d$, where $d$ is even. Let $d'$ be an even integer in $\{2,\ldots,d-2\}$. If $p(x^*)>0$ for some $x^*\in\reals^n$, then for $\hat{t}>0$ sufficiently large, $p(x)$ can be written as
\[p(x)=\left(\frac{p(x^*)}{2}-\hat{t}g_{d'}(x-x^*)\right)\left(1+g_{d-d'}(x-x^*)\right)+\hat{\sigma}(x),\]
where $\hat{\sigma}(x)\in\SOS_{n,d}$. In particular, the above statement holds for any 
% $t$ such that
\[\hat{t}\ge \tilde{T}(p;x^*)+\frac{p(x^*)}{2}.\]
Here, $g_{d'}(x), g_{d-d'}(x)$, and $\tilde{T}(p;x^*)$ are defined as in Lemma~\ref{thm:degd_nonhomo}.
% Consider a (not necessarily homogeneous) polynomial~$p$ in~$n$ variables and degree at most~$d$, where~$d$ is even. Suppose there exist scalars $A>0$ and $t>0$ such that $p(x)$ can be written as
% \[p(x)=\left(A-tg_d(x-x^*)\right)+\sigma(x),\]
% where $\sigma(x)\in\SOS_{n,d}$ and $g_d(x)$ is defined in Lemma~\ref{thm:degd_nonhomo}.
% Then, for any even degree $d'\in\{2,\ldots,d-2\}$, $p(x)$ can be written as
% \[p(x)=\left(A-\hat{t}g_{d'}(x-x^*)\right)\left(1+g_{d-d'}(x-x^*)\right)+\hat{\sigma}(x),\]
% where $\hat{t}=t+A$, and $\hat{\sigma}(x)\in\SOS_{n,d}$.
\end{lemma}

\begin{proof}
Without loss of generality, we assume $x^*=0$; otherwise, applying the statement of the lemma to the polynomial $p(x+x^*)$ and recalling that the existence of an sos decomposition is invariant under an affine change of variables, yields the desired result. By Lemma~\ref{thm:degd_nonhomo}, for any $t\ge \tilde{T}(p;0)$, $p(x)$ admits the representation
\[p(x)=\frac{p(0)}{2}-tg_d(x)+\sigma(x),\]
where $\sigma(x)\in\SOS_{n,d}$. Let $A=p(0)/2$ and $\hat{t}=t+A$. Then 
\[\begin{aligned}
p(x)&=A-tg_d(x)+\sigma(x)\\
&=\left(A-\hat{t}g_{d'}(x)\right)(1+g_{d-d'}(x))-Ag_{d-d'}(x)+\hat{t}g_{d'}(x)(1+g_{d-d'}(x))-tg_d(x)+\sigma(x)\\
&=\left(A-\hat{t}g_{d'}(x)\right)(1+g_{d-d'}(x)) + A\left(g_{d}(x)-g_{d-d'}(x)\right)\\
&\qquad +\hat{t}\big(g_{d'}(x)(1+g_{d-d'}(x))-g_d(x)\big)+ \sigma(x)\\
&=\left(A-\hat{t}g_{d'}(x)\right)(1+g_{d-d'}(x))+\hat{\sigma}(x),
\end{aligned}\]
where \[\hat{\sigma}(x)=\sigma(x)+A\left(g_{d}(x)-g_{d-d'}(x)\right)+\hat{t}\big(g_{d'}(x)(1+g_{d-d'}(x))-g_d(x)\big).\]
It remains to show that $\hat{\sigma}(x)\in\SOS_{n,d}$.
Observe that
\[g_{d}(x)-g_{d-d'}(x)=\sum_{k=(d-d')/2+1}^{d/2}\|x\|_2^{2k},\]
and
\[\begin{aligned}
   g_{d'}(x)(1+g_{d-d'}(x))-g_d(x)&=g_{d'}(x)+\left(g_{d'-2}(x)+\|x\|_2^{d'}\right)g_{d-d'}(x)-g_d(x)\\
% &=\sum_{k=1}^{d'/2}\|x\|_2^{2k}+\sum_{k=1}^{(d-d')/2}\|x\|_2^{2k+d'}-\sum_{k=1}^{d/2}\|x\|_2^{2k} +\left(g_{d'}(x)-\|x\|_2^{d'}\right)g_{d-d'}(x)\\
   % &=\left(g_{d'}(x)-\|x\|_2^{d'}\right)g_{d-d'}(x)\\
   &=g_{d'-2}(x)g_{d-d'}(x).
\end{aligned}\]
Hence, both expressions are sos, which implies $\hat{\sigma}(x)$ is sos.
\end{proof}



The proof of Theorem~\ref{thm:degd_homo_d'} follows the same approach as the proof of Theorem~\ref{thm:degd_homo} and is therefore only outlined. We first consider the special case where $x^* = e_1$. We dehomogenize~$p$ by defining $\tilde{p}(z)\defeq p(1,z)$, apply Lemma~\ref{thm:degd_nonhomo_d'} to $\tilde{p}$, and homogenize back to show that whenever  $t \geq T'(p; e_1)$, $p(x)$ can be written as 
\[
    p(x)=\left(\frac{p(e_1)}{2}(x^Te_1)^{d'}-th_{d'}(x;e_1)\right)\left((x^Te_1)^{d-d'}+h_{d-d'}(x;e_1)\right)+\sigma(x),
\]
where $\sigma(x) \in \HSOS_{n,d}$. We then establish the claim for a general point $x^* \in \mbS^{n-1}$ via an orthogonal change of variables, as in the proof of Theorem~\ref{thm:degd_homo}.


% \begin{proof}[Proof of Theorem~\ref{thm:degd_homo_d'}]
% \textcolor{blue}{We first consider the special case where $x^* = e_1$.
% By the same proof technique as in Theorem~\ref{thm:degd_homo},
% we dehomogenize~$p$ by defining $\tilde{p}(z)\defeq p(1,z)$, and apply Lemma~\ref{thm:degd_nonhomo_d'} to $\tilde{p}$, and homogenize back to show that whenever  $t \geq T'(p; e_1)$, $p(x)$ can be written as 
% \[
%     p(x)=\left(\frac{p(e_1)}{2}(x^Te_1)^{d'}-th_{d'}(x;e_1)\right)\left((x^Te_1)^{d-d'}+h_{d-d'}(x;e_1)\right)+\sigma(x),
% \]
% where $\sigma(x) \in \HSOS_{n,d}$.}

% \textcolor{blue}{We then establish the claim for a general point $x^* \in \mbS^{n-1}$ via an orthogonal change of variables as in the proof of Theorem~\ref{thm:degd_homo}.}

% Using a similar argument as in the proof of Theorem~\ref{thm:degd_homo}, one can extend Lemma~\ref{thm:degd_nonhomo_d'} to the homogeneous case, which gives the claimed result. 
% Details are analogous and omitted.
% Specifically, for a polynomial $p\in H_{n,d}$ with $d$ even, suppose there exists $A>0$ and $t>0$ such that $p(x)$ can be written as
% \[p(x)=\left(A(x^Tx^*)^d-th_d(x;x^*)\right)+\sigma(x),\]
% where $\sigma(x)\in\HSOS_{n,d}$ and $h_d(x;x^*)$ is defined in Theorem~\ref{thm:degd_homo}.
% Then, for any even degree $d'\in\{2,\ldots,d-2\}$, $p(x)$ can be written as
% \begin{equation}\label{eq:degd_homo_d'}
%     p(x)=\left(A(x^Tx^*)^{d'}-\hat{t}h_{d'}(x;x^*)\right)\left((x^Tx^*)^{d'}+h_{d-d'}(x;x^*)\right)+\hat{\sigma}(x),
% \end{equation}
% where $\hat{t}=t+A$, and $\hat{\sigma}(x)\in\HSOS_{n,d}$. Finally, by Theorem~\ref{thm:degd_homo}, when $t\ge T(p;x^*)$, $p(x)$ admits the representation
% \[p(x)=\left(\frac{p(x^*)}{2}(x^Tx^*)^d-th_d(x;x^*)\right)+\sigma(x).\]
% Taking $A=p(x^*)/2$ in~\eqref{eq:degd_homo_d'} yields the desired result.
% \end{proof}


\subsection{From Local Positivstellens\"atze to a Global Positivstellensatz}\label{sec:global_psatz}

% Let~$p\in H_{n,d}$ be positive definite, with~$d$ even. Then, for any $x^*\in \mbS^{n-1}$, Theorem~\ref{thm:degd_homo} and Theorem~\ref{thm:degd_homo_d'} guarantee the existence of an algebraic proof of nonnegativity of~$p$ over the spherical cap centered at $x^*$. Since the polar angle of each spherical cap is uniformly bounded below, the unit sphere can be covered by finitely many of them. 

Building on the results established in the previous subsection, we now present the proof of Theorem~\ref{cor:strictly_pos}.

% As a direct consequence of Theorem~\ref{thm:degd_homo} and~\ref{thm:degd_homo_d'}, if a polynomial $p\in H_{n,d}$ (with $d$ even) is positive definite, we can always find an algebraic disjunction $\mcD$ such that each subset in the collection $\mcD$ contains exactly one polynomial of degree less then or equal to $d$,
% and $p\in \DiSOS_{n,d}^{d}(\mcD)$. This is because for any $x_0\in \mbS^{n-1}$, we can find an algebraic proof of nonnegativity of $p$ over the spherical cap centered at $x^*$. Note that the polar angle of the spherical cap is uniformly lower bounded. Thus, we can select a finite number of such spherical caps to cover the unit sphere, thereby forming the desired algebraic disjunction. Based on the above observations, we now proceed to prove Theorem~\ref{cor:strictly_pos}.

% \begin{theorem}[Disjunctive Positivstellensatz]\label{cor:strictly_pos}
% For fixed dimension $n$ and even degree $d$, there exists an explicit family of algebraic disjunctions
% $\mcD^r$, indexed by $r$, such that for any positive
%  definite $p\in H_{n,d}$, we have $p\in \DiSOS_{n,d}(\mcD^r)$ for some $r$. Furthermore, each subregion in $\mcD^r$ is defined by a single polynomial of degree less then or equal to $d$, and the total degree of the identities in the disjunctive sos proof of nonnegativity of $p$ is exactly $d$.
% \end{theorem}

\begin{proof}[Proof of Theorem~\ref{cor:strictly_pos}]
% For any integer $r\ge 2$, we
% construct an $\epsilon$-net of the unit sphere, denoted by $\mcN_r\subset \mbS^{n-1}$, such that for any $y\in \mbS^{n-1}$, there exists $x\in \mcN_r$ with $\langle x,y\rangle\ge r/\sqrt{r^2+1}$. Write $\mcN_r=\{x^1,\ldots,x^{n_r}\}$, where  $n_r=|\mcN_r|$. 
For any integer $m\ge 2$, we
construct a net $\mcN_m=\{x^1,\ldots,x^{r_m}\}\subset\mbS^{n-1}$ such that for any $y\in \mbS^{n-1}$, there exists $x\in \mcN_m$ with $\langle x,y\rangle\ge m/\sqrt{m^2+1}$. 
% Here, $n_r=|\mcN_r|$ and $x^i\in\mbS^{n-1}$ for all $i=1,\ldots,n_r$. 
Fix an even degree $d'\in\{2,\ldots,d\}$ and let the initial algebraic disjunction be $\mcD_{d'}^1=\{\{0\}\}$. For $m\ge 2$, we construct the subsequent algebraic disjunctions as 
\begin{equation}\label{eq:strictly_pos_ad}
 \mcD^{m}_{d'}=\big\{\{(x^Tx^i)^{d'}-t_mh_{d'}(x;x^i)\},\quad i=1,\ldots,r_m\big\},
\end{equation}
where $t_m=\left(\sum_{k=1}^{d'/2}m^{-2k}\right)^{-1}$ and $h_{d}(x;x^i)$ is defined as in~\eqref{eq:degd_homo_h}. 
% Observe that for any $x\in\mbS^{n-1}$, 
% \[(x^Tx^i)^{d'}-t_rh_{d'}(x;x^i)\ge 0\quad\Leftrightarrow\quad \langle x,x^i\rangle\ge r/\sqrt{r^2+1}.\]
% This follows from the remark after Theorem~\ref{thm:degd_homo}. 
For any $x\in\mbS^{n-1}$ and any $x^i\in\mcN_m$, let $\theta\in[0,\pi]$ be the unique angle satisfying $\langle x,x^i\rangle=\cos\theta$. Then, it follows that
\[\begin{aligned}
    (x^Tx^i)^{d'}-t_mh_{d'}(x;x^i)&=(\cos\theta)^{d'}-t_m\sum_{k=1}^{d'/2}(\sin\theta)^{2k}(\cos\theta)^{d'-2k}\\
    &=(\cos\theta)^{d'}\left(1-t_m\sum_{k=1}^{d'/2}(\tan\theta)^{2k}\right).
\end{aligned}\]
Therefore, the inequality $(x^Tx^i)^{d'}-t_mh_{d'}(x;x^i)\ge 0$ holds if and only if 
% is equivalent to
\[
  \sum_{k=1}^{d'/2}(\tan\theta)^{2k}\le t_m^{-1}
  \quad\Leftrightarrow\quad |\tan\theta|\le m^{-1} \quad
  \Leftrightarrow\quad\langle x,x^i\rangle\ge \frac{m}{\sqrt{m^2+1}}. 
\]
Hence $\mcD_{d'}^{m}$ forms a valid algebraic disjunction of $\reals^n$ by the construction of the net. 

Finally, note that $t_m\to\infty$ as $m\to\infty$. 
Therefore, keeping in mind that $\mbS^{n-1}$ is a compact set, for any positive definite form $p\in H_{n,d}$, letting $p_{\min}\defeq \min_{x\in\mathbf{S}^{n-1}}p(x)$,  
there exists $m$ such that $t_m$ exceeds $2T(p;x^*)/p_{\min}$ (resp. $2T'(p;x^*)/p_{\min}$) for any $x^*\in\mbS^{n-1}$, where $T(p;x^*)$ (resp. $T'(p;x^*)$) is defined in Theorem~\ref{thm:degd_homo} (resp. Theorem~\ref{thm:degd_homo_d'}). 
This guarantees that $p\in \DiSOS_{n,d}^{d}(\mcD_{d'}^{m})$ {for} any fixed even degree $d'\in\{2,\ldots,d\}$.
\end{proof}

The preceding proof in fact yields an explicit bound on the number of subregions required to certify nonnegativity of a polynomial with a disjunctive sos certificate. As an illustration, we present such a bound for the case $d'=d$ in Theorem~\ref{thm:num_ad} below. Given the NP-hardness of certifying polynomial nonnegativity, any such bound is expected to grow exponentially with the dimension.

\begin{theorem}\label{thm:num_ad}
Consider a positive definite form $p\in H_{n,d}$. Let $p_{\min}\defeq\min_{x\in\mbS^{n-1}} p(x)$, and define
 \begin{equation}\label{eq:thm_counting}
T^*(p)\defeq \sup_{U\in\reals^{n\times n}:U^TU=I_n}\left( \frac{n}{2p_{\min}}\|p(Ux)\|_{\infty}^2 + \|p(Ux)\|_1\right),\quad
M(p)\defeq \left\lceil \sqrt{\frac{T^*(p)d}{p_{\min}}}\;\right\rceil.
\end{equation}
Let $\mcD_d^m$ be the family of algebraic disjunctions constructed in~\eqref{eq:strictly_pos_ad}, where each subset of $\mcD_d^m$ contains exactly one degree-$d$ form. Then $p\in\DiSOS_{n,d}^d(\mcD_d^m)$ for all $m\ge M(p)$. In particular, $\left(1+2 M(p)\right)^n$ subregions, each defined by a single degree-$d$ form, suffice to establish a disjunctive sos proof of nonnegativity of~$p$.
% when $r=\lceil R(p)\rceil$, 
% the minimal possible cardinality of $\mcD_d^r$ is at most $\left(3+2R(p)\right)^n$. 
\end{theorem}

\begin{proof}
Recall the definition of the algebraic disjunction
\[
 \mcD_d^{m}=\big\{\{(x^Tx^i)^{d}-t_mh_{d}(x;x^i)\},\quad i=1,\ldots,r_m\big\},   
\]
where $t_m=\left(\sum_{k=1}^{d/2}m^{-2k}\right)^{-1}$. For each $i=1,\ldots,r_m$, by Theorem~\ref{thm:degd_homo}, there exists a polynomial $\sigma_i\in\HSOS_{n,d}$ such that
 \[    p(x)=\left(\frac{p(x^i)}{2}(x^Tx^i)^d-T(p;x^i)h_d(x;x^i)\right)+\sigma_i(x),\]
 where $T(p;\cdot)$ is defined as in~\eqref{eq:degd_homo_t}. 
Rewriting gives
  \[\begin{aligned}    
  p(x)&=\frac{p(x^i)}{2}(x^Tx^i)^d+\sigma_i(x)-T(p;x^i)h_d(x;x^i)\\
  % &=\frac{p(x^i)}{2}\left((x^Tx^i)^d-t_mh_d(x;x^i)\right)+\sigma_i(x)\\
  &=\frac{p(x^i)}{2}\left((x^Tx^i)^d-t_mh_d(x;x^i)\right)+\left(\sigma_i(x)+\left(\frac{p(x^i)}{2}t_m-T(p;x^i)\right)h_d(x;x^i)\right).
  \end{aligned}\]
   Note that $t_m\ge 2m^2/d$. Therefore, whenever $m\ge M(p)$,
  \[\frac{p(x^i)}{2}t_m\ge \frac{p_{\min}}{2}\cdot\frac{2M^2(p)}{d}\ge T^*(p)\ge T(p;x^i),\quad i=1,\ldots,r_m.\]
 % for all $i=1,\ldots,r_m$. 
 Since $h_d(x;x^*)\in\HSOS_{n,d}$ for any $x^*\in\reals^n$, it follows that
 \[\sigma_i(x)+\left(\frac{p(x^i)}{2}t_m-T(p;x^i)\right)h_d(x;x^i)\in\HSOS_{n,d}.\]
 Thus, $p\in\DiSOS_{n,d}^d(\mcD_d^m)$ whenever $m\ge M(p)$. We now bound the size of $\mcN_m$ when $m=M(p)$. Recall that $\mcN_m=\{x^1,\ldots,x^{r_m}\}\subset\mbS^{n-1}$ is a net such that for any $y\in \mbS^{n-1}$, there exists $x\in \mcN_m$ with $\langle x,y\rangle\ge m/\sqrt{m^2+1}$. This implies that
 % For every $x\in\mbS^{n-1}$, there exists $y\in \mcN_r$ with $\langle x,y\rangle\ge r/\sqrt{r^2+1}$, which implies that
 \[\begin{aligned}
\|x-y\|_2^2&=2-2\langle x,y\rangle\le 2\left(1-\frac{m}{\sqrt{m^2+1}}\right)\\
&= \frac{2(\sqrt{m^2+1}-m)}{\sqrt{m^2+1}}=\frac{2}{(\sqrt{m^2+1}+m)\sqrt{m^2+1}}\\
&\le \frac{1}{m^2}.
 \end{aligned}\]
By \cite[Corollary 4.2.11]{vershynin2018high}, we conclude that 
$r_m\le (1+2m)^n= (1+2M(p))^n$.
\end{proof}

% In addition, we present the following theorem to demonstrate that in theory, we still may need exponential number of subregions to certify nonnegativity of a polynomial.

% \begin{theorem}\label{thm:num_ad}
% %Given dimension $n$,even degree $d$, and a positive definite
%  Given a positive definite form $p\in H_{n,d}$, let $p_{\min}\defeq\inf_{x\in\mbS^{n-1}} p(x)$, and
%  \[\begin{aligned}
% T_d(p)&\defeq\max\left\{ \sup_{U:U^TU=I_n}\left( \frac{n}{2p_{\min}}\|p(Ux)\|_{\infty}^2 + \|p(Ux)\|_1\right),\ p_{\min}\right\},\\    N(p)&\defeq\left(1+\frac{4T_d(p)}{p_{\min}}\right)^n.
%  \end{aligned}\]
%  Then, there exist $r\le N(p)$ points $x^1,\ldots,x^{r}\in\mbS^{n-1}$, and polynomials $\sigma_1,\ldots,\sigma_{r}\in\HSOS_{n,d}$, such that a disjunctive sum of squares proof of nonnegativity of $p$ can be given as 
%  \[    p(x)=\left(\frac{p(x^i)}{2}(x^Tx^i)^d-T_d(p)h_d(x;x^i)\right)+\sigma_i(x),\quad i=1,\ldots, r,\]
% where $h_d(x;x^i)$ is defined as in~\eqref{eq:degd_homo_h}.
% Consequently, we have $p\in\DiSOS_{n,d}^d(\mcD)$, where 
% \[ \mcD=\left\{\{p(x^i)(x^Tx^i)^{d}-2T_d(p)h_{d}(x;x^i)\},\quad i=1,\ldots,r\right\}.  \]   
% \end{theorem}

% \begin{proof}
% Let $\epsilon={p_{\min}}/{(2T_d(p))}$. 
% By~\cite[Lemma 5.2]{vershynin2018high}, there exists an $\epsilon$-net $\mcN$ of $\mbS^{n-1}$ such that $|\mcN|\le (1+2/\epsilon)^n= N(p)$. Let $|\mcN|=r$, and $\mcN=\{x^1,\ldots,x^r\}$. We claim that
% \[\mcD=\{\{(x^Tx^i)^{d}-T_d(p)h_{d}(x;x^i)\},\quad i=1,\ldots,r\}\]
% forms a valid algebraic disjunction of $\reals^n$. 

% By definition of an $\epsilon$-net, for any point $x\in\mbS^{n-1}$, there exists a point $x^i\in\mcN$ such that $\|x-x^i\|_2\le \epsilon$. Let the angle between $x$ and $x^i$ be $\theta$, then $\cos\theta=x^Tx^i$, and \[\|x-x^i\|_2^2=2-2\cos\theta=4\sin^2\frac{\theta}{2}.\] 
% Thus, we have
% \[\begin{aligned}
% \tan^2\theta&=\frac{\sin^2\theta}{\cos^2\theta}
% =\frac{4\sin^2\frac{\theta}{2}\cos^2\frac{\theta}{2}}{\left(1-\sin^2\frac{\theta}{2}\right)^2}\\
% &=\frac{4\sin^2\frac{\theta}{2}\left(1-\sin^2\frac{\theta}{2}\right)}{\left(1-2\sin^2\frac{\theta}{2}\right)^2}\\
% &\le \frac{\epsilon^2(1-\epsilon^2/4)}{(1-\epsilon^2/2)^2}\\
% &\le \frac{\epsilon}{1+\epsilon}.
% \end{aligned}\]
% One can verify numerically that the last inequality holds for all $\epsilon\in[0,1/2]$. Furthermore, 
% \[\begin{aligned}
% \epsilon(x^Tx^i)^d-h_d(x;x^i)&=\epsilon\cos^d\theta-\sum_{k=1}^{d/2}\sin^{2k}\theta\cos^{d-2k}\theta\\
% &=\cos^d{\theta}\left(\epsilon-\sum_{k=1}^{d/2}\tan^{2k}\theta\right)\\
% &=\cos^d{\theta}\left(\epsilon-\frac{\tan^2\theta-\tan^{d+2}\theta}{1-\tan^2\theta}\right)\\
% &\ge \cos^d{\theta}\left(\epsilon-\frac{\tan^2\theta}{1-\tan^2\theta}\right)\\
% &\ge 0
% \end{aligned}\]
% Therefore, the subregions defined by $p(x^i)(x^Tx^i)^{d}-2T_d(p)h_{d}(x;x^i)\ge 0$ for $i=1,\ldots,r$ cover $\mbS^{n-1}$ and, by homogeneity, also $\reals^n$. Now for each $i=1,\ldots,r$, since $T_d(p)\ge T_{d}(p;x^i)$, by Theorem~\ref{thm:degd_homo}, there exists a polynomial $\sigma_i\in\HSOS_{n,d}$ such that
%  \[    p(x)=\left(\frac{p(x^i)}{2}(x^Tx^i)^d-T_d(p)h_d(x;x^i)\right)+\sigma_i(x).\]
%  Thus, we complete the proof.

% \end{proof}

% \begin{theorem}\label{thm:lb_bound}
%  Given a polynomial $p\in H_{n,d}$, let $p_{\min}\defeq\inf_{x\in\mbS^{n-1}} p(x)$. Then for any positive integer $r$, there exists an algebraic disjunction $\mcD^{r}$ with $r$ subsets, such that 
% \[p_{\min}-p^{(r)}=O\left(r^{-\frac{1}{2n}}\right),\] 
% where 
%  \[\begin{array}{rl}
%     		p^{(r)}:=\max\limits_{\gamma\in{\scriptsize\reals}} & \gamma\\
% 		\emph{\st} & p(x)-\gamma\|x\|_2^d\in\DiSOS_{n,d}^{d}(\mcD^r).

%  \end{array}\]
% \end{theorem}
% \begin{proof}
% Let $q(x)=\|x\|_2^d$. Given $\epsilon<p_{\min}$, by Theorem~\ref{thm:num_ad}, we have that $ p^{(k)}\ge \epsilon$ 
% for all $k\ge N(p-\epsilon q)$. Note that 
% \[\begin{aligned}
% T_d(p-\epsilon q)&\le \sup_{U:U^TU=I_n}\left( \frac{n}{2(p_{\min}-\epsilon)}\|(p-\epsilon q)(Ux)\|_{\infty}^2 + \|(p-\epsilon q)(Ux)\|_1\right) \\
% &\le \sup_{U:U^TU=I_n}\left( \frac{n}{p_{\min}-\epsilon}\left(\|p(Ux)\|_{\infty}^2+\epsilon^2\|q(Ux)\|_{\infty}^2\right) + \|p(Ux)\|_1+\epsilon \|q(Ux)\|_1\right)\\
% &\le \frac{2p_{\min}}{p_{\min}-\epsilon} T_d(p)+\frac{c_1n\epsilon^2}{p_{\min}-\epsilon}+c_2\epsilon,
% \end{aligned}\]
% where $c_1=\|q\|_{\infty}^2,c_2=\|q\|_1$ are constants independent of $p$. Therefore,
% \[\begin{aligned}
% N(p-\epsilon q)&=\left(1+\frac{4T_d(p-\epsilon q)}{p_{\min}-\epsilon}\right)^n\\
% &\le \left(1+\frac{8p_{\min}T_d(p)+4c_1n\epsilon^2+4c_2\epsilon(p_{\min}-\epsilon)}{(p_{\min}-\epsilon)^2}\right)^n\\
% &\le \left(1+\frac{8p_{\min}T_d(p)+(4c_1n+c_2)p_{\min}^2}{(p_{\min}-\epsilon)^2}\right)^n.    
% \end{aligned}\]
% Let $\ell$ be the upper bound of $N(p-\epsilon q)$ above. We have 
% \[p_{\min}-\epsilon\le \sqrt{\frac{8p_{\min}T_d(p)+(4c_1n+c_2)p_{\min}^2}{\ell^{1/n}-1}}.\]
% Therefore, we conclude that for any positive integer $\ell$,
% \[0\le p_{\min}- p^{(\ell)} \le \sqrt{\frac{8p_{\min}T_d(p)+(4c_1n+c_2)p_{\min}^2}{\ell^{1/n}-1}}.\]
% \end{proof}

% \subsection{Spherical-Caps Covering Algorithm for Polynomial Minimization} \label{sec:disos_grid}
\subsection{An Algorithm Based on Spherical-Cap Subregions} \label{sec:disos_grid}
Building on Theorem~\ref{cor:strictly_pos}, we propose an algorithm to approximate the minimum value $p_{\min}$ of a form $p\in H_{n,d}$ over the unit sphere $\mbS^{n-1}$. This quantity is also equal to the optimal value of the following optimization problem:
\begin{equation}\label{eq:poly_min}
	p_{\min} =\; \begin{array}[t]{cl}
		\max\limits_{\gamma\in{\scriptsize\reals}} & \gamma\\
		\st & p(x)-\gamma\|x\|_2^d\ge 0,\quad\forall x\in\reals^n.
	\end{array}
\end{equation}
For any even degree $d'\in\{2,\ldots,d\}$, we can present a complete SDP-based hierarchy of lower bounds on $p_{\min}$. At level $m$ of this hierarchy, we construct a net $\mcN_m$ of the unit sphere\footnote{There are different techniques for constructing such a net; see, e.g.,~\cite{leopardi_partition_2006,chakraborty2022generating}.} and the corresponding algebraic disjunction $\mcD_{d'}^m$ defined in~\eqref{eq:strictly_pos_ad}. Using the notation introduced in Definition~\ref{def:disos}, we compute a lower bound $p^{(m)}$ by solving the following SDP:
\begin{equation}\label{eq:poly_min_disos_relax}
	p^{(m)} \defeq \begin{array}[t]{cl}
		\max\limits_{\gamma\in{\scriptsize\reals}} & \gamma\\
		\st & p(x)-\gamma\|x\|_2^d\in\DiSOS_{n,d}^{d}(\mcD_{d'}^m).
	\end{array}
\end{equation}
{This SDP can be decomposed into $r_m$ independent SDPs, one for each subset in $\mcD_{d'}^m$, and solved in parallel; see Remark~\ref{remark:uniform_grid} for details.} 
The next theorem utilizes Theorem~\ref{thm:num_ad} to establish a convergence rate for the sequence of lower bounds $p^{(m)}$ produced by the SDP hierarchy in~\eqref{eq:poly_min_disos_relax} (in the case $d'=d$ as a representative example).

% that this hierarchy produces increasingly tight lower bounds on $p_{\min}$ as the level $m$ increases.

\begin{theorem}\label{thm:poly_min_disos_lb}
Let~$p\in H_{n,d}$ and~$p_{\min}\defeq\min_{x\in\mbS^{n-1}} p(x)$. Let $\mcD^m_d$ be the family of algebraic disjunctions constructed in~\eqref{eq:strictly_pos_ad}, where each subset of $\mcD_d^m$ contains exactly one degree-$d$ form. 
% For each positive integer $m$, let $p^{(m)}$ be the optimal value of problem~\eqref{eq:poly_min_disos_relax}.
Then, we have
\[p_{\min}-p^{(m)}= O(1/m).\] 
\end{theorem}
\begin{proof}
Let $q(x)\defeq\|x\|_2^d$ and fix $\epsilon\in(0,p_{\min})$. By Theorem~\ref{thm:num_ad}, we have $p-\epsilon q\in\DiSOS_{n,d}^d(\mcD_d^m)$ for all $m\ge M(p-\epsilon q)$, where $M(\cdot)$ is defined in~\eqref{eq:thm_counting}. Thus, $ p^{(m)}\ge \epsilon$ 
whenever $m\ge M(p-\epsilon q)$. Let $c_{\infty}\defeq \|q\|_{\infty},c_1\defeq\|q\|_1$, which are constants depending only on $n$ and $d$. We now bound $T^*(p - \epsilon q)$ from above using the definition in~\eqref{eq:thm_counting}:
\[\begin{aligned}
T^*(p-\epsilon q)&= \sup_{U\in\reals^{n\times n}:U^TU=I_n}\left( \frac{n}{2(p_{\min}-\epsilon)}\|(p-\epsilon q)(Ux)\|_{\infty}^2 + \|(p-\epsilon q)(Ux)\|_1\right) \\
&\le \sup_{U\in\reals^{n\times n}:U^TU=I_n}\left( \frac{n}{p_{\min}-\epsilon}\left(\|p(Ux)\|_{\infty}^2+\epsilon^2\|q(Ux)\|_{\infty}^2\right) + \|p(Ux)\|_1+\epsilon \|q(Ux)\|_1\right)\\
&= \sup_{U\in\reals^{n\times n}:U^TU=I_n}\left( \frac{n}{p_{\min}-\epsilon}\left(\|p(Ux)\|_{\infty}^2+c_{\infty}^2\epsilon^2\right) + \|p(Ux)\|_1+c_1\epsilon\right)\\
&\le \frac{2p_{\min}}{p_{\min}-\epsilon} T^*(p)+\frac{c_{\infty}^2n\epsilon^2}{p_{\min}-\epsilon}+c_1\epsilon\\
&= \frac{2p_{\min}T^*(p)+c_{\infty}^2n\epsilon^2+c_1\epsilon(p_{\min}-\epsilon)}{p_{\min}-\epsilon}\\
&\le \frac{2p_{\min}T^*(p)+(c_{\infty}^2n+c_1/4)p_{\min}^2}{p_{\min}-\epsilon}.
\end{aligned}\]
Here, in the first inequality we used the triangle inequality for the polynomial norms. 
Therefore, for any $\epsilon\in(0,p_{\min})$, 
\[M(p-\epsilon q)\le \sqrt{\frac{2T^*(p-\epsilon q)d}{p_{\min}-\epsilon}}+1\le \frac{\sqrt{2(2p_{\min}T^*(p)+(c_{\infty}^2n+c_1/4)p_{\min}^2)d}}{p_{\min}-\epsilon}+1.\]
Finally, since the above bound holds for every $\epsilon\in(0,p_{\min})$, when~$m$ is sufficiently large, we may choose $\epsilon\in(0,p_{\min})$ such that
\[m = \frac{\sqrt{2(2p_{\min}T^*(p)+(c_{\infty}^2n+c_1/4)p_{\min}^2)d}}{p_{\min}-\epsilon}+1.\]
With this choice, it follows that
\[p_{\min}-p^{(m)}\le p_{\min}-\epsilon= \frac{\sqrt{2(2p_{\min}T^*(p)+(c_{\infty}^2n+c_1/4)p_{\min}^2)d}}{m-1}=O(1/m).\]
\end{proof}

Since the net $\mcN_m$ already appears in the formulation of our SDP, it is natural to use the points in $\mcN_m$ to obtain upper bounds on $p_{\min}$ via direct point evaluation. 
As the level $m$ increases, the resolution of the net improves and its covering radius tends to zero, yielding increasingly tight upper bounds on $p_{\min}$. In particular, by adapting the argument in the comments following~\cite[Proposition 4]{vianello_subperiodic_2018}, one can construct a net such that the error between the resulting upper bound at level $m$ and $p_{\min}$ is of order $O(1/m^2)$. Combined with Theorem~\ref{thm:poly_min_disos_lb}, this implies that the gap between the lower and upper bounds on $p_{\min}$ converges to zero as $m$ increases, with an overall rate of $O(1/m)$. Therefore, for any prescribed tolerance $\epsilon_{\text{tol}}>0$ to approximate $p_{\min}$, Algorithm~\ref{alg:uniform_grid} is guaranteed to terminate within $O(1/\epsilon_{\text{tol}})$ iterations. 

% To the best of our knowledge, however, there is no scalable implementation that deterministically constructs such nets with increasing resolution. Although Algorithm~\ref{alg:uniform_grid} is theoretically sound, it quickly becomes inefficient in practice. To improve tractability, we propose an alternative algorithm in Section~\ref{sec:bnb} for computing $p_{\min}$ that builds on the spatial branch-and-bound framework.

% While Algorithm~\ref{alg:uniform_grid} is theoretically sound, 


\begin{algorithm}[!tb]
\caption{Spherical-Cap-Based Algorithm for Minimizing a Form}
\label{alg:uniform_grid}

\newlength{\algmathwidth}
\setlength{\algmathwidth}{\dimexpr\linewidth-\algorithmicindent\relax}
\begin{algorithmic}
\Require A polynomial $p\in H_{n,d}$, an even number $d'\in\{2,\ldots,d\}$, a tolerance $\epsilon_{\text{tol}}>0$.
\Ensure  A lower bound $L$ on the minimum value of $p$ over $\mbS^{n-1}$ and a point $\bar{x}\in\mbS^{n-1}$ such that 
$p(\bar{x})-L\le \epsilon_{\text{tol}}$.
\Initialization Let $L=-\infty, U=\infty, m=1$.
\While{$U-L> \epsilon_{\text{tol}}$}
\State Generate a net $\mcN_m=\{x^1,\ldots,x^{r_m}\}\subset\mbS^{n-1}$ such that for all $y\in \mbS^{n-1}$, there exists\\
\hspace*{\algorithmicindent} $x\in \mcN_m$ with $\langle x,y\rangle\ge m/\sqrt{m^2+1}$.
\State $t_m\gets\left(\sum_{k=1}^{d'/2}m^{-2k}\right)^{-1}$.
\State 
$L^m\gets$ \parbox[t]{\algmathwidth}{%
$\begin{array}[t]{cl}
 \max\limits_{\gamma,s_i,\hat{s}_i} & \gamma\\
 \st & p(x)-\gamma\|x\|_2^d= s_{i}(x)+\left((x^Tx^i)^{d'}-t_mh_{d'}(x;x^i)\right)\hat{s}_{i}(x),\quad i=1,\ldots,r_m\\
 & s_{i}\in\HSOS_{n,d},\hat{s}_{i}\in\HSOS_{n,d-d'},\quad i=1,\ldots,r_m.
\end{array}$
}
\State $L\gets\max\{L,L^m\}$.
\State $U^m\gets\min_{1\le i\le r_m}\{p(x^i)\}$, $j\gets\argmin_{1\le i\le r_m}\{p(x^i)\}$
\Comment{Break ties arbitrarily.}
\If{$U^m < U$} $U\gets U^m$, $\bar{x}\gets x^{j}$.
\EndIf
\State $m\gets m+1$.
\EndWhile
\State \Return{$L, \bar{x}$}
\end{algorithmic}
\end{algorithm}

\begin{remark}\label{remark:uniform_grid}
{The computation of $L^m$ in Algorithm~\ref{alg:uniform_grid} can be decomposed into $r_m$ independent SDPs that can be solved in parallel.
Indeed, for each $k\in\{1,\ldots,r_m\}$, one can compute
\begin{equation*}
\gamma_k^*\defeq
\begin{array}[t]{cl}
\max\limits_{\gamma_k,s_k,\hat{s}_k} & \gamma_k\\
\st & p(x)-\gamma_k\|x\|_2^d= s_{k}(x)+\left((x^Tx^k)^{d'}-t_mh_{d'}(x;x^k)\right)\hat{s}_{k}(x)\\
& s_{k}\in\HSOS_{n,d},\ \hat{s}_{k}\in\HSOS_{n,d-d'},
\end{array}
\end{equation*}
and set $L^m\gets\min_{1\le k\le r_m}\gamma_k^*$.}
\end{remark}


% In practice, one can further improve efficiency by refining the disjunction 

In practice, it is often more efficient to refine the subregions of the unit sphere adaptively instead of increasing the resolution of the net uniformly. Such an adaptive algorithm can be implemented within a spatial branch-and-bound framework, as discussed in Section~\ref{sec:bnb}.

% \begin{remark}
% From the remark after the proof of Theorem~\ref{thm:degd_homo}, the computation of $L^r$ can be implemented using linear programming combined with bisection when $d'=d$. Concretely, for a fixed scalar $\gamma$ and a fixed $\epsilon$-net $\mcN_r=\{x^1,\ldots,x^{n_r}\}$, we test if $\gamma$ is a valid lower bound on~$p$ as follows. For each $i=1,\ldots,n_r$, we first check that $p(x^i)>\gamma$. We then search for a positive scalar $t_i$ such that
% \[p(x)-\gamma\|x\|_2^d-\left(\frac{p(x^i)-\gamma}{2}(x^Tx^i)^d-t_ih_d(x;x^i)\right)\]
% can be written as the sum of an explicit sos polynomial and a dsos polynomial with respect to an appropriate monomial basis as guaranteed by the remark. The search for such a scalar $t_i$ and the dsos polynomial can be formulated as a linear program. If such $t_i$ and dsos polynomial exist for every $i$, then $\gamma$ is certified as a valid lower bound. Finally, bisection on $\gamma$ yields the largest lower bound that can be certified with the current $\epsilon$-net.
% \end{remark}
\section{An Optimization-Free Disjunctive Positivstellensatz}\label{sec:disos_new}
In this section, we introduce a new disjunctive approach to certify polynomial nonnegativity. Unlike our previous approach, which relies on sos polynomials and {solutions} to SDPs, the new approach only requires checking nonnegativity of the coefficients of some polynomials we can explicitly write down. In this framework, \emph{checking} polynomial nonnegativity does not require any optimization solver at all, and \emph{imposing} a nonnegativity constraint on a polynomial leads to a linear program. Our main result in this section is an optimization-free disjunctive Positivstellensatz (Theorem~\ref{cor:strictly_pos_new}), analogous to Theorem~\ref{cor:strictly_pos}. 
% We note that 
% Specifically, we construct a family of special algebraic disjunctions and show that every positive definite form satisfies the sufficient condition for some disjunction in this family.
%We note that the results of this section are mainly of theoretical interest, as the number of polynomials whose coefficients should be checked can be prohibitively large in practice. 




% \subsection{Simplicial Disjunctions}
\subsection{Simplicial Disjunctions and Optimization-Free Certificates}
% In Section~\ref{sec:disos}, we introduce the concept of a disos proof of nonnegativity, which involves solving multiple small, fixed-size SDPs. In this section, we present an optimization-free approach that certifies nonnegativity of polynomials. Specifically, we utilize a particular class of algebraic disjunctions $\mcD=\{\{q_{k,j}\}_{j=1,\ldots, n_k},\quad k=1,\ldots,r\}$, where all polynomials $q_{k,j}$ are linear. We refer to this class as simplicial disjunctions.
% Definition~\ref{def:disjunction} introduces algebraic disjunctions in Definition~\ref{def:disjunction}, 
% Recall the notion of an algebraic disjunction $\mcD=\{\{q_{k,j}\}_{j=1,\ldots, n_k},\quad k=1,\ldots,r\}$ in Definition~\ref{def:disjunction}. 
% written as $\mcD=\{\{q_{k,j}\}_{j=1,\ldots, n_k},\quad k=1,\ldots,r\}$. 
% We now consider a specific subclass
In this section, we work with algebraic disjunctions where the corresponding subregions of the space are polyhedral cones. 
Recall that a polyhedral cone in $\reals^n$ can be represented as a conic combination of a finite number of vectors $v_1,\ldots,v_{\ell}\in\reals^n$: 
\[\cone\{v_1,\ldots,v_{\ell}\}=\{x\in\reals^n\mid x=\alpha_1v_1+\ldots+\alpha_{\ell}v_{\ell}\;\text{ for some }\alpha_1,\ldots,\alpha_{\ell}\ge 0\}.\]
We say that a polyhedral cone is simplicial if it is generated by linearly independent vectors.

\begin{definition}
Let $r$ be a positive integer. We say that a collection of $n\times n$ invertible matrices $\mcV=\{V_k\}_{k=1,\ldots,r}$ forms a \emph{simplicial disjunction} of $\reals^n$ if 
\[\mathbb{R}^n=\bigcup\limits_{k=1}^r\Omega_k,\quad \text{where }\quad\Omega_k=\cone\{V_k\}\defeq\cone\{v_{k,1}, \ldots, v_{k,n}\},\]
and $v_{k,j}$ is the $j$-th column of $V_k$.
\end{definition}

Note that any vector $x\in\cone\{V_k\}$ can be written as $x=V_kz$ for some elementwise nonnegative vector $z\in\reals^n$. Since $V_k$ is invertible, this implies that $$\cone\{V_k\}=\{x\in\reals^n\mid V_k^{-1}x\ge 0\}.$$ 
Therefore, the simplicial disjunction $\mcV$ can be viewed as a special case of an algebraic disjunction introduced in Definition~\ref{def:disjunction} with $n_k=n$ for $k=1,\ldots,r$, and polynomials $q_{k,j}(x)=c_{k,j}^Tx$ for $k=1,\ldots,r,j=1,\ldots,n$, where $c_{k,j}\in\reals^n$ are the columns of $V_k^{-T}$. 
% , i.e., 
% $V_k^{-T}=[c_{k,1}\ \ldots\ c_{k,n}]$.
We now explain how simplicial disjunctions can help certify polynomial nonnegativity. Because of the homogenization argument given in Section~\ref{sec:local_psatz}, we restrict attention to forms.
% Building on the notion of simplicial disjunctions, we now formalize a sufficient condition for certifying polynomial nonnegativity.
\begin{definition}\label{def:disos_new}
A polynomial~$p\in H_{n,d}$ is said to \emph{have nonnegative coefficients with respect to a simplicial disjunction} $\mcV=\{V_k\}_{k=1,\ldots,r}$ if $p(V_kx)$ has nonnegative coefficients for all $k=1,\ldots, r$. We denote the set of such polynomials by $\DiSOSL_{n,d}(\mcV)$. 
\end{definition}
Analogous to the cone $\DiSOS_{n,d}^{\bar{d}}(\mcD)$ 
in Definition~\ref{def:disos}, the set $\DiSOSL_{n,d}(\mcV)$ is a convex cone.  
Moreover, any polynomial $p\in\DiSOSL_{n,d}(\mcV)$ has an explicit disjunctive proof of nonnegativity. Indeed, for $k=1,\ldots,r$, let $\sigma_k(x)\defeq p(V_k x)$, which, by definition, has nonnegative coefficients. We can therefore write 
\[p(x)=\sigma_k(V_k^{-1}x),\quad k=1,\ldots,r.\]
Since $\sigma_k$ has nonnegative coefficients, we have the implication
\[V_k^{-1}x\ge 0\ \Rightarrow\ \sigma_k(V_k^{-1}x)\ge 0,\]
for all $k=1,\ldots,r$. Equivalently, $x\in\cone\{V_k\}\ \Rightarrow\ p(x)\ge 0$. Since $\bigcup_{k=1}^r\cone\{V_k\}=\reals^n$ by the definition of a simplicial disjunction,~$p$ is nonnegative over $\reals^n$.

We now establish the following optimization-free disjunctive Positivstellensatz, which runs parallel to the disjunctive Positivstellensatz in Theorem~\ref{cor:strictly_pos}. It shows that every positive definite form has a disjunctive proof of nonnegativity of the above type.
\begin{theorem}[Optimization-Free Disjunctive Positivstellensatz]\label{cor:strictly_pos_new}
For any fixed dimension $n$ and degree $d$, there exists an explicit family of simplicial disjunctions
$\mcV^m$, indexed by $m$, such that for any positive
 definite form $p\in H_{n,d}$, we have $p\in \DiSOSL_{n,d}(\mcV^m)$ for some $m$.
\end{theorem}
\subsection{A Local Optimization-Free Positivstellensatz}
 To prove Theorem~\ref{cor:strictly_pos_new}, we first establish {a} local optimization-free Positivstellensatz in Theorem~\ref{thm:degd_homo_new}, which runs parallel to Theorem~\ref{thm:degd_homo}. Recall that $\|\cdot\|_{\infty}$ denotes the largest coefficient in absolute value of a polynomial in the standard monomial basis. For an $n$-dimensional multi-index $\alpha$ with $|\alpha|=d$, we define
 \[\binom{d}{\alpha}\defeq \frac{d!}{\alpha_1!\ldots\alpha_n!}.\]
We denote the $n$-dimensional vector of all ones by $\mathbf{1}$.



% In Section~\ref{sec:disos_psatz}, we present a disjunctive Positivstellensatz, motivated by continuity of polynomials. With the same spirit, we establish the following optimization-free disjunctive Positivstellensatz. 
% Consider the polynomial $q(x;V)=p(Vx)$, parameterized by a matrix $V\in\reals^{n\times n}$. When $V=x^*\mathbf{1}^T$, i.e., all columns of $V$ are $x^*$, we have
% \[ 
% q(x;V)=p(x^*\mathbf{1}^Tx)=(\mathbf{1}^Tx)^d p(x^*)= p(x^*)\left(\sum_{i=1}^n x_i\right)^d.\]
% It is a homogeneous polynomial whose coefficient of each monomial of degree $d$ is positive. Since the coefficients of $q(x;V)$ are continuous in the entries of $V$, they remain nonnegative when all columns of $V$ are sufficiently close to $x^*$. The following theorem provides a quantitative measure of this closeness.

\begin{theorem}\label{thm:degd_homo_new}
For a polynomial $p\in H_{n,d}$ with $d$ even, if $p(x^*)>0$ for some $x^*\in \mbS^{n-1}$, then there exists $\delta >0$ such that for any matrix $V\in\reals^{n\times n}$ with $\|V-x^*\mathbf{1}^T\|_{F}\le \delta$, the polynomial $p(Vx)$ has nonnegative coefficients.
In particular, the above statement holds for any
\[\delta \le \delta(p;x^*)\defeq \sup_{\epsilon>1/p(x^*)}\min\left\{\frac{(p(x^*)\epsilon)^{1/d}-1}{( p(x^*)\epsilon)^{1/d}+1},\frac{1}{(1+\eta(p)\epsilon)n}\right\},\]
where
\[\eta(p)\defeq\sup_{U\in\reals^{n\times n}:U^TU=I_n} \|p(Ux)\|_{\infty}.\]
\end{theorem}

\begin{remark}
The polynomial $p(Vx)$ having nonnegative coefficients implies that $p(x)\ge 0$ for all $x\in\cone\{V\}$. In particular, when $V$ is invertible, $\cone\{V\}$ is simplicial and full-dimensional, which provides a proof of nonnegativity of~$p$ over a full-dimensional neighborhood around~$x^*$.
\end{remark}

% The proof of Theorem~\ref{thm:degd_homo_new} follows a similar approach to that of Theorem~\ref{thm:degd_homo}. We begin by considering the case where $x^*=e_1$, and then generalize the result to an arbitrary point $x^*\in \mbS^{n-1}$ using a rotation argument.
% \begin{lemma}\label{lem:degd_homo_new_e1}
% For a polynomial $p\in H_{n,d}$ with $d$ even, if $p(e_1)>0$, then there exists $\delta >0$, such that for any matrix $V\in\reals^{n\times n}$ satisfying $\|V-e_1\mathbf{1}^T\|_{F}\le \delta$, the polynomial $p(Vx)$ has nonnegative coefficients.
% In particular, the above statement holds for any
% \[\delta \le \delta(p;e_1)=\sup_{\epsilon>1/p(e_1)}\min\left\{\frac{(\epsilon p(e_1))^{1/d}-1}{(\epsilon p(e_1))^{1/d}+1},\frac{1}{(1+\epsilon)n}\right\}.\]
% \end{lemma}

\begin{proof}
We first consider the special case where $x^*=e_1$. Let $p(x)=\sum_{\alpha:|\alpha|=d}p_{\alpha}x^{\alpha}$. For each multi-index $\beta$ with $|\beta|=d$, we estimate the coefficient of the monomial $x^{\beta}$ in the expansion $p(Vx)=\sum_{\alpha:|\alpha|=d}p_{\alpha}(Vx)^{\alpha}$. 
% For each multi-index $\alpha$ with $|\alpha|=d$, 
Let $c(\alpha,\beta)$ denote the coefficient of $x^{\beta}$ arising from the expansion of $(Vx)^{\alpha}$. Then, the coefficient of $x^{\beta}$ in $p(Vx)$ is given by $$\sum_{\alpha:|\alpha|=d}p_{\alpha}c(\alpha,\beta).$$
In the case where $\alpha\neq de_1$ (i.e., $x^{\alpha}\neq x_1^d$), we bound $c(\alpha,\beta)$ by first noting that $$(Vx)^{\alpha}=\prod_{k=1}^n(V_{k1}x_1+\ldots+V_{kn}x_n)^{\alpha_k}.$$ 
Since $\|V-e_1\mathbf{1}^T\|_{F}\le \delta$, it follows that $V_{1j}\in[1-\delta,1+\delta]$ for $j=1,\ldots, n$, and $V_{kj}\in[-\delta,\delta]$ for $k=2,\ldots, n$ and $j=1,\ldots, n$. Hence, $\max_{1\le j\le n}|V_{1j}|\le 1+\delta$ and $\max_{1\le j\le n}|V_{kj}|\le \delta$ for  $k=2,\ldots,n$. 
Therefore, recalling the multinomial theorem, we have
\begin{equation}\label{eq:optfree_b1}
\begin{aligned}
|c(\alpha,\beta)|&\le \sum_{\substack{\gamma_1+\ldots+\gamma_n=\beta \\ |\gamma_i|=\alpha_i,\; i=1,\ldots,n}}\prod_{k=1}^n \left(\max_{1\le j\le n}|V_{kj}|\right)^{\alpha_k}\binom{\alpha_k}{\gamma_k}\\
&\le (1+\delta)^{\alpha_1}\delta^{d-\alpha_1}\sum_{\substack{\gamma_1+\ldots+\gamma_n=\beta \\ |\gamma_i|=\alpha_i,\; i=1,\ldots,n}}\prod_{k=1}^n \binom{\alpha_k}{\gamma_k}\\
&=(1+\delta)^{\alpha_1}\delta^{d-\alpha_1} \binom{d}{\beta},
\end{aligned}    
\end{equation}
%where we have used the multinomial theorem in the first and third lines.
The last equality follows from matching the coefficients of $x^{\beta}$ on both sides of the identity $(x_1+\dots+x_n)^d=\prod_{k=1}^n(x_1+\ldots+x_n)^{\alpha_k}$ and employing the multinomial theorem. 

We now turn to the case $\alpha=de_1$, where we have $(Vx)^{\alpha}=(V_{11}x_1+\ldots+V_{1n}x_n)^d$. For any $\delta\le \delta(p;e_1)$, we have $\delta < 1/n\le 1$. Thus, $\min_{1\le j\le n} V_{1j}\ge1-\delta\ge 0$, and
\begin{equation}\label{eq:optfree_b2}
c(\alpha,\beta)\ge \left(\min_{1\le j\le n}V_{1j}\right)^{d}\binom{d}{\beta}\ge (1-\delta)^d\binom{d}{\beta}.
\end{equation}
Combining the bounds~\eqref{eq:optfree_b1} for $\alpha\neq de_1$ and~\eqref{eq:optfree_b2} for $\alpha=de_1$, the coefficient of $x^{\beta}$ in $p(Vx)$ can be bounded from below as
% \[\begin{aligned}
% \sum_{\alpha:|\alpha|=d}p_{\alpha}c(\alpha,\beta)
% &=\sum_{l=0}^d\sum_{\alpha:|\alpha|=d,\alpha_1=l}p_{\alpha}c(\alpha,\beta)\\
% &\ge \binom{d}{\beta}\left(p(e_1)(1-\delta)^d-\|p\|_{\infty}\sum_{l=0}^{d-1}\sum_{\alpha:|\alpha|=d,\alpha_1=l}(1+\delta)^{\alpha_1}\delta^{d-\alpha_1}\right).
% \end{aligned}\]
% \[\begin{aligned}
% \sum_{l=0}^{d-1}\sum_{\alpha:|\alpha|=d,\alpha_1=l}(1+\delta)^{\alpha_1}\delta^{d-\alpha_1}&\le\sum_{l=0}^{d-1}n^{d-l}(1+\delta)^{l}\delta^{d-l}\\
% &=\left(\frac{1+\delta}{n\delta}-1\right)^{-1}\left((1+\delta)^d-n^{d}\delta^d\right)\\
% &\le\left(\frac{1}{n\delta}-1\right)^{-1}(1+\delta)^d.
% \end{aligned}\]
% Therefore,
% \[
% \sum_{\alpha:|\alpha|=d}p_{\alpha}c(\alpha,\beta)\ge\binom{d}{\beta}\left(\frac{1}{n\delta}-1\right)^{-1}\left(\left(\frac{1}{n\delta}-1\right)p(e_1)(1-\delta)^d-\|p\|_{\infty}(1+\delta)^d\right).
% \]
\[\begin{aligned}
\sum_{\alpha:|\alpha|=d}p_{\alpha}c(\alpha,\beta)&=p_{de_1}c(de_1,\beta)+\sum_{\alpha:|\alpha|=d,
\alpha \neq de_1}p_{\alpha}c(\alpha,\beta)\\
&=p(e_1)c(de_1,\beta)+\sum_{l=0}^{d-1}\sum_{\alpha:|\alpha|=d,\alpha_1=l}p_{\alpha}c(\alpha,\beta)\\
&\ge \binom{d}{\beta}\left(p(e_1)(1-\delta)^d-\|p\|_{\infty}\sum_{l=0}^{d-1}\sum_{\alpha:|\alpha|=d,\alpha_1=l}(1+\delta)^{\alpha_1}\delta^{d-\alpha_1}\right)\\
&\ge\binom{d}{\beta}\left(p(e_1)(1-\delta)^d-\|p\|_{\infty}\sum_{l=0}^{d-1}n^{d-l}(1+\delta)^{l}\delta^{d-l}\right)\\
&=\binom{d}{\beta}\left(p(e_1)(1-\delta)^d-\|p\|_{\infty}\left(\frac{1+\delta}{n\delta}-1\right)^{-1}\left((1+\delta)^d-n^{d}\delta^d\right)\right)\\
&\ge\binom{d}{\beta}\left(p(e_1)(1-\delta)^d-\|p\|_{\infty}\left(\frac{1}{n\delta}-1\right)^{-1}(1+\delta)^d\right)\\
% &=\binom{d}{\beta}\left(\frac{1}{n\delta}-1\right)^{-1}\left(\left(\frac{1}{n\delta}-1\right)p(e_1)(1-\delta)^d-\|p\|_{\infty}(1+\delta)^d\right),
\end{aligned}\]
% \[\begin{aligned}
% \sum_{\alpha:|\alpha|=d}p_{\alpha}c(\alpha,\beta)
% &\ge \binom{d}{\beta}\left(p(e_1)(1-\delta)^d-\sum_{l=0}^{d-1}\sum_{\substack{\alpha:|\alpha|=d\\\alpha_1=l}}(1+\delta)^{\alpha_1}\delta^{d-\alpha_1}\right)\\
% &\ge\binom{d}{\beta}\left(p(e_1)(1-\delta)^d-\sum_{l=0}^{d-1}n^{d-l}(1+\delta)^{l}\delta^{d-l}\right)\\
% % &=\binom{d}{\beta}\left(p(e_1)(1-\delta)^d-
% % \frac{(1+\delta)^d-n^{d}\delta^d}{\frac{1+\delta}{n\delta}-1}\right)\\
% &=\binom{d}{\beta}\left(p(e_1)(1-\delta)^d-\left(\frac{1+\delta}{n\delta}-1\right)^{-1}\left((1+\delta)^d-n^{d}\delta^d\right)\right)\\
% &\ge\binom{d}{\beta}\left(p(e_1)(1-\delta)^d-\left(\frac{1}{n\delta}-1\right)^{-1}(1+\delta)^d\right)\\
% % &=\binom{d}{\beta}\left(\frac{1}{n\delta}-1\right)^{-1}\left(\left(\frac{1}{n\delta}-1\right)p(e_1)(1-\delta)^d(1+\delta)^d\right)\\
% % &\ge \binom{d}{\beta}\left(\frac{1+\delta}{n\delta}-1\right)^{-1}\left(\left(\frac{1}{n\delta}-1\right)p(e_1)(1-\delta)^d-(1+\delta)^d\right).
% \end{aligned}\]
where the second inequality follows from the fact that the number of distinct $n$-dimensional multi-indices $\alpha$ satisfying $|\alpha|=d$ and $\alpha_1=l$ is given by $\binom{d-l+n-2}{n-2}\le n^{d-l}$. 

For any $\epsilon>1/p(e_1)>0$, since
$$\delta\le \frac{1}{(1+\eta(p)\epsilon)n}\le \frac{1}{(1+\|p\|_{\infty}\epsilon)n},$$
we have
\[\sum_{\alpha:|\alpha|=d}p_{\alpha}c(\alpha,\beta)\ge 
\binom{d}{\beta}\left(p(e_1)(1-\delta)^d-\epsilon^{-1}(1+\delta)^d\right).
\]
This expression is nonnegative because
$$\delta\le \frac{(p(e_1)\epsilon)^{1/d}-1}{(p(e_1)\epsilon)^{1/d}+1}.$$
Therefore, the coefficient of $x^{\beta}$ in $p(Vx)$ remains nonnegative for every multi-index $\beta$ with $|\beta|=d$ whenever $\delta\le\delta(p;e_1)$. 
% \[\delta\le \min\left\{\frac{( p(e_1)\epsilon)^{1/d}-1}{( p(e_1)\epsilon)^{1/d}+1},\frac{1}{(1+\eta(p)\epsilon)n}\right\}.\]
% Since this bound holds for any $\epsilon>1/p(e_1)$, we can take the supreme over all admissible $\epsilon$ to obtain the bound $\delta(p;e_1)$ stated in the theorem. 
Therefore, the theorem follows for the case where $x^*=e_1$.

For a general point $x^*\in \mbS^{n-1}$, there exists an orthogonal matrix $U\in\reals^{n\times n}$, such that $Ue_1=x^*$. Let  $p_U(x)=p(Ux)$. Then~$p_U$ is a polynomial with $p_U(e_1)=p(x^*)>0$. Applying the statement of the theorem to~$p_U$, we have that for every matrix  $\tilde{V}\in\reals^{n\times n}$ with $\|\tilde{V}-e_1\mathbf{1}^T\|_{F}\le \delta(p_U;e_1)$, the polynomial $p_U(\tilde{V}x)=p(U\tilde{V}x)$ has nonnegative coefficients. Hence, for every matrix $V\in\reals^{n\times n}$ with $\|V-x^*\mathbf{1}^T\|_F\le \delta(p_U;e_1)$, the matrix $U^TV$ satisfies
\[\|U^TV-e_1\mathbf{1}^T\|_{F}=\|U(U^TV-e_1\mathbf{1}^T)\|_F=\|V-x^*\mathbf{1}^T\|_F\le \delta(p_U;e_1).\]
This implies that $p(Vx)=p(U(U^TV)x)$ has nonnegative coefficients. Finally, observe that $\eta(p)$ is unitarily invariant, i.e., $\eta(p)=\eta(p_U)$ for any orthogonal matrix $U$. Consequently,
\[\begin{aligned}
\delta(p_U;e_1)
    &=\sup_{\epsilon>1/p_U(e_1)}\min\left\{\frac{(p_U(e_1)\epsilon)^{1/d}-1}{(p_U(e_1)\epsilon)^{1/d}+1},\frac{1}{(1+\eta(p_U)\epsilon)n}\right\}\\
    &=\sup_{\epsilon>1/p(x^*)}\min\left\{\frac{(p(x^*)\epsilon)^{1/d}-1}{( p(x^*)\epsilon)^{1/d}+1},\frac{1}{(1+\eta(p)\epsilon)n}\right\}\\
    &=\delta(p;x^*),
\end{aligned}\]
which completes the proof.
\end{proof}

% Theorem~\ref{thm:degd_homo_new} provides a local Positivstellensatz, similar to Theorem~\ref{thm:degd_homo}. As a consequence, we explicitly construct a family of simplicial disjunctions in Theorem~\ref{cor:strictly_pos_new} and show that any positive definite form has nonnegative coefficients with respect to some simplicial disjunction in the family.

% As a consequence, if $p\in H_{n,d}$ is positive definite, we can always find a simplicial disjunction $\mcV=\{V_k\}_{k=1,\ldots,r}$ such that $p\in\DiSOSL_{n,d}(\mcV)$. This is because when all the columns of $V_k$ belong to the unit sphere and the spherical polygon formed by these columns, i.e., the set $\cone\{V_k\}\cap \mbS^{n-1}$, has small enough diameter, we have $V_k$ is invertible and $p(V_k x)$ has nonnegative coefficients. Since the diameter is uniformly lower bounded by Theorem~\ref{thm:degd_homo_new}, we can select a finite number of simplicial cones that covers the unit sphere. This guarantees that $p\in \DiSOSL_{n,d}(\mcV)$. Additionally, similar to Theorem~\ref{cor:strictly_pos}, we can explicitly construct a family of simplicial disjunctions, such that any positive definite form has nonnegative coefficients with respect to some simplicial disjunction in the family.
\subsection{Proof of the Optimization-Free Disjunctive Positivstellensatz}

With the theorem from the previous subsection in place, we are now ready to prove Theorem~\ref{cor:strictly_pos_new}.
\begin{proof}[Proof of Theorem~\ref{cor:strictly_pos_new}]
We construct 
$\mcV^{1}$ using $2^n$ matrices, each of the form $\left[\pm e_1\, \ldots\, \pm e_n\right]$, so that the conic hulls of their columns correspond to the $2^n$ orthants of $\reals^n$. We describe how to further refine the partitioning of the space by focusing on the nonnegative orthant (the other orthants can be handled analogously by flipping signs). 
The columns of the matrix $\left[e_1\, \ldots\,  e_n\right]$ form the vertices of the unit simplex. 
For $m\ge 2$, we divide each edge of the unit simplex into $m$ equal segments of length $\sqrt{2}/m$ and partition the simplex into $m^{n-1}$ smaller simplices. Taking the conic hulls of these simplices partitions the nonnegative orthant into $m^{n-1}$ simplicial cones, whose extreme rays associate with the columns of $m^{n-1}$ invertible $n\times n$ matrices. Specifically, if $w_1,\ldots,w_n$ are the vertices of one of the smaller simplices, we include the invertible matrix 
$$V=\left[\frac{w_1}{\|w_1\|_2} \ \cdots\ \frac{w_n}{\|w_n\|_2}\right]$$
as a member of the simplicial disjunction $\mcV^m$. Repeating this procedure for each of the $2^n$ orthants, we obtain a total of $m^{n-1}2^n$ matrices that together define $\mcV^m$.

We now estimate how fast the maximum pairwise distance between the columns of each matrix $V$ constructed as above tends to zero as $m\to\infty$. Note that at level $m$, the vertices of the small simplices have the form $\alpha/m$, where $\alpha=(\alpha_1,\ldots,\alpha_n)$ is an $n$-tuple of nonnegative integers such that $\sum_{i=1}^n\alpha_i=m$. Let $\alpha/m$ and $\beta/m$ be two vertices of the same simplex. Since $\|\alpha/m-\beta/m\|_2=\sqrt{2}/m$, we have $\|\alpha-\beta\|_2=\sqrt{2}$. Note that their corresponding columns in $V$ are $\alpha/\|\alpha\|_2$ and $\beta/\|\beta\|_2$. The distance between these two columns can be upper bounded as
% Each conic hull intersects $\mbS^{n-1}$ and forms a spherical simplex with exactly $n$ vertices. To establish the desired result, it suffices to show that the maximum pairwise distance between the vertices of each spherical simplex tends to zero as $m\to\infty$. For each small simplex, its vertices are in the form of $\alpha/m$, where $\alpha\in\reals^n$ has nonnegative integer entries such that $\sum_{i=1}^n\alpha_i=m$. If $\alpha/m$ and $\beta/m$ are two vertices of the same simplex, then $\|\alpha/m-\beta/m\|_2=\sqrt{2}/m$ and thus $\|\alpha-\beta\|_2=\sqrt{2}$. Observe that the corresponding vertices of the spherical simplex are $\alpha/\|\alpha\|_2$ and $\beta/\|\beta\|_2$. Their distance can be upper bounded as:
\begin{equation}\label{eq:dis_vertices}
   \begin{aligned}
\left\|\frac{\alpha
}{\|\alpha\|_2}-\frac{\beta
}{\|\beta\|_2}\right\|_2&=\left\|\frac{\alpha-\beta}{\|\alpha\|_2}+\beta\left(\frac{1
}{\|\alpha\|_2}-\frac{1
}{\|\beta\|_2}\right)\right\|_2\\
&=\left\|\frac{\alpha-\beta}{\|\alpha\|_2}+\frac{\beta}{\|\beta\|_2}\cdot\frac{\|\beta\|_2-\|\alpha\|_2
}{\|\alpha\|_2}\right\|_2\\
&\le \frac{\|\alpha-\beta\|_2}{\|\alpha\|_2}+\frac{|\|\beta\|_2-\|\alpha\|_2|}{\|\alpha\|_2}\\
&\le \frac{2\|\alpha-\beta\|_2}{\|\alpha\|_2}\le \frac{2\sqrt{2n}}{m},
\end{aligned} 
\end{equation}
where we have used the triangle inequality and the fact that 
\[\|\alpha\|_2^2\ge \frac{1}{n}\left(\sum_{i=1}^n\alpha_i\right)^2=\frac{m^2}{n}.\] 
% Let the $n$ vertices of a spherical simplex be the columns of $V\in\reals^{n\times n}$, and 
Let $v\in\mbS^{n-1}$ denote any column of a matrix $V$ constructed as above. Then, it follows from~\eqref{eq:dis_vertices} that $\|V-v\mathbf{1}^T\|_F\le 2\sqrt{2}n/m$. 
Let $p_{\min}=\min_{x\in\mbS^{n-1}}p(x)$, and let $\eta(p)$ and $\delta(p;v)$ be as in Theorem~\ref{thm:degd_homo_new}. Define
\[\tilde{M}(p)\defeq \left\lceil2\sqrt{2}n\left(\sup_{\epsilon>1/p_{\min}}\min\left\{\frac{(p_{\min}\epsilon)^{1/d}-1}{( p_{\min}\epsilon)^{1/d}+1},\frac{1}{(1+\eta(p)\epsilon)n}\right\}\right)^{-1}\right\rceil.\]
% Let $\delta(p;v)$ be as in Theorem~\ref{thm:degd_homo_new}. 
It follows that $\|V-v\mathbf{1}^T\|_F\le 2\sqrt{2}n/\tilde{M}(p)
\le \delta(p;v)$ whenever $m\ge \tilde{M}(p)$. By Theorem~\ref{thm:degd_homo_new}, the polynomial $p(Vx)$ has nonnegative coefficients. Since the same argument applies to every matrix $V$ constructed as above, we conclude that $p\in \DiSOSL_{n,d}(\mcV^{m})$ whenever $m\ge \tilde{M}(p)$. 
\end{proof}

% \begin{remark}
% As a byproduct of the proof of Theorem~\ref{cor:strictly_pos_new}, certifying nonnegativity of a polynomial~$p$ without solving optimization problems may require up to a super-exponential number of subregions---namely, $\tilde{M}(p)^{n-1}2^n$. 
% % In comparison with Theorem~\ref{thm:num_ad}, it reveals a trade-off between the number of hierarchy levels and the computational complexity per level. 
% \end{remark}

Analogous to Algorithm~\ref{alg:uniform_grid}, which is based on Theorem \ref{cor:strictly_pos}, the optimization-free disjunctive Positivstellensatz can be used to construct 
an optimization-free algorithm for certifying polynomial nonnegativity. In practice, instead of uniformly refining each simplicial cone as in the proof, it is more efficient to adaptively refine these cones within a spatial branch-and-bound framework. Even with such a modification, we expect the optimization-free approach to require a significantly larger number of regions to obtain a disjunctive certificate of nonnegativity compared to the disjunctive sum of squares approach based on semidefinite programming (i.e., the approach of Section~\ref{sec:disos_psatz}). Studying these tradeoffs computationally is a direction for future research.

% Theorem~\ref{cor:strictly_pos_new} can also be used to construct an algorithm for certifying polynomial nonnegativity without solving optimization problems, analogous to Algorithm~\ref{alg:uniform_grid} for the SDP-based Positivstellensatz of Theorem~\ref{cor:strictly_pos}.
% In practice, instead of uniformly refining each simplicial cone as in the proof, it is more efficient to adaptively refine these cones within a spatial branch-and-bound framework; see Section~\ref{sec:bnb} for a related approach.
% However, we expect the optimization-free approach to require a significantly larger number of regions compared to the semidefinite-programming-based approach of Section~\ref{sec:disos_psatz}.

%We note that the results of this section are mainly of theoretical interest, as the number of polynomials whose coefficients should be checked can be prohibitively large in practice.

%See Section~\ref{sec:bnb} for related computational experiments.


%To further improve practical efficiency, we also adopt an adaptive splitting strategy within a branch-and-bound framework, as discussed in Section~\ref{sec:bnb}.


% To improve practical efficiency, instead of relying on the fixed splitting strategy implied by Theorem~\ref{cor:strictly_pos_new}, one may adopt an adaptive splitting strategy. This approach uses a branch-and-bound framework and is described in Section~\ref{sec:bnb}.

% Given a polynomial $p\in H_{n,d}$ and a simplicial disjunction $\mcV=\{V_k\}_{k=1,\ldots,r}$, membership in $\DiSOSL_{n,d}(\mcV)$ can be verified by substituting $x$ with $V_kx$ in~$p$ and checking whether all coefficients of the resulting polynomial are nonnegative. This procedure yields an optimization-free hierarchy for certifying nonnegativity of~$p$: one sequentially tests whether $p\in\DiSOSL_{n,d}(\mcV^r)$ for the simplicial disjunctions ${\mcV^r}$ in the family. 
% By Theorem~\ref{cor:strictly_pos_new}, for any pd form $p$, there is always a simplicial disjunction in the family, such that $p$ has nonnegative coefficients with respect to it. 




\section{Extensions to Constrained Optimization Problems}
\label{sec:constrained_opt}
In this section, we extend the disjunctive sum of squares method from certifying global nonnegativity of polynomials to certifying nonnegativity over a closed basic semialgebraic set, i.e., a subset of the Euclidean space defined by finitely many polynomial inequalities. 
We present two approaches. The first, discussed in Section~\ref{sec:hierarchy}, builds on a polynomial-time reduction from~\cite{aaa2019} that converts constrained polynomial optimization problems (POPs) into unconstrained ones, to which our disjunctive sum of squares method then applies. The second, given in Section~\ref{sec:copositive}, extends Definition~\ref{def:disjunction} so that the subregions defined by the algebraic disjunction cover the feasible set of interest rather than all of $\reals^n$. We illustrate this approach for the problem of testing matrix copositivity, which reduces to minimizing a quadratic form over the nonnegative orthant.

% This extension builds on a polynomial-time reduction proposed in~\cite{aaa2019} from constrained polynomial optimization problems to unconstrained ones. 
% Although this reduction provides a general recipe for POPs, it can be inefficient in practice. In some important cases, one can instead partition the feasible region directly using algebraic disjunctions, in the same spirit as the partitions of $\reals^n$ used in the unconstrained setting. As an illustration of this idea, Section~\ref{sec:copositive} applies the approach to certifying matrix copositivity, which reduces to minimizing a quadratic form over the nonnegative orthant.
% \section{Converging hierarchies based on disos}\label{sec:hierarchy}
\subsection{Lower Bounds on the Optimal Values of Polynomial Optimization Problems}\label{sec:hierarchy}
Consider a POP of the form
\begin{equation}\label{eq:general_pop}
	\begin{array}{rl}
    \min\limits_{x\in{\scriptsize\reals^n}} & p(x)\\
		\st & g_j(x)\ge 0, \quad j=1,\ldots,J,
	\end{array}    
\end{equation}
with a compact feasible set. Throughout this section, we let $p^*$ denote the optimal value of problem~\eqref{eq:general_pop}.  
We show that our disjunctive Positivstellens\"atze can be combined with the reduction from~\cite[Theorem 2.1]{aaa2019} to produce a converging hierarchy of lower bounds on $p^*$. 
% a converging hierarchy of lower bounds on $p^*$ following the reduction given in~\cite[Theorem 2.4]{aaa2019}. 

We first describe the reduction from~\cite{aaa2019}. 
Let $d$ be the smallest even integer larger or equal to the maximum degrees of the polynomials $p$ and $g_1,\ldots,g_J$. Following~\cite[Theorem 2.1]{aaa2019}, one can explicitly construct a form $f_{\gamma}(z)$ in $N\defeq n+J+3$ variables $z$ and of degree $2d$ using a given scalar $\gamma$, the coefficients of $p, g_1,\ldots,g_J$, and an upper bound on the radius of a ball containing the feasible set of problem~\eqref{eq:general_pop}. This construction has the property that $\gamma$ is a strict lower bound on $p^*$ if and only if $f_{\gamma}(z)$ is positive definite. 

Then,~\cite[Theorem 2.4]{aaa2019} shows how to use this reduction to solve problem~\eqref{eq:general_pop}. 
Recall that $H_{n,d}$ denotes the set of homogeneous polynomials in~$n$ variables and degree~$d$, and $\HNN_{n,d}$ denotes the set of nonnegative forms in $H_{n,d}$. 
Consider a sequence of subsets of $H_{n,d}$, denoted by $K_{n,d}^i$ and indexed by $i$, that satisfies the following properties:
% Consider a sequence of sets of polynomials in $H_{n,d}$, denoted by $K_{n,d}^m$, that satisfies the following properties:
\begin{enumerate}[label=(\arabic*)]
    \item $K_{n,d}^i\subseteq \HNN_{n,d}$ for all $i$, and there exists a positive definite form $s_{n,d}\in K^0_{n,d}$,
    \item if $p\in H_{n,d}$ is positive definite, then there exists an integer $i$ such that $p\in K_{n,d}^i$,
    \item $K_{n,d}^i\subseteq K_{n,d}^{i+1}$ for all $i$,
    \item if $p\in K_{n,d}^i$, then $p + \epsilon s_{n,d}\in K_{n,d}^i$ for all $\epsilon\in[0,1]$.
\end{enumerate}

By~\cite[Theorem 2.4]{aaa2019}, 
for any such sequence $K_{n,d}^i$ satisfying properties~(1)--(4), the sequence $l_i$ defined by
\begin{equation}\label{eq:general_hierarchy}
     	\begin{array}{rl}
		l_i \defeq \max\limits_{\gamma\in\reals} & \gamma\\
		\st & f_{\gamma}(z)-\dfrac{1}{i}s_{N,2d}(z)\in K_{N,2d}^i,
	\end{array}  
\end{equation}
is non-decreasing, satisfies $l_i\le p^*$ for all~$i$, and $\lim_{i\to\infty} l_i=p^*$. 
As shown in~\cite{aaa2019}, 
Artin, Reznick, and Poly\'a's Positivstellens\"atze all lead to sequences $K_{n,d}^i$ satisfying properties~(1)--(4).
In Theorem~\ref{thm:hier_pop} below, we show that our disjunctive sum of squares method fits into this framework as well. 
% with a key advantage of keeping the polynomial degrees in~\eqref{eq:general_hierarchy} fixed (and low) throughout the hierarchy.
Moreover, in contrast to the above Positivstellens\"atze, our method has the 
advantage that the degrees of the polynomials required to certify membership in $K_{n,d}^i$ remain fixed (and low) throughout the hierarchy.

We recall that $\HSOS_{n,d}$ denotes the set of sos forms in~$n$ variables and degree~$d$.


% degrees of the polynomials in~\eqref{eq:general_hierarchy} fixed (and low) throughout the hierarchy.


% is that our method keeps the polynomial degrees in~\eqref{eq:general_hierarchy} fixed (and low) throughout the hierarchy.

% the increasing degrees required by the Positivstellens\"atze of Artin, Reznick, and Poly\'a.
% This general framework encompasses multiple Positivstellens\"atze, including those by Artin, Reznick, and Poly\'a. The proof of~\cite[Theorem 2.4]{aaa2019} relies on the fact that the sets $K_{n,d}^m$, constructed from these Positivstellens\"atze, can approximate the cone of positive definite forms arbitrarily well. Notably, the hierarchy also extends to our disjunctive Positivstellensatz presented in Theorem~\ref{cor:strictly_pos}, 
% as established in Theorem~\ref{thm:hier_pop} below.
% since the sequence of cones $\DiSOS_{n,d}^d(\mcD^r)$, with the family of algebraic disjunctions $\mcD^r$ defined therein, satisfies the four properties listed above.
\begin{theorem}\label{thm:hier_pop}
Fix a dimension $n$ and an even degree $d$. Let $\mcD^m$ be any family of algebraic disjunctions, indexed by $m$, that satisfies the guarantees of Theorem~\ref{cor:strictly_pos}. 
Let $\widehat{K}_{n,d}^0\defeq\HSOS_{n,d}$. For each integer $i\ge 1$, define
    \begin{equation}\label{eq:hier_pop}
      \widehat{K}_{n,d}^i\defeq  \bigcup_{m=1}^i\DiSOS_{n,d}^d\left(\mcD^m\right) .
    \end{equation}
Then, the sequence $\widehat{K}_{n,d}^i\cap H_{n,d}$ satisfies properties~(1)--(4) stated above.
\end{theorem}

\begin{proof}
We verify properties~(1)--(4). 
% Recall that $\HSOS_{n,d}$ is the set of homogeneous sos polynomials in~$n$ variables and degree~$d$.
\begin{enumerate}[label=(\arabic*)]
        \item When $i=0$, $\widehat{K}_{n,d}^0=\HSOS_{n,d}\
        \subseteq \HNN_{n,d}$. 
        % the property holds trivially
        Let $i\ge 1$.         
        Since $\SOS_{n,d}\subseteq\DiSOS^d_{n,d}(\mcD^m)\subseteq\NN_{n,d}$ for all $m$, we have
        $\HSOS_{n,d}\subseteq\DiSOS^d_{n,d}(\mcD^m)\cap H_{n,d}\subseteq\HNN_{n,d}$ for all $m$. Taking unions gives 
        $\HSOS_{n,d}\subseteq \widehat{K}_{n,d}^i\cap H_{n,d}\subseteq \HNN_{n,d}$ for all $i\ge 1$. Moreover, the positive definite form~$\|x\|_2^d\in\HSOS_{n,d}=\widehat{K}_{n,d}^0$.
        \item By Theorem~\ref{cor:strictly_pos}, for any positive definite form~$p$, there exists an integer $\bar{m}$ such that $p\in \DiSOS^d_{n,d}(\mcD^{\bar{m}})$, and therefore $p\in\widehat{K}_{n,d}^{\bar{m}}$.
        \item The sequence $\widehat{K}_{n,d}^i$ is nested by construction.        
        % By definition of $\widehat{K}_{n,d}^i$ in~\eqref{eq:hier_pop}, the sequence is nested.
        \item When $i=0$, the property holds trivially. 
        When $i\ge 1$, suppose $p\in\widehat{K}_{n,d}^i$. Then, there exists $m\in\{1,\ldots,i\}$ such that $p\in\SOS_{n,d}^d(\mcD^m)$. 
        Recall from Theorem~\ref{cor:strictly_pos} that each subset in $\mcD^m$ contains exactly one form, i.e., we can write
        \[\mcD^m=\big\{\{q_{j}(x)\},\quad j=1,\ldots,r_m\big\}.\]
 For each $j\in\{1,\ldots,r_m\}$, 
        there exist sos polynomials $s_{j}$ and $\sigma_{j}$ such that
        \[p(x)=s_{j}(x)+\sigma_{j}(x)q_{j}(x),\]
        with $\deg s_{j}\le d$ and $\deg (\sigma_{j}q_{j})\le d$. 
        Consequently, 
        \[p(x) + \epsilon \|x\|_2^d=(s_{j}(x)+\epsilon \|x\|_2^d)+ \sigma_{j}(x)q_{j}(x).\]
        Since $\|x\|_2^d$ is sos, the polynomial $s_{j}(x)+\epsilon \|x\|_2^d$ remains sos for all $\epsilon\in[0,1]$ and $j=1,\ldots,r_m$. Therefore, 
        $p(x) + \epsilon \|x\|_2^d\in \SOS_{n,d}^d(\mcD^m)\subseteq\widehat{K}_{n,d}^i$  for all $\epsilon\in[0,1]$.
    \end{enumerate}
\vspace{-\baselineskip}
% Therefore, by~\cite[Theorem 2.4]{aaa2019}, the statement of the theorem follows.
\end{proof}
Since the sequence $\widehat{K}_{n,d}^i$ satisfies properties~(1)--(4), replacing\footnote{
To compute the optimal value $l_i$ of problem~\eqref{eq:general_hierarchy} with $K_{N,2d}^i$ replaced by $\widehat{K}_{N,2d}^i$, one can take the maximum between $l_{i-1}$ and the optimal value of problem~\eqref{eq:general_hierarchy} with $K_{N,2d}^i$ replaced by $\DiSOS^{2d}_{N,2d}(\mcD^i)$. As discussed earlier, a membership constraint to $\DiSOS^{2d}_{N,2d}(\mcD^i)$ amounts to solving a semidefinite program.
} $K_{N,2d}^i$ with $\widehat{K}_{N,2d}^i$ in~\eqref{eq:general_hierarchy} and applying~\cite[Theorem 2.4]{aaa2019} yields a non-decreasing sequence $l_i$ that satisfies 
$l_i\le p^*$ for all $i$ and $\lim_{i\to\infty} l_i=p^*$. 
Furthermore, one can show that 
 the theorem remains valid if the cones $\DiSOS_{n,d}^d(\mcD^m)$ in~\eqref{eq:hier_pop} are replaced by the cones $\DiSOSL_{n,d}(\mcV^m)$ introduced in Theorem~\ref{cor:strictly_pos_new}, yielding an optimization-free hierarchy for general POPs with a compact feasible set.
 % analogous to the hierarchy based on Poly\'a's Positivstellensatz~\cite[Theorem 4.1]{aaa2019}, which similarly involves checking nonnegativity of polynomial coefficients.
% \end{remark}
% \begin{remark}
% The theorem remains valid if the cones $\DiSOS_{n,d}^d(\mcD^m)$ in~\eqref{eq:hier_pop} are replaced by the cones $\DiSOSL_{n,d}(\mcV^m)$ introduced in Theorem~\ref{cor:strictly_pos_new}. In particular, this yields an optimization-free hierarchy for general POPs, analogous to the hierarchy based on Poly\'a's Positivstellensatz~\cite[Theorem 4.1]{aaa2019}, which similarly involves checking nonnegativity of polynomial coefficients.
% \end{remark}

\subsection{An Application to Copositive Programming}\label{sec:copositive}
In this section, we give an example of how the disjunctive sum of squares approach can be adapted to certify nonnegativity over a subset of the Euclidean space. We take this subset to be the nonnegative orthant, which arises naturally in copositive programming.
% which is relevant to a broad range of applications, e.g., in copositive programming. 

A symmetric matrix $Q \in \reals^{n \times n}$ is said to be \emph{copositive} if $x^T Q x \geq 0$ for all $x \geq 0$. Let $\mcC_n$ denote the cone of $n\times n$ copositive matrices. A copositive program is a linear optimization problem over the intersection of $\mcC_n$ with an affine subspace. 
% Copositive programming is remarkably expressive: 
In recent years, copositive programming has found many applications in both discrete and continuous optimization; see, e.g.,~\cite{burer2012copositive,dur2010copositive}. 
% many NP-hard problems, including quadratic optimization over the unit simplex, maximum clique, graph coloring, and stable set problems, can be formulated exactly as copositive programs~\cite{burer2012copositive,dur2010copositive}.
Since checking copositivity of a given matrix is NP-hard~\cite{murty1987some}, tractable inner approximations to $\mcC_n$ are of interest in most applications of copositive programming. 
% In the majority of the applications,
% A fundamental obstacle, however, is that deciding whether a given matrix is copositive is NP-hard~\cite{murty1987some}.
% This motivates our study of tractable inner approximations of the copositive cone $\mcC_n$.

A well-known SDP-based sufficient condition for copositivity is the so-called {``$P+N$'' criterion}. For a symmetric matrix $A \in \reals^{n \times n}$, let us use the standard notation $A \succeq 0$ (resp. $A \geq 0$) to denote that $A$ is positive semidefinite (resp. elementwise nonnegative). Observe that if a matrix $Q\in\reals^{n\times n}$ can be written as $Q = P + N$, where $P \succeq 0$ and $N \geq 0$, then $Q$ is copositive. 
This condition is known to also be necessary for copositivity if $n\le 4$~\cite{Diananda1962}. For $n\ge 5$, however, there are copositive matrices that do not admit a $P+N$ decomposition. A classical example is the Horn matrix~\cite{Diananda1962}:
\begin{equation}\label{eq:horn}
  H=\begin{bmatrix}
    \phantom{-}1 & -1 & \phantom{-}1 & \phantom{-}1 & -1 \\
    -1 & \phantom{-}1 & -1 & \phantom{-}1 & \phantom{-}1 \\
    \phantom{-}1 & -1 & \phantom{-}1 & -1 & \phantom{-}1 \\
    \phantom{-}1 & \phantom{-}1 & -1 & \phantom{-}1 & -1 \\
    -1 & \phantom{-}1 & \phantom{-}1 & -1 & \phantom{-}1 \\
  \end{bmatrix}.
\end{equation}

% which is copositive but admits no $P+N$ decomposition.

To bridge the gap between the $P+N$ criterion and copositivity, a complete SDP-based hierarchy of sufficient conditions
for copositivity was introduced in~\cite{Parrilo2000}.
Let $x.^2\defeq(x_1^2,\ldots,x_n^2)$. At level $m$ of this hierarchy, one checks if the polynomial $(x.^2)^TQ(x.^2)\cdot(x_1^2 + \cdots + x_n^2)^m$ is sos, and if so, copositivity is certified. The zeroth level of this hierarchy recovers exactly the $P+N$ criterion described above~\cite{Parrilo2000}.
Moreover, this hierarchy is complete: for any strictly copositive matrix $Q$ (i.e., a matrix $Q$ satisfying $x^TQx> 0$ for all $x\ge 0$ with $x\neq 0$), there exists a level $m$ at which copositivity is certified~\cite{Parrilo2000, reznick1995uniform}.
However, the size of the semidefinite constraint at level $m$ of this hierarchy grows as $\binom{n+m+2}{n} \times \binom{n+m+2}{n}$. While one can reduce this size using symmetry reduction techniques~\cite[Section 8.1]{Gatermann2004}, the higher levels of this hierarchy often remain too expensive to solve. 
% , making higher levels of the hierarchy increasingly expensive to solve, even after applying symmetry reduction techniques~\cite{Gatermann2004}.
%Here, we utilize the disjunctive sos approach to propose an alternative SDP-based complete hierarchy for certifying copositivity of an $n\times n$ matrix where the size of the semidefinite constraint remains $n\times n$ throughout the hierarchy. 

Here, we utilize the disjunctive sos approach to propose an alternative SDP-based complete hierarchy for certifying copositivity of an $n\times n$ matrix, where the size of the semidefinite constraint remains $n\times n$ throughout the hierarchy. To this end, we extend the notion of a simplicial disjunction in Definition~\ref{def:disos_new} and require the subregions to cover only the nonnegative orthant rather than all of $\reals^n$.

\begin{definition}\label{def:sd_simplex}
Let $r$ be a positive integer. We say that a collection of $n\times n$ invertible matrices $\mcV=\{V_k\}_{k=1,\ldots,r}$ forms a \emph{simplicial disjunction} of a region $\Omega\subseteq\reals^n$ if 
\[\Omega\subseteq \bigcup\limits_{k=1}^r\Omega_k,\quad \text{where }\quad\Omega_k=\cone\{V_k\}\defeq\cone\{v_{k,1}, \ldots, v_{k,n}\},\]
and $v_{k,j}$ is the $j$-th column of $V_k$.
\end{definition}

\begin{definition}\label{def:copos_disos}
A symmetric matrix $Q\in\reals^{n\times n}$ is \emph{disjunctive $P+N$ with respect to a simplicial disjunction} $\mcV=\{V_k\}_{k=1,\ldots,r}$ of the nonnegative orthant if there exist matrices $P_k$ and $N_k$, for $k=1,\ldots,r$, such that
\[V_k^TQV_k=P_k+N_k,\;P_k\succeq 0,\; N_k\ge 0,\quad k=1,\ldots,r.\]
We denote the set of such matrices by $\mcPN_{n}(\mcV)$. 
\end{definition}

% \begin{remark}
% Setting $P_k = 0$ in Definition~\ref{def:copos_disos} reduces the condition 
% $V_k^T Q V_k = P_k + N_k$ to $V_k^T Q V_k \geq 0$ elementwise, 
% which coincides with the copositivity detection approach of~\cite{BUNDFUSS2008}. 
% In this case, linear optimization over $\mcC_n(\mcV)$ reduces to a linear program, 
% making it cheaper to solve but potentially requiring a finer partition of the unit simplex to certify 
% copositivity compared to the our approach; see Section~\ref{sec:num} for numerical experiments.
% \end{remark}

Observe that for any simplicial disjunction $\mcV$ of the nonnegative orthant, the set $\mcPN_n(\mcV)$ is a semidefinite representable convex cone. Moreover, it is straightforward to see that $\mcPN_n(\mcV)\subseteq\mcC_n$. Indeed, 
if $Q\in\mcPN_n(\mcV)$,  
the polynomial $y^TV_k^TQV_ky$ is nonnegative for all $y\ge 0$ and $k=1,\ldots,r$. This implies that $x^TQx$ is nonnegative on each subregion $\Omega_k=\cone\{V_k\}$ as any point $x\in\cone\{V_k\}$ can be written as $x=V_ky$ with $y\ge 0$. Since these subregions cover the nonnegative orthant, we have $x^TQx\ge 0$ 
for all $x\ge 0$, and hence $Q$ is copositive.

Let us revisit the Horn matrix $H$ in~\eqref{eq:horn} in view of Definition~\ref{def:copos_disos}. 
Consider the simplicial disjunction $\mcV=\{V_1,V_2\}$ of the nonnegative orthant, where
\[V_1=\begin{bmatrix}
     1   & 0&    0&     0&   0\\
    0 &    1&     0&    0&     0\\
   0&   0&     1&   0&     0\\
    0 &   0&    0&     1&    0.5\\
    0 &    0&     0&    0&     0.5\\
\end{bmatrix}, \quad 
V_2=\begin{bmatrix}     
     1   & 0&    0&     0&   0\\
    0 &    1&     0&    0&     0\\
   0&   0&     1&   0&     0\\
    0 &   0&    0&     0.5&    0\\
    0 &    0&     0&    0.5&     1\\ \end{bmatrix}.\]
One can readily verify that
\[
V_1^THV_1=\begin{bmatrix}
    \phantom{-}1 & -1            & \phantom{-}1  & \phantom{-}1  & \phantom{-}0 \\
    -1            & \phantom{-}1  & -1            & \phantom{-}1  & \phantom{-}1 \\
    \phantom{-}1 & -1            & \phantom{-}1  & -1            & \phantom{-}0 \\
    \phantom{-}1 & \phantom{-}1  & -1            & \phantom{-}1  & \phantom{-}0 \\
    \phantom{-}0 & \phantom{-}1  & \phantom{-}0  & \phantom{-}0  & \phantom{-}0 \\
\end{bmatrix}=\begin{bmatrix}
    \phantom{-}1 & -1            & \phantom{-}1  & -1            & \phantom{-}0 \\
    -1            & \phantom{-}1  & -1            & \phantom{-}1  & \phantom{-}0 \\
    \phantom{-}1 & -1            & \phantom{-}1  & -1            & \phantom{-}0 \\
    -1            & \phantom{-}1  & -1            & \phantom{-}1  & \phantom{-}0 \\
    \phantom{-}0 & \phantom{-}0  & \phantom{-}0  & \phantom{-}0  & \phantom{-}0 \\
\end{bmatrix}+\begin{bmatrix}
     0 & 0 & 0 & 2 & 0 \\
     0 & 0 & 0 & 0 & 1 \\
     0 & 0 & 0 & 0 & 0 \\
     2 & 0 & 0 & 0 & 0 \\
     0 & 1 & 0 & 0 & 0 \\
\end{bmatrix},
\]
\[
V_2^THV_2=\begin{bmatrix}
    \phantom{-}1 & -1            & \phantom{-}1  & \phantom{-}0  & -1            \\
    -1            & \phantom{-}1  & -1            & \phantom{-}1  & \phantom{-}1 \\
    \phantom{-}1 & -1            & \phantom{-}1  & \phantom{-}0  & \phantom{-}1 \\
    \phantom{-}0 & \phantom{-}1  & \phantom{-}0  & \phantom{-}0  & \phantom{-}0 \\
    -1            & \phantom{-}1  & \phantom{-}1  & \phantom{-}0  & \phantom{-}1 \\
\end{bmatrix}=\begin{bmatrix}
    \phantom{-}1 & -1            & \phantom{-}1  & \phantom{-}0  & -1            \\
    -1            & \phantom{-}1  & -1            & \phantom{-}0  & \phantom{-}1 \\
    \phantom{-}1 & -1            & \phantom{-}1  & \phantom{-}0  & -1            \\
    \phantom{-}0 & \phantom{-}0  & \phantom{-}0  & \phantom{-}0  & \phantom{-}0 \\
    -1            & \phantom{-}1  & -1            & \phantom{-}0  & \phantom{-}1 \\
\end{bmatrix}+\begin{bmatrix}
     0 & 0 & 0 & 0 & 0 \\
     0 & 0 & 0 & 1 & 0 \\
     0 & 0 & 0 & 0 & 2 \\
     0 & 1 & 0 & 0 & 0 \\
     0 & 0 & 2 & 0 & 0 \\
\end{bmatrix}.
\]
Let $e_1,\ldots,e_5$ be the standard basis vectors in $\reals^5$. The first identity proves nonnegativity of $x^THx$ over $\cone\{V_1\}=\cone\{e_1,e_2,e_3,e_4,(e_4+e_5)/2\}$. Similarly, the second identity proves nonnegativity of $x^THx$ over $\cone\{V_2\}=\cone\{e_1,e_2,e_3,(e_4+e_5)/2,e_5\}$. 
Together, these two identities certify that the Horn matrix $H\in\mcPN_5(\mcV)$ and is hence copositive. 

We show in Theorem~\ref{thm:psatz_copositive_full} below that one can a priori fix a family of simplicial disjunctions of the nonnegative orthant such that any strictly copositive matrix is disjunctive $P+N$ with respect to some simplicial disjunction in the family. 
% a proof of copositivity can be given via membership in $\mcPN_n(\mcV^m)$ for some $m$. 
% proofs of copositivity via membership in $\mcPN_n(\mcV)$ are always possible for any strictly copositive matrix, using an explicit, fixed family of simplicial disjunctions $\mcV$ of the nonnegative orthant. 
% one can a priori fix a family of low-degree algebraic disjunctions with respect to which there always exists a low-degree disjunctive sos proof of nonnegativity. 
In fact, the proof shows that the disjunctive proofs do not necessarily need to involve positive semidefinite matrices; that is, one can set $P_k = 0$ in Definition~\ref{def:copos_disos}. This result follows as a simple corollary of what we have proven in the previous section.

%This result follows as a simple corollary of the more general results in Theorem~\ref{cor:strictly_pos_new} and Theorem~\ref{thm:degd_homo_new}.

In the proof of the following theorem, we use the standard notation~$\textcolor{black}{\Delta^{n-1}}$ to denote the unit simplex in~$\reals^n$, i.e., $\textcolor{black}{\Delta^{n-1}}=\{x\in\reals^n\mid \sum_{i=1}^n x_i=1,\; x\ge 0\}$, and we let $\conv\{V\}$ denote the convex hull of the columns of a matrix
$V\in\reals^{n\times n}$.

% We now consider the problem of testing copositivity of a symmetric matrix~$Q\in\reals^{n\times n}$, which is equivalent to certifying that the quadratic form~$p(x) = x^TQx$ is nonnegative for all $x\ge 0$. This is a special case of a constrained POP, where the feasible set is the nonnegative orthant.
% {\bf TO CONNECT to next sentence}
% Equivalently, this amounts to certifying that the quadratic form $p(x)=x^\tpose Qx$ is nonnegative over the nonnegative orthant.

% We denote the set of $n\times n$ copositive matrices as $\mcC_n$. 

% It is straightforward to observe some sufficient conditions to certify copositivity. Firstly, a positive semidefinite matrix or an elementwise nonnegative matrix is copositive. Secondly, copositive matrices are additive. Thus if a matrix $Q$ can be decomposed as $Q=P+N$, where $P\succeq 0$ (positive semidefinite) and $N\ge 0$ (elementwise nonnegative), then $Q\in \mcC_n$. Denote $\mcS_n^+$ as the set of positive semidefinite matrices and $\mcN_n$ as the set of elementwise nonnegative matrices, we have 
% $\mcS_n^++\mcN_n\subseteq \mcC_n$.

% To see the difference between $\mcS_n^++\mcN_n$ and $\mcC_n$, we can view copositivity from the polynomial perspective. On the one hand, it follows from the definition that $Q\in\mcC_n$ is equivalent to the degree-4 polynomial $(x.^2)^TQ(x.^2)\in P_{n,4}$. On the other hand, Parrilo proved his thesis~\cite{Parrilo2000} that $Q\in \mcS_n^++\mcN_n$ if and only if the degree-4 polynomial $(x.^2)^TQ(x.^2)\in\HSOS_{n,4}$, where $(x.^2)=(x_1^2,\ldots,x_n^2)^T$. Thus if $(x.^2)^TQ(x.^2)\in P_{n,4}\backslash\HSOS_{n,4}$, we cannot cerify copositivity using the $P+N$ condition. In fact, the Horn matrix, previously appeared in~\cite{Diananda1962}, is known to be copositive but cannot be decomposed as $P+N$.
% \[H=\begin{bmatrix}
%     1 &-1& 1& 1& -1\\
%      -1& 1& -1& 1& 1\\
%      1& -1& 1& -1& 1\\
%      1& 1& -1& 1& -1\\
%      -1& 1& 1& -1& 1\\
% \end{bmatrix}\]
% In the unconstrained setting, an algebraic disjunction (Definition~\ref{def:disjunction}) requires the subregions $\Omega_1,\ldots,\Omega_r$ to cover all of $\reals^n$. Here, we extend this definition by requiring the subregions to cover only the feasible set, namely the nonnegative orthant. 
% As we will see below, simplicial disjunctions of the nonnegative orthant lead to particularly simple formulations. 
% Given the quadratic nature of~$p$, simplicial disjunctions of the nonnegative orthant provide a natural family of such subregions. 
% We certify global nonnegativity of polynomials via algebraic disjunctions of $\reals^n$. In the constrained setting, we seek disjunctive proofs of nonnegativity of the quadratic form using algebraic disjunctions of the nonnegative orthant. 
% Given the quadratic nature of the objective, simplicial disjunctions of the nonnegative orthant are a natural choice. 


% Let $\mcV=\{V_k\}_{k=1,\ldots,r}$ be a simplicial disjunction of the nonnegative orthant. By Definition~\ref{def:disos_new}, the quadratic form~$p$ has nonnegative coefficients with respect to $\mcV$ if $p(V_kx)=x^TV_k^TQV_kx$ has nonnegative coefficients for all $k=1,\ldots,r$. Equivalently, 
% \[V_k^TQV_k\ge 0,\quad k=1,\ldots,r.\]
% We therefore define the set of disjunctive copositive matrices below.


% Definition~\ref{def:copos_disos} coincides with the copositivity detection approach in~\cite{BUNDFUSS2008}, where~\cite[Definition 1]{BUNDFUSS2008} introduces a simplicial partition of the unit simplex $\textcolor{black}{\Delta^{n-1}}\defeq\{x\in\reals^n\mid\mathbf{1}^Tx=1,\ x\ge 0\}$.
% Definition~\ref{def:copos_disos} coincides with the copositivity detection approach in~\cite{BUNDFUSS2008}. In~\cite[Definition 1]{BUNDFUSS2008}, the authors introduce a simplicial partition of the unit simplex $\textcolor{black}{\Delta^{n-1}}\defeq\{x\in\reals^n\mid1^Tx=1,x\ge 0\}$, which corresponds directly to the simplicial disjunction in our framework.
% The constructions are equivalent: the conical hull of each simplex in a simplicial partition yields a simplicial disjunction of the nonnegative orthant, while intersecting a simplicial disjunction of the nonnegative orthant with the unit simplex yields a simplicial partition.

% In~\cite[Theorem 2]{BUNDFUSS2008}, for any strictly copositive matrix $Q$, there exists $\delta>0$ such that if a simplicial partition of the unit simplex has maximum diameter less than $\delta$, then $Q$ is disjunctive copositive with respect to the corresponding simplicial disjunction.
% Theorem~\ref{thm:copositive_new} below provides an alternative explicit estimate of $\delta$ and leads to a disjunctive Positivstellensatz for matrix copositivity in Theorem~\ref{thm:psatz_copositive_full}. It also shows how the broader disjunctive sum of squares framework developed in this paper applies naturally to this setting.
% In~\cite[Theorem 2]{BUNDFUSS2008}, the authors establish that for any strictly copositive matrix $Q$ (i.e., $x^TQx>0$ for all $x\in\textcolor{black}{\Delta^{n-1}}$), there exists $\delta>0$ such that if a simplicial partition of the unit simplex has a maximum diameter less than $\delta$, then $Q$ is disjunctive copositive with respect to the corresponding simplicial disjunction. We provide an alternative estimate of $\delta$ in Theorem~\ref{thm:copositive_new}. By presenting this alternative perspective, we supplement the theoretical results in~\cite{BUNDFUSS2008}. Moreover, our derivation arises from the broader framework developed in this paper, which shows how the disjunctive sum of squares method applies to this setting.

% aligning with their findings while embedding them into a larger context of the disjunctive sum of squares method, rather than drawing direct comparisons. ; the two formulations are equivalent. However, our derivation arises from the broader framework developed in this paper, which shows how the disjunctive sum of squares method applies to this setting.

% The proof of Theorem~\ref{thm:copositive_new} is similar to that of Theorem~\ref{thm:degd_homo_new}, with two main differences: 1) the points belong to the unit simplex rather than the unit sphere; and 2) the quadratic nature of the objective simplifies the proof, as we can explicitly expand the quadratic form.

\begin{theorem}[Disjunctive Positivstellensatz for Matrix Copositivity]
\label{thm:psatz_copositive_full}
For any fixed dimension $n$, there exists an explicit family of simplicial disjunctions
$\mcV^m$ of the nonnegative orthant, indexed by $m$, such that for any strictly 
copositive matrix $Q \in \reals^{n \times n}$, we have $Q \in \mcPN_{n}(\mcV^m)$ 
for some $m$.
\end{theorem}
\begin{proof}
The construction of the family of simplicial disjunctions $\mcV^m$ is similar to the construction for Theorem~\ref{cor:strictly_pos_new}. We let $\mcV^1=\{I_n\}$. For $m\ge 2$, we divide each edge of the unit simplex $\Delta^{n-1}$ into $m$ equal segments of length $\sqrt{2}/m$ and partition the simplex into $m^{n-1}$ smaller simplices. For each smaller simplex, we form a matrix whose columns are the vertices of this simplex and include this matrix as a member of $\mcV^m$. 

% Let $\mcV^m$ be the collection of matrices constructed in the proof of Theorem~\ref{cor:strictly_pos_new} that are associated with the nonnegative orthant. 
% \textcolor{blue}{Let $\mcV^m$ be the collection of matrices arising from the partition
% of the unit simplex in the proof of Theorem~\ref{cor:strictly_pos_new},
% where the columns of each $V\in\mcV^m$ are the vertices of the corresponding smaller simplex.}
Let $q(x)\defeq x^TQx$. 
Since $q$ is homogeneous and $Q$ is strictly copositive, we have $q(x^*/\|x^*\|_2) > 0$ for all $x^* \in \textcolor{black}{\Delta^{n-1}}$. Letting $\delta(\cdot;\cdot)$ be defined as in Theorem~\ref{thm:degd_homo_new}, it follows that $\delta(q;\, x^*/\|x^*\|_2) > 0$ for all
$x^* \in \textcolor{black}{\Delta^{n-1}}$. It is straightforward to check that $\delta_{\min}\defeq \inf_{x^* \in \textcolor{black}{\Delta^{n-1}}}
\|x^*\|_2\,\delta(q; x^*/\|x^*\|_2) > 0$. 
For any integer $m \geq \lceil\sqrt{2n}/\delta_{\min}\rceil$ and any matrix $V \in \mcV^m$, $\conv\{V\}$ is a polytope with edge length $\sqrt{2}/m$, so for any column $v \in \textcolor{black}{\Delta^{n-1}}$ of $V$, we have
\[
\|V - v\mathbf{1}^T\|_F \leq \sqrt{n} \cdot \frac{\sqrt{2}}{m} \leq \delta_{\min}
\leq \|v\|_2\,\delta(q; v/\|v\|_2).
\]
Thus, $\left\|V/\|v\|_2 - (v/\|v\|_2)\mathbf{1}^T\right\|_F \leq \delta(q; v/\|v\|_2)$, and by Theorem~\ref{thm:degd_homo_new}, the polynomial $q(V x/\|v\|_2) = x^TV^TQVx/\|v\|_2^2$ 
has nonnegative coefficients, hence $V^TQV \geq 0$.
Since the same argument holds for each matrix in $\mcV^m$, we conclude that 
$Q \in \mcPN_n(\mcV^m)$.
\end{proof}

% \begin{theorem}\label{thm:copositive_new}
% For a symmetric matrix $Q\in\reals^{n\times n}$, if $(x^*)^TQx^*>0$ for some $x^*\in \textcolor{black}{\Delta^{n-1}}$, then there exists $\delta>0$ such that for any matrix $V\in\reals^{n\times n}$ satisfying $\|V-x^*\mathbf{1}^T\|_{F}\le \delta$, we have $V^TQV\ge 0$.
% In particular, the above statement holds for any
% \[\delta\le \hat{\delta}(Q;x^*)\defeq \|x^*\|_2\delta(q;\hat{x}),\]
% where $q(x)\defeq x^TQx, \hat{x}\defeq x^*/\|x^*\|_2$, and $\delta(p;x)$ is defined as in Theorem~\ref{thm:degd_homo_new}.
% \end{theorem}

% \begin{proof}
% Let $q(x)\defeq x^TQx$ and $\hat{x}\defeq x^*/\|x^*\|_2$. Since $q$ is homogeneous and $(x^*)^TQx^*>0$, we have $q(\hat{x})>0$. For any matrix $V\in\reals^{n\times n}$ satisfying $\|V-x^*\mathbf{1}^T\|_{F}\le \|x^*\|_2\delta(q;\hat{x})$, 
% we have
% $\|\hat{V}-\hat{x}\mathbf{1}^T\|_{F}\le \delta(q;\hat{x})$, where $\hat{V}=V/{\|x^*\|_2}$. By Theorem~\ref{thm:degd_homo_new}, the polynomial $q(\hat{V}x)~=~x^T\hat{V}^TQ\hat{V}x$ has nonnegative coefficients. Therefore, $\hat{V}^TQ\hat{V}\ge 0$, and hence $V^TQV\ge 0$.

% \end{proof}
% % Based on the above theorem, we further establish a disjunctive Positivstellensatz for matrix copositivity.
% \begin{theorem}[Disjunctive Positivstellensatz for Matrix Copositivity]\label{thm:psatz_copositive_full}
% For any fixed dimension $n$, there exists an explicit family of simplicial disjunctions
% $\mcV^m$ of the nonnegative orthant, indexed by $m$, such that:
% \begin{enumerate}
%     \item[(i)] For any matrix $Q$ admitting a $P+N$ decomposition, i.e., $Q = P + N$ with 
%     $P \succeq 0$ and $N \geq 0$, we have $Q \in \mcPN_{n}(\mcV^m)$ for any $m$.
%     \item[(ii)] For any strictly copositive matrix $Q \in \reals^{n \times n}$, we have 
%     $Q \in \mcPN_{n}(\mcV^m)$ for some $m$.
% \end{enumerate}
% \end{theorem}

% \begin{proof}
% Let $\mcV^m$ be constructed as in the proof of Theorem~\ref{cor:strictly_pos_new}, where we present $m^{n-1}$ matrices so that the corresponding simplicial cones cover the nonnegative orthant. Write $\mcV^m = \{V_1, \ldots, V_{r_m}\}$.

% \emph{Proof of (i).} Write $Q = P + N$ with $P \succeq 0$ and $N \geq 0$.
% For each $k = 1, \ldots, r_m$, we have $V_k^T Q V_k = V_k^T P V_k + V_k^T N V_k$.
% Since $V_k^T P V_k \succeq 0$ and $V_k^T N V_k \geq 0$ elementwise, because $V_k$ is elementwise nonnegative by construction, we have $Q \in \mcPN_{n}(\mcV^m)$.

% \emph{Proof of (ii).} Let $\hat{\delta}(Q; x^*)$ be as in Theorem~\ref{thm:copositive_new}.
% Since $Q$ is strictly copositive, the quantity $\delta_{\min} := \min_{x^* \in \textcolor{black}{\Delta^{n-1}}}
% \hat{\delta}(Q; x^*) > 0$.
% For any matrix $V \in \mcV^m$, the simplex $\cone\{V\} \cap \textcolor{black}{\Delta^{n-1}}$
% has edge length $\sqrt{2}/m$. Let $v\in\textcolor{black}{\Delta^{n-1}}$ be any column of $V$, we have
% \[
%     \|V - v\mathbf{1}^T\|_F \leq \sqrt{n} \cdot \frac{\sqrt{2}}{m} \leq \delta_{\min} 
%     \leq \hat{\delta}(Q; v),
% \]
% whenever $m \geq \lceil \sqrt{2n} / \delta_{\min} \rceil$. 
%  By Theorem~\ref{thm:copositive_new}, this implies $V^TQV\ge 0$. Since the same argument holds for each matrix in $\mcV^m$, we conclude that $Q\in \mcPN_n(\mcV^{m})$. 
% \end{proof}

\begin{remark}
We make two observations regarding Theorem~\ref{thm:psatz_copositive_full} and Definition~\ref{def:copos_disos}:
\begin{enumerate}
\item [(i)]
By the construction in the proof of Theorem~\ref{thm:psatz_copositive_full}, the first simplicial disjunction $\mcV^1$ contains only the identity matrix, 
and hence checking membership in $\mcPN_n(\mcV^1)$ reduces exactly to the $P+N$ test.
Moreover, the tests for copositivity in the higher levels of the hierarchy all dominate the $P+N$ test. Indeed, consider an arbitrary integer $m\ge 1$. If $Q = P + N$ with $P \succeq 0$ and $N \geq 0$, then for each matrix $V$ in $\mcV^m$, $V^T Q V = V^T P V + V^T N V$, where $V^T P V \succeq 0$, and $V^T N V \geq 0$ as the entries of $V$ are nonnegative by construction. Hence, $Q \in \mcPN_n(\mcV^m)$.

\item [(ii)]
Setting $P_k = 0$ in Definition~\ref{def:copos_disos} reduces the condition 
$V_k^T Q V_k= P_k + N_k$ to $V_k^T Q V_k~\geq~0$, 
which coincides with the sufficient condition for copositivity in~\cite[Theorem 1]{BUNDFUSS2008}.
However, by doing so, no level of the resulting hierarchy would dominate the $P+N$ test. See also Section~\ref{sec:num_cop} for some related comparisons in numerical experiments.
% linear optimization over $\mcPN_n(\mcV)$ reduces to a linear program, 
% making it cheaper to solve but potentially requiring a finer partition of the unit simplex to certify 
% copositivity compared to our approach; see Section~\ref{sec:num_cop} for numerical experiments.
\end{enumerate}
\end{remark}

% \begin{proof}
% Initialize with $r=1$ and $\mcV^1=\{I_n\}$. Suppose we have already constructed $\mcV^r=\{V_k\}_{k=1,\ldots,r}$. For each $k$, let $\Delta^k=\conv\{V_k\}$. Select an index $\ell\in\{1,\ldots,r\}$ such that $\Delta^{\ell}$ has the longest edge, denoted by $(v_{\ell,i},v_{\ell,j})$. Set $w=(v_{\ell,i}+v_{\ell,j})/2$ and define
% \[\begin{aligned}
% V_{\ell}^+&=[v_{\ell,1}\ \ldots\ v_{\ell,i-1}\ v_{\ell,i+1}\ \ldots\ v_{\ell,n}\ w],\\
% V_{\ell}^-&=[v_{\ell,1}\ \ldots\ v_{\ell,j-1}\ v_{\ell,j+1}\ \ldots\ v_{\ell,n}\ w].\\
% \end{aligned}\]
% This bisects $\Delta^{\ell}$ along its longest edge into two smaller simplices $\Delta^\ell_+=\conv\{V_\ell^+\}$ and $\Delta^\ell_{-}=\conv\{V_\ell^-\}$. We then update 
% \[\mcV^{r+1}=\left(\mcV^r-\{V_{\ell}\}\right)\cup\{V_{\ell}^+,V_{\ell}^-\}.\]
% Since $Q$ is strictly copositive, the quantity
% $\delta_{\min}\defeq\min_{x\in \textcolor{black}{\Delta^{n-1}}} \delta(Q;x)>0$. Therefore, for $r$ sufficiently large, each simplex $\Delta^k$ induced by the matrix $V^k$ in $\mcV^r$ has diameter at most $\delta_{\min}/\sqrt{n}$. Fix any $k$ and $x^*\in\Delta^k$. We have $\|V_k-x^*\mathbf{1}^T\|_F\le \delta_{\min}\le \delta(Q;x^*)$. By Theorem~\ref{thm:copositive_new}, it follows that $V_k^TQV_k\ge 0$ for all $k=1,\ldots,r$, and thus $Q\in\mcC_n(\mcV^r)$.
% \end{proof}

% \begin{remark}
% Theorem~\ref{thm:copositive_new} parallels~\cite[Lemma 7]{BUNDFUSS2008}. We make a step further in Theorem~\ref{thm:psatz_copositive_full} by constructing an explicit family of simplicial disjunctions. Consequently, the algorithm proposed in~\cite{BUNDFUSS2008} can adopt a fixed partitioning strategy for the unit simplex with guaranteed convergence. This finding aligns with~\cite[Proposition 1]{BUNDFUSS2008}.
% \end{remark}

\section{A Spatial Branch-and-Bound Framework}\label{sec:bnb}

% While the Positivstellens\"atze in previous sections show that one can always find low-degree proofs of nonnegativity of a polynomial by uniformly dividing the space, in practice it is more efficient to perform this division \emph{adaptively}. In this section, we show how to do this by integrating the disjunctive sum of squares method into a spatial branch-and-bound (SBB) framework~\cite[Chapter~IV]{horst_global_1996}.
% Our framework prioritizes the division of the subregions with the weakest lower bound.
While the Positivstellens\"atze in previous sections show that one can always find low-degree disjunctive sos proofs of nonnegativity of a polynomial by uniformly dividing the space, in practice it is more efficient to perform this division \emph{adaptively}. In this section, we describe how this can be done by integrating our framework into a spatial branch-and-bound (SBB) algorithm.
% integrate our framework into a spatial branch-and-bound (SBB) algorithm (see, e.g.,~\cite[Chapter~IV]{horst_global_1996}), which adaptively refines the subregion with the weakest lower bound.
% In this section, we integrate the disjunctive sum of squares method into a spatial branch-and-bound (SBB) framework~\cite{horst_global_1996} to \emph{adaptively} partition the region where nonnegativity of a polynomial needs to be certified. While the Positivstellens\"atze in previous sections show that one can always find low-degree proofs of nonnegativity by uniformly dividing the space, in practice it is more efficient to perform this division adaptively. 
% Our SBB approach prioritizes the division of the subregions with the weakest lower bound. 
% the theorems in earlier sections provide families of algebraic disjunctions which uniformly divide the region into subregions, this branch-and-bound framework iteratively 
% refines one subregion at a time. 
% such partition in only one subregion at a time. 
% adaptive region splitting based on various criteria. 
% Our approach is inherently flexible and supports diverse strategies for partitioning the feasible set and bounding over subregions using different Positivstellens\"atze. 
We present SBB algorithms to certify polynomial nonnegativity and matrix copositivity in Section~\ref{sec:adaptive_grid} and Section~\ref{sec:bnb_simplex}, respectively, and evaluate their numerical performance in Section~\ref{sec:num}. 
% Although we focus on the algorithm that performs best in our experiments, 


% computing lower bounds over subregions, based on different Positivstellens\"atze.

% As a brief overview, the SBB algorithm is a widely-used method to solve optimization problems of the form:
% \begin{equation*}
% 	\begin{array}{cl}
% 		\min\limits_{x\in S} & p(x)\\
% 		\st & g_i(x)\ge 0, \ i=1,\ldots,m,
% 	\end{array}    
% \end{equation*}
% where $S$ is the set of candidate solutions. 

% The algorithm iteratively partitions the solution space into smaller regions, a process known as branching. Branching strategies determine how regions are split, with one common approach selecting the region with the smallest lower bound. Other strategies may depend on specific properties of the decision variables or the problem structure. In the bounding step, the algorithm estimates upper and lower bounds of the objective function over the subregions, often through convex relaxations for the lower bound and feasible solutions for the upper bound. This iterative refinement generates a sequence of approximate solutions converging to the global optimum. By iteratively refining these bounds and pruning subregions, the algorithm generates a sequence of approximate solutions for the optimization problem.

% Other important steps include pruning, which prevents exploration of the sub-optimal space, and convergence check.


% We demonstrate the combination of the disjunctive sum of squares method with spatial branch-and-bound algorithm below. Although we focus on the algorithm that performs best in our experiments, this combination is inherently flexible to support diverse strategies for partitioning the feasible region and estimating lower bounds in each subregion, based on different sufficient conditions for certifying nonnegativity of polynomials.

% There are various ways to combine the algebraic disjunctions with the branching process. Instead of using the fixed choice of algebraic disjunctions constructed by the theorems, the algorithms smartly partition the regions with the smallest lower bound. Additionally, there are different ways to estimate lower bounds on each region. We present algorithms that are easy to implement and achieve good performance.

\subsection{A SBB Framework for Polynomial Minimization}
\label{sec:adaptive_grid}
In this section, we apply a branch-and-bound algorithm to approximate the minimum value $p_{\min}$ of a form $p\in H_{n,d}$ over the unit sphere $\mbS^{n-1}$ by adaptively dividing the sphere into subregions using simplicial disjunctions.
% since bisecting a simplicial cone is easy to implement and always yields two
% simplicial cones that exactly partition the original one.
% \textcolor{black}{
% In this section, we apply a branch-and-bound algorithm to approximate the minimum value $p_{\min}$ of a form $p\in H_{n,d}$ over the unit sphere $\mbS^{n-1}$, using simplicial disjunctions to adaptively partition the sphere into subregions.}
{We use simplicial disjunctions because it is straightforward to partition a simplicial cone into two smaller ones.
% Simplicial disjunctions are a natural choice for this purpose: it is straightforward to exactly partition a simplicial cone
% bisect a simplicial cone in a way that produces two simplicial cones that exactly partition the original one.
% bisecting a simplicial cone is straightforward and always produces two simplicial cones that exactly partition the original one.
}
We consider two choices for the initial simplicial disjunction:
\begin{enumerate}[label=(\arabic*)]
    \item $2^{n-1}$ matrices, each of the form\footnote{Since even-degree homogeneous polynomials are symmetric with respect to the origin, it suffices to certify nonnegativity on a hemisphere.} $\begin{bmatrix}\pm e_1 & \ldots & \pm e_{n-1} & e_n\end{bmatrix}$.
    % Construct $2^{n-1}$ matrices, each of the form $[\pm e_1\ \ldots\ \pm e_{n-1}\ e_n]$. Since even-degree homogeneous polynomials are symmetric with respect to the origin, it suffices to certify nonnegativity on a hemisphere.
    \item $n+1$ matrices, constructed as follows: Consider $n+1$ vertices of a regular $n$-simplex
\[c_i=\sqrt{1+n^{-1}}e_i-n^{-3/2}(\sqrt{n+1}-1)\mathbf{1},\quad i=1,\ldots,n,\quad c_{n+1}=-n^{-1/2}\mathbf{1}.\]
Each matrix is formed by taking $n$ of the $n+1$ vectors above as its columns.
% \[V_i=[v_{1}\ \ldots\ v_{i-1}\ v_{i+1}\ \ldots\ v_{n+1}],\quad i=1,\ldots,n+1.\]
% $$\{v_1,\ldots,v_{n+1}\}-\{v_i\}$ for $i=1,\ldots,n+1$.
% in the initialization represents a spherical polygon with vertices $\{v_1,\ldots,v_{n+1}\}-\{v_i\}$ for $i=1,\ldots,n+1$.
\end{enumerate}
% At each iteration, we update global lower and upper bounds as subregions are refined. 
% the global lower and upper bounds are updated dynamically as the branching and bounding processes refine subregions. 
% For branching, we select the subregion with the smallest lower bound, 

\begin{figure}[!tb]
    \centering
    \includegraphics[width=0.85\textwidth]{figs/branching.pdf}
    \caption{Branching step for $\cone\{V\}=\cone\{v_1,v_2,v_3\}$. Left: $\mcQ=\cone\{V\}\cap\mbS^{2}$ is the spherical triangle bounded by geodesic arcs (solid); planar chords of~$\cone\{V\}$ shown dashed. Right: the longest arc~$(v_1,v_3)$ is bisected at its geodesic midpoint~$w$, splitting~$\mcQ$ into~$\mcQ^+$ and~$\mcQ^-$.}
    \label{fig:branching}
\end{figure}

\paragraph{Bounding.}
Each subregion of the sphere takes the form $\mcQ = \cone\{V\} \cap \mbS^{n-1}$ for some invertible
matrix $V\in \reals^{n \times n}$. We compute {a} lower bound on $p$ over $\mcQ$ by solving
\begin{equation}\label{eq:poly_bnb}
	\phi(V) \defeq \begin{array}[t]{cl}
		\max\limits_{\gamma\in\reals} & \gamma\\
		\st & p(V(x.^2))-\gamma\|V(x.^2)\|_2^d\in\HSOS_{n,2d}.
	\end{array}
\end{equation} 
Note that any feasible $\gamma$ in problem~\eqref{eq:poly_bnb} is a valid lower bound on~$p$ over~$\mcQ$. 
{Indeed, any point $z\in\cone\{V\}$ can be written as $z=V(x.^2)$ for some $x\in\reals^n$, so the constraint in~\eqref{eq:poly_bnb} implies $p(z)-\gamma\|z\|_2^d\ge 0$ for all $z\in\cone\{V\}$. In particular, $p(z)-\gamma\ge 0$ for all $z\in\cone\{V\} \cap \mbS^{n-1}$.} 
% Indeed, for any polynomial $q\in H_{n,d}$ and any matrix $V\in\reals^{n\times n}$, the condition $q(V(x.^2))\in\HSOS_{n,2d}$ implies nonnegativity of $q$ over $\cone\{V\}$. 
% This is because any point $z\in\cone\{V\}$ can be written as $z=V\bar{z}$ with $\bar{z}\ge 0$, and such vector $\bar{z}$ can be written as $\bar{z}=x.^2$ for some $x\in\reals^n$. 

% We run projected gradient steps to obtain an upper bound, specifically, we iterate
 % on $p_{\min}$ over $\mcQ$ by running projected gradient steps: starting from a column of~$V$, we iterate
{To obtain an upper bound on~$p$ over~$\mcQ$, we run $K$ projected gradient descent steps (where $K$ is chosen by the user): starting from a column of~$V$, we iterate}
\begin{equation}\label{eq:pgd_poly}
   y^k=\Pi_{\cone\{V\}}(x^k - \beta \nabla p(x^k)),\quad x^{k+1} = \frac{y^k}
    {\|y^k\|_2},\quad k=0,\ldots,K-1,
\end{equation}
where $\Pi_{\cone\{V\}}$ denotes the projection onto $\cone\{V\}$ and $\beta>0$ is a chosen stepsize.

% the stepsize $\beta$ is 
% set proportional to $\max_{1\le k,l\le n}\|v_k - v_l\|_2$. 
% We compute an upper bound on $p_{\min}$ by evaluating~$p$ at points within~$\mcQ$; specifically, we run a few projected gradients descent steps on $p$ over $\mcQ$ starting from a column of~$V$. 
% The lower bound is obtained by solving
% \begin{equation}\label{eq:poly_bnb}
% 	\phi(V) \defeq \begin{array}[t]{cl}
% 		\max\limits_{\gamma\in\reals} & \gamma\\
% 		\st & p(V(x.^2))-\gamma\|V(x.^2)\|_2^d\in\HSOS_{n,2d}.
% 	\end{array}
% \end{equation} 


\paragraph{Branching.}
At each iteration, we select a subregion
\begin{equation*}
\mcQ=\cone\{V\}\cap\mbS^{n-1}
\end{equation*}
 with the smallest lower bound, where $V = [v_1 \,\ldots\, v_n] \in\reals^{n\times n}$ is some invertible matrix. Letting $(i,j)\in \operatornamewithlimits{argmax}_{1 \le k, l \le n} \|v_k-v_l\|_2$, we 
% and split it by bisecting its longest arc. 
% Specifically, we let the selected subregion be
% \begin{equation*}
% \mcQ=\cone\{V\}\cap\mbS^{n-1},
% \end{equation*}
% for some invertible matrix $V = [v_1 \,\ldots\, v_n] \in\reals^{n\times n}$.
% Denoting the longest arc by the pair $(v_i, v_j)$, we 
set~$w= (v_{i}+v_{j})/\|v_{i}+v_{j}\|_2$ and form two further subregions
\[
    \mcQ^+=\cone\{V^+\}\cap \mbS^{n-1},\quad \mcQ^-=\cone\{V^-\}\cap \mbS^{n-1},\]
where $V^+=[v_{1}\, \ldots\, v_{i-1}\, v_{i+1}\, \ldots\, v_{n}\, w]$ and $V^-=[v_{1}\, \ldots\, v_{j-1}\, v_{j+1}\, \ldots\, v_{n}\, w]$.
{
% The point~$w$ is the geodesic midpoint of the arc~$(v_i,v_j)$ on~$\mbS^{n-1}$, since the great-circle arc between two unit vectors admits the parametrization~$\gamma(t) = ((1-t)v_i+t v_j)/\|(1-t)v_i + t v_j\|_2$ for $t \in [0,1]$, and $\gamma(1/2) = w$.
Figure~\ref{fig:branching} illustrates an example of this construction when $n=3$.}

% \paragraph{Convergence.}
\textcolor{black}{The 
overall branch-and-bound scheme is presented in Algorithm~\ref{alg:bnb}. 
% The convergence of Algorithm~\ref{alg:bnb} follows from the abstract branch-and-bound framework of~\cite[Theorem~5.26]{locatelli2013global}. 
One can utilize Theorem~\ref{thm:degd_homo_new} together with convergence results for branch-and-bound methods (see, e.g.~\cite[Theorem~5.26]{locatelli2013global} or~\cite[Chapter~IV]{horst_global_1996}) to show that for any prescribed tolerance $\epsilon_{\text{tol}}>0$, Algorithm~\ref{alg:bnb} is guaranteed to terminate.}



%It suffices to check the exhaustiveness and exactness in the limit properties~\cite[Definitions~5.5 and~5.4]{locatelli2013global}.Our bisection rule for branching guarantees exhaustiveness, because the subregions shrink to a single point along any infinite branch of the branch-and-bound tree. Exactness follows from Theorem~\ref{thm:degd_homo_new}: as a subregion shrinks, $\phi(V)$ converges to the minimum of $p$ over that subregion. Hence, Algorithm~\ref{alg:bnb} terminates in finitely many iterations for any prescribed tolerance $\epsilon_{\text{tol}}>0$.}
% which guarantees finite-time termination at any prescribed tolerance $\epsilon_{\text{tol}}>0$
% provided that the subdivision process is exhaustive
% and the bounding procedure is exact in the limit.
% The longest-arc bisection rule ensures exhaustiveness~\cite[Definition~5.5 and Remark~5.7]{locatelli2013global},
% because the subregions shrink to a single ray along any infinite branch of the branch-and-bound tree.

 % Theorems~\ref{cor:strictly_pos_new} and~\ref{thm:degd_homo_new} together imply that any positive definite form~$p$ has nonnegative coefficients with respect to a simplicial disjunction~$\mcV$ with sufficiently fine subregions. 
 % It follows that $p(V(x.^2))\in\HSOS_{n,2d}$ for any matrix~$V$ in such simplicial disjunction~$\mcV$. 
 % \textcolor{black}{Bisecting the longest chord~$\|v_i-v_j\|_2$ is an exhaustive simplicial subdivision rule, so any nested sequence of subcones generated by the algorithm narrows to a single ray~\cite[Proposition~IV.2 and~\S IV.3.2]{horst_global_1996}; the standard branch-and-bound convergence theorem~\cite[Theorem~IV.3]{horst_global_1996} then implies that the global lower bound converges to $p_{\min}$ as the algorithm progresses.}
 
 % this guarantees that the minimum lower bound across all subregions becomes arbitrarily close to $p_{\min}$.
 % the minimum value of $p$.
 


% $p(Vx)$ has nonnegative coefficients implies $p(V(x.^2))\in\HSOS_{n,2d}$, so~\eqref{eq:poly_bnb} provides a tighter lower. Moreover, $p(V(x.^2))\in\HSOS_{n,2d}$ still implies nonnegativity of $p$ over $\cone\{V\}$, since any point $z\in\cone\{V\}$ can be written as $z=V\alpha$ with $\alpha\ge 0$, and such $\alpha$ can be represented componentwise as $\alpha=x.^2$ for some $x\in\reals^n$. 

% using two approaches:
% \begin{enumerate}
%     \item Following Definition~\ref{def:disos_new}, we solve
% \begin{equation}\label{eq:poly_bnb_free}
% 	\begin{array}{cl}
% 		\max\limits_{\gamma\in\reals} & \gamma\\
% 		\st & \text{all coefficients of } p(Vx)-\gamma\|Vx\|_2^d\ \text{ are nonnegative.}\\
% 	\end{array}
% \end{equation} 
% This can be done by bisection and is thus optimization-free. 
% \item To obtain tighter bounds, we incorporate the sos technique and solve
% \begin{equation}\label{eq:poly_bnb}
% 	\begin{array}{cl}
% 		\max\limits_{\gamma\in\reals} & \gamma\\
% 		\st & p(V(x.^2))-\gamma\|(V(x.^2)\|_2^d\in\HSOS_{n,2d},
% 	\end{array}
% \end{equation} 
% where $x.^2\defeq (x_1^2,\ldots,x_n^2)$ for a vector $x\in\reals^n$. For any polynomial $p\in H_{n,d}$ and any matrix $V\in\reals^{n\times n}$, the condition that 
% $p(Vx)$ has nonnegative coefficients implies $p(V(x.^2))\in\HSOS_{n,2d}$, so~\eqref{eq:poly_bnb} provides a tighter lower bound than~\eqref{eq:poly_bnb_free}. Moreover, $p(V(x.^2))\in\HSOS_{n,2d}$ still implies nonnegativity of $p$ over $\cone\{V\}$, since any point $z\in\cone\{V\}$ can be written as $z=V\alpha$ with $\alpha\ge 0$, and such $\alpha$ can be represented componentwise as $\alpha=x.^2$ for some $x\in\reals^n$. 
% % Therefore, we solve 
% % the following optimization problem, which provides a tighter lower bound than~\eqref{eq:poly_bnb_free}:
% \end{enumerate}

% We now present the full branch-and-bound scheme in Algorithm~\ref{alg:bnb} below.



\begin{algorithm}[H]
\caption{Spatial Branch-and-Bound Algorithm for Minimizing a Form}
\label{alg:bnb}
\begin{algorithmic}
\Require A polynomial $p\in H_{n,d}$, a tolerance $\epsilon_{\text{tol}}>0$.
\Ensure A lower bound $L$ on the minimum value of $p$ over $\mbS^{n-1}$ and a point $\bar{x}\in\mbS^{n-1}$ such that 
$p(\bar{x})-L\le \epsilon_{\text{tol}}(1+|L|+|p(\bar{x})|)$.
\Initialization Generate an initial simplicial disjunction $\{V_i\}_{i=1,\ldots,r_0}$.
\State $\mcT\gets\{V_i\mid i=1,\ldots,r_0\}$. 
\Comment{Create a branch-and-bound tree}
% with $V_i\in\reals^{n\times n}$ invertible. 
\State $l_i\gets \phi(V_i)$, $\quad i=1,\ldots,r_0$. \Comment{Local lower bound computed by solving~\eqref{eq:poly_bnb}}

\State $u_i\gets\min_{1\le j\le n}\ p(V_ie_j)$, $\quad i=1,\ldots,r_0$. \Comment{Local upper bound}
\State $L\gets\min_{1\le i\le r_0}\{l_i\},\quad U\gets\min_{1\le i\le r_0}\{u_i\}$. 
\Comment{Global lower and upper bounds}
\While{$U-L> \epsilon_{\text{tol}}(1+|L|+|U|)$}
\State $V\gets \argmin_{\hat{V}\in \mcT} \phi(\hat{V}),\quad \mcT\gets \mcT\setminus\{V\}$. 
% \Comment{Update the global lower bound.}
\State $v_1, \ldots, v_n \gets$ columns of $V$.
\State $(i, j) \gets \operatornamewithlimits{argmax}_{1 \le k, l \le n} \|v_k-v_l\|_2,\quad w \gets (v_i + v_j)/\|v_i + v_j\|_2$.
\State $V^+ \gets [v_1\, \ldots\, v_{i-1}\, v_{i+1}\, \ldots\, v_n\, w],\quad V^- \gets [v_1\, \ldots\, v_{j-1}\, v_{j+1}\, \ldots\, v_n\, w]$. \Comment{Branch}
\State $\mcT \gets \mcT \cup \{V^+,V^-\}$.
% \State \Call{CreateNode\underline{ }Poly}{$V^+$}. \Comment{See Procedure~\ref{alg:create_node}.}
% \State \Call{CreateNode\underline{ }Poly}{$V^-$}.
\State $L\gets \min_{\hat{V}\in \mcT} \phi(\hat{V})$.  \Comment{Update global lower bound}
\For{$\hat{V}\in\{V^+,V^-\}$}\Comment{Update global upper bound}
\State $x^K \gets$ $K$-th iterate of~\eqref{eq:pgd_poly} with $x^0=w$.
\State $u\gets p(x^K)$.
% \State $u,x\gets \psi(V)$. \Comment{See Procedure~\ref{alg:create_node}.}
\If{$u<U$} $U\gets u,\quad \bar{x}\gets x^K$.\EndIf
\EndFor
% \State $U\gets \min\{U,\psi(V^+),\psi(V^-)\}$. 
% \State $V^+,V^- \gets$ split $\cone\{V\}\cap\mbS^{n-1}$ along its longest arc into \[\mcQ^+=\cone\{V^+\}\cap\mbS^{n-1},\quad \mcQ^-=\cone\{V^-\}\cap\mbS^{n-1}.\]
% \State $x^+ \gets$ run $T$ steps of~\eqref{eq:pgd_poly} 
% with $V=V^+$ and $\alpha = 0.5\delta$, 
% starting from~$w$.
% \State $l^+\gets \phi(V^+)$, \quad $u^+ \gets p(x^+)$.
% \If{$l^+< U$} $\mcS\gets \mcS\cup\{(V^+, l^+, u^+)\}$. 
% \EndIf
% \If{$u^+<U$} $U\gets u^+,\quad \hat{x}\gets x^+$
% \EndIf
% \State $x^- \gets$ run $T$ steps of~\eqref{eq:pgd_poly} 
% with $V=V^-$ and $\alpha = 0.5\delta$, 
% starting from~$w$.
% \State $l^-\gets \phi(V^-)$, \quad $u^- \gets p(x^-)$.
% \If{$l^-< U$} $\mcS\gets \mcS\cup\{(V^-, l^-, u^-)\}$.
% \EndIf
% \If{$u^-<U$} $U\gets u^-,\quad \hat{x}\gets x^-$
% \EndIf
\EndWhile
\State \Return $L,\bar{x}$
\end{algorithmic}
\end{algorithm}

% \vspace{-4ex}
% \floatname{algorithm}{Procedure}
% \begin{algorithm}[H]
% \caption{$\psi(V)$: An upper bound on~$p$ over $\cone\{V\}\cap\mbS^{n-1}$ for Algorithm~\ref{alg:bnb}.}
% \label{alg:create_node}
% \begin{algorithmic}
% \Require An invertible matrix $V\in\reals^{n\times n}$.
% % \Ensure Updated $U$, $\mcS$ and a point $\hat{x}\in\mbS^{n-1}$ satisfying $p(\hat{x})=U$.
% \State $\delta\gets\max_{1\le k,l\le n}\|v_k - v_l\|_2$.

% \State $x \gets$ run $T$ steps of~\eqref{eq:pgd_poly} with $\beta=0.5\delta$, starting from $w$.
% \State $u\gets p(x)$.
% \State \Return $u,x$
% % \State $\hat{l} \gets \phi(\hat{V})$, as in~\eqref{eq:poly_bnb}.
% % \If{$\hat{l} < U$} $\mcT \gets \mcT \cup \{(\hat{V}, \hat{l}, \hat{u})\}$.
% % \EndIf
% % \If{$\hat{u} < U$} $U \gets \hat{u}$, \quad $\hat{x} \gets x$.
% % \EndIf
% % % \State \Return $U$, $\hat{x}$, $\mcS$.
% \end{algorithmic}
% \end{algorithm}
% \floatname{algorithm}{Algorithm}


\subsection{A SBB Framework for Matrix Copositivity}\label{sec:bnb_simplex}
In this section, we focus on the problem of testing copositivity of a matrix $Q\in\reals^{n\times n}$, which reduces to computing the minimum value $p_{\min}$ of a quadratic form $p(x)=x^TQx$ over the unit simplex $\textcolor{black}{\Delta^{n-1}}$.
% Recall that certifying matrix copositivity reduces to minimizing a quadratic form over the unit simplex $\textcolor{black}{\Delta^{n-1}}$.
We develop a branch-and-bound algorithm analogous to Algorithm~\ref{alg:bnb}, tailored to the unit simplex. This algorithm can be used more generally to solve copositive programs. We initialize the algorithm with a single subregion, the unit simplex $\textcolor{black}{\Delta^{n-1}}=\conv\{I_n\}$, corresponding to the simplicial disjunction $\{I_n\}$.

% , then divide it into smaller simplicies of the form $\conv\{V\}$ for some invertible matrix $V\in\reals^{n\times n}$.

\begin{figure}[!ht]
    \centering
    \includegraphics[width=0.85\textwidth]{figs/branching_simplex.pdf}
    \caption{Branching step for $\conv\{V\}=\conv\{v_1,v_2,v_3\}$. Left: $\mcQ=\conv\{V\}$ is the blue triangle; the unit simplex~$\textcolor{black}{\Delta^2}$ shown dashed. Right: the longest edge~$(v_1,v_2)$ is bisected at its midpoint~$w=(v_1+v_2)/2$, splitting~$\mcQ$ into~$\mcQ^+$ and~$\mcQ^-$.}
    \label{fig:branching_simplex}
\end{figure}

\paragraph{Bounding.}
Each subregion of the unit simplex takes the form $\mcQ = \conv\{V\}$ for some invertible
matrix $V~=~[v_1 \,\ldots\, v_n]\in\reals^{n \times n}$. Note that the minimum of~$p$ over~$\mcQ$ equals the optimal value of
\begin{equation*}
    \begin{array}{cl}
        \max\limits_{t\in\reals} & t\\
        \st & V^T(Q-tJ_n)V\in\mcC_n,
    \end{array}
\end{equation*}
where $J_n$ is the $n\times n$ all-one matrix and $\mcC_n$ is the cone of $n\times n$ copositive matrices. Using the inner approximations of $\mcC_n$ from Section~\ref{sec:copositive}, we compute a lower bound on~$p$ over $\mcQ$ by solving the following semidefinite program:
\begin{equation}\label{eq:DiSOS_lb_mat}
    \psi(V)\defeq \begin{array}[t]{cl}
        \max\limits_{t\in\reals,N\in\reals^{n\times n}} & t\\
        \st & V^T(Q-tJ_n)V-N\succeq 0,\\
            & N\ge 0.
    \end{array}
\end{equation}
% We run projected gradient steps to obtain an upper bound, similar to Section~\ref{sec:adaptive_grid}: starting from a column of~$V$, we iterate
\textcolor{black}{To obtain an upper bound on~$p$ over $\mcQ$, we run $K$ projected gradient descent steps starting from a column of~$V$:}
\begin{equation}\label{eq:pgd_copos}
    x^{k+1} = \Pi_{\conv\{V\}}(x^k - 2\beta Qx^k),\quad k=0,\ldots,K-1.
\end{equation}
Here, $\Pi_{\conv\{V\}}$ denotes the projection onto $\conv\{V\}$ 
and $\beta>0$ is the stepsize.
% the stepsize $\beta$ is set proportional to $\max_{1\le k,l\le n}\|v_k - v_l\|_2$. 
% One may also apply cheaper classes of conic programs at the cost of loosening the lower bound.
% For example, one may replace the positive semidefinite cone in~\eqref{eq:DiSOS_lb_mat} by the cone of diagonally dominant matrices, or remove it entirely and instead impose $V^T(Q-tJ)V\ge 0$.
\paragraph{Branching.}
At each iteration, we select a subregion with the smallest lower bound and split it by bisecting its longest edge. More specifically, if 
\begin{equation*}
\mcQ=\conv\{V\}
\end{equation*}
denotes the selected subregion, where $V = [v_1 \,\ldots\, v_n] \in\reals^{n\times n}$ is some invertible matrix, we let $(i,j)\in \operatornamewithlimits{argmax}_{1 \le k, l \le n} \|v_k-v_l\|_2$, and set~$w= (v_{i}+v_{j})/2$, forming two further subregions
\[
    \mcQ^+=\conv\{V^+\},\quad \mcQ^-=\conv\{V^-\},\]
where $V^+=[v_{1}\, \ldots\, v_{i-1}\, v_{i+1}\, \ldots\, v_{n}\, w]$ and $V^-=[v_{1}\, \ldots\, v_{j-1}\, v_{j+1}\, \ldots\, v_{n}\, w]$.
\textcolor{black}{Figure~\ref{fig:branching_simplex} illustrates an example of this construction when $n=3$.}

% For branching, we select the subregion with the smallest lower bound. 
% Since we partition $\textcolor{black}{\Delta^{n-1}}$ using simplicial disjunctions, we denote the selected subregion by $\mcQ=\conv\{V\}$ and split it by bisecting its longest edge. 
% denoted $\mcQ=\conv\{V\}$. We split $\mcQ$ by bisecting its longest edge: 
% Denoting the longest edge by the pair $(v_i,v_j)$, we set the point $w=(v_i+v_j)/2$ and form two subregions $\Delta^+=\conv\{V^+\}$ and $\Delta^-=\conv\{V^-\}$, where
% \[\begin{aligned}
% V^+&=[v_{1}\, \ldots\, v_{i-1}\, v_{i+1}\, \ldots\, v_{n}\, w],\\
% V^-&=[v_{1}\, \ldots\, v_{j-1}\, v_{j+1}\, \ldots\, v_{n}\, w].\\
% \end{aligned}\]


% \paragraph{Convergence.}
\textcolor{black}{
The 
overall branch-and-bound scheme is presented in Algorithm~\ref{alg:bnb_simplex}. 
By the same argument as in Section~\ref{sec:adaptive_grid}, one can show that for any prescribed tolerance $\epsilon_{\text{tol}}>0$, 
Algorithm~\ref{alg:bnb_simplex} terminates. If the algorithm returns a nonnegative lower bound $L$, then $Q$ is certified to be copositive. If it returns a point $\bar{x}\in\Delta^{n-1}$ with $\bar{x}^TQ\bar{x}<0$, then $Q$ is not copositive. In the remaining case, copositivity is inconclusive with the prescribed tolerance.}
 % \textcolor{black}{Bisecting the longest edge~$\|v_i-v_j\|_2$ is an exhaustive simplicial subdivision rule~\cite[Proposition~IV.2]{horst_global_1996}, so the diameters of the generated subsimplices tend to zero; combined with Theorem~\ref{thm:psatz_copositive_full}, the standard branch-and-bound convergence theorem~\cite[Theorem~IV.3]{horst_global_1996} implies that the global lower bound converges to $p_{\min}$ as the algorithm progresses.}
 % the minimum value of the quadratic form $x^TQx$ over the unit simplex.
 

% \begin{enumerate}
%     \item Following Definition~\ref{def:disos_new}, we solve
%     \begin{equation}\label{eq:cp_bnb_free}
%         \begin{array}{cl}
%             \max\limits_{t\in\reals} & t\\
%             \st & V^T(Q-tJ)V\ge 0.
%         \end{array}
%     \end{equation}
%     % which can be solved by bisection and is thus optimization-free.
%     \item To obtain tighter bounds, we incorporate the sos technique. As discussed in Section~\ref{sec:copositive}, we solve
%     \begin{equation}\label{eq:DiSOS_lb_mat}
%         \begin{array}{cl}
%             \max\limits_{t\in\reals,N\in\reals^{n\times n}} & t\\
%             \st & V^T(Q-tJ)V-N\succeq 0,\\
%                 & N\ge 0.
%         \end{array}
%     \end{equation}
% This follows from substituting $p(x)=x^TQx$ into the condition $p(V(x.^2))\in\HSOS_{n,4}$ to obtain $(x.^2)^TV^TQV(x.^2)\in\HSOS_{n,4}$, which is further equivalent to the decomposition $V^TQV=P+N$ with $P\succeq 0$ and $N\ge 0$ by~\cite[Section 5.3]{Parrilo2000}.
% \end{enumerate}
% We now present the full branch-and-bound scheme in Algorithm~\ref{alg:bnb_simplex} below.
% With this specialized branching and bounding procedure, the disjunctive sum of squares method is effectively adapted to this copositivity setting. We present the full branch-and-bound scheme below.



\begin{algorithm}[H]
\caption{Spatial Branch-and-Bound Algorithm for Testing Matrix Copositivity}
\label{alg:bnb_simplex}
\begin{algorithmic}
\Require A symmetric matrix $Q\in\reals^{n\times n}$, a tolerance $\epsilon_{\text{tol}}>0$.
\Ensure A lower bound $L$ on the minimum value of $x^TQx$ over $\textcolor{black}{\Delta^{n-1}}$ and a point $\bar{x}\in\textcolor{black}{\Delta^{n-1}}$ such that
$\bar{x}^TQ\bar{x}-L\le \epsilon_{\text{tol}}(1+|L|+|\bar{x}^TQ\bar{x}|)$.
\Initialization
$\mcT\gets\{I_n\}$. 
\Comment{Create a branch-and-bound tree}
\State $L\gets \psi(I_n),\quad U\gets \min_{1\le i\le n} e_i^TQe_i$. 
\Comment{Initialize global lower and upper bounds}
\While{$U-L> \epsilon_{\text{tol}}(1+|L|+|U|)$}
\State $V\gets \argmin_{\hat{V}\in\mcT} \psi(\hat{V}),\quad \mcT\gets \mcT\setminus\{V\}$. 
% \Comment{Update the global lower bound.}
\State $v_1, \ldots, v_n \gets$ columns of $V$.
\State $(i, j) \gets \operatornamewithlimits{argmax}_{1 \le k, l \le n} \|v_k-v_l\|_2,\quad w \gets (v_i + v_j)/2$.
\State $V^+ \gets [v_1\, \ldots\, v_{i-1}\, v_{i+1}\, \ldots\, v_n\, w],\quad V^- \gets [v_1\, \ldots\, v_{j-1}\, v_{j+1}\, \ldots\, v_n\, w]$. \Comment{Branch}
\State $\mcT \gets \mcT \cup \{V^+,V^-\}$.
% \State \Call{CreateNode\underline{ }Poly}{$V^+$}. \Comment{See Procedure~\ref{alg:create_node}.}
% \State \Call{CreateNode\underline{ }Poly}{$V^-$}.
\State $L\gets \min_{\hat{V}\in \mcT} \psi(\hat{V})$.  \Comment{Update global lower bound}
\For{$\hat{V}\in\{V^+,V^-\}$}\Comment{Update global upper bound}
\State $x^K \gets$ $K$-th iterate of~\eqref{eq:pgd_copos} with $x^0=w$.
\State $u\gets (x^K)^TQx^K$.
% \State $u,x\gets \psi(V)$. \Comment{See Procedure~\ref{alg:create_node}.}
\If{$u<U$} $U\gets u,\quad \bar{x}\gets x^K$.\EndIf
\EndFor
% \State $x^+\gets$ run $T$ steps of~\eqref{eq:pgd_copos} with $V=V^+$, $\alpha = 0.5\delta$, starting from $w$.
% \State $l^+\gets \psi(V^+),\quad u^+\gets (x^+)^TQx^+$.
% \If{$l^+< U$} $\mcS\gets \mcS\cup\{(V^+, l^+, u^+)\}$. 
% \EndIf
% \If{$u^+<U$} $U\gets u^+,\quad \hat{x}\gets x^+$
% \EndIf
% \State $x^-\gets$ run $T$ steps of~\eqref{eq:pgd_copos} with $V=V^-$, $\alpha = 0.5\delta$, starting from $w$.
% \State $l^-\gets \psi(V^-),\quad u^-\gets (x^-)^TQx^-$.
% \If{$l^-< U$} $\mcS\gets \mcS\cup\{(V^-, l^-, u^-)\}$.
% \EndIf
% \If{$u^-<U$} $U\gets u^-,\quad \hat{x}\gets x^-$
% \EndIf
% \State $U\gets\min\{U,u^+,u^-\}$. \Comment{Update the global upper bound.}
\EndWhile
\State \Return $L, \bar{x}$
\end{algorithmic}
\end{algorithm}

% \vspace{-4ex}
% \floatname{algorithm}{Procedure}
% \begin{algorithm}[H]
% \caption{\textsc{CreateNode\underline{ }Cop}: Create and insert a new node for Algorithm~\ref{alg:bnb_simplex}}
% \label{alg:create_node_cop}
% \begin{algorithmic}
% \Require An invertible matrix $V\in\reals^{n\times n}$.
% % \Ensure Updated $U$, $\mcS$ and a point $\hat{x}\in\mbS^{n-1}$ satisfying $p(\hat{x})=U$.
% \State $\delta\gets\max_{1\le k,l\le n}\|v_k - v_l\|_2$.
% \State $x \gets$ run $T$ steps of~\eqref{eq:pgd_copos} with $\beta=0.5\delta$, starting from $w$;\quad $u \gets x^TQx$.
% \State $l \gets \psi(V)$, as in~\eqref{eq:DiSOS_lb_mat}.
% \If{$l < U$} $\mcT \gets \mcT \cup \{(V, l, u)\}$.
% \EndIf
% \If{$u < U$} $U \gets u$, \quad $\hat{x} \gets x$.
% \EndIf
% % \State \Return $U$, $\hat{x}$, $\mcS$.
% \end{algorithmic}
% \end{algorithm}
% \floatname{algorithm}{Algorithm}

\subsection{Numerical Experiments}\label{sec:num}
In this subsection, we present some numerical experiments focusing on certifying nonnegativity of non-sos polynomials, and solving copositive programs with applications to nonconvex quadratic programming and the maximum clique problem in combinatorial optimization. 
% solving standard quadratic \textcolor{black}{programs} via the copositive cone, and computing the clique number of graphs. 
The code and data used for the numerical experiments are available in our GitHub repository.\footnote{\url{https://github.com/YH7422/DisjunctiveSOS}}
\subsubsection{Polynomial Minimization}\label{sec:num_poly}
We apply Algorithm~\ref{alg:bnb} to minimize 
several classical polynomials for which the standard sos test (see, e.g.,~\cite{Parrilo2003}) fails to return the minimum value. The examples, drawn from~\cite{Reznick2000, Delzell1980, ahmadi2017sum}, are listed in the following table. 

% In particular, the polynomial $P_m$ is a homogenized version of
% \[P(x_1,\ldots,x_5)=\left(\sum_{i=1}^5x_i\right)^2+\sum_{i=1}^5(x_i^2-1)^2,\]
% studied in~\cite[Section 5.1]{ahmadi2017sum}. This example corresponds to the infeasible \textsc{Partition} instance $\{1,1,1,1,1\}$ and is known to satisfy that there is no $\epsilon>0$ such that $P(x)-\epsilon$ is sos. The polynomial $D$ from~\cite{Delzell1980} is said to have a ``bad point", meaning that $\left(\sum_{i=1}^4 x_i^2\right)^k D(x)$ is not sos for any integer $k\ge 0$.
% proved to be not ``sos-refutable", i.e., there is no $\epsilon>0$ such that $P(x)-\epsilon$ is sos.
\begin{table}
  \centering
  \small
  \begin{tabular}{llll}\toprule
Name & $n$ & $d$ &\hfil Expression\\\midrule
    Motzkin & 3 & 6 & $M_h(x_1,x_2,x_3)=x_1^4x_2^2+x_1^2x_2^4-3x_1^2x_2^2x_3^2+x_3^6$\\
        Robinson-1 & 3 & 6 & \hspace{-1.6ex} $\begin{array}[t]{ll}R_1(x_1,x_2,x_3) =&x_1^6+x_2^6+x_3^6- (x_1^4x_2^2+x_1^2x_2^4+x_1^4x_3^2\\ &+x_1^2x_3^4+x_2^4x_3^2+x_2^2x_3^4)+3x_1^2x_2^2x_3^2\end{array}$\\
  Robinson-2 & 4 & 4& \hspace{-1.6ex} $\begin{array}[t]{ll}R_2(x_1,\ldots,x_4)=&x_1^2(x_1-x_4)^2+x_2^2(x_2-x_4)^2+x_3^2(x_3\\&-x_4)^2+2x_1x_2x_3(x_1+x_2+x_3-2x_4)\end{array}$\\
  Choi-Lam-1 & 4 & 4 & $        CL_1(x_1,\ldots,x_4)=x_1^2x_2^2+x_1^2x_3^2+x_2^2x_3^2+x_4^4-4x_4x_1x_2x_3$\\
  Choi-Lam-2 & 3 & 6 & $CL_2(x_1,x_2,x_3)=x_1^4x_2^2+x_2^4x_3^2+x_3^4x_1^2-3x_1^2x_2^2x_3^2$\\
  % Horn & 5&4&\hspace{-1.8ex} $\begin{array}[t]{ll}H(x_1,\ldots,x_5)=&\left(\sum_{i=1}^5x_i^2\right)^2-4(x_1^2x_2^2+x_2^2x_3^2+x_3^2x_4^2\\&+x_4^2x_5^2+x_5^2x_1^2)\end{array}$\\
  Lax & 5&4 &$L(x_1,\ldots,x_5)=\sum_{i=1}^5\prod_{j\neq i}(x_i-x_j)$ \\
  Schm\"udgen & 3 & 6&\hspace{-1.6ex} $\begin{array}[t]{ll}SC(x_1,x_2,x_3)=&200(x_1^3-4x_1x_3^2)^2+200(x_2^3-4x_2x_3^2)^2\\&+(x_2^2-x_1^2)x_1(x_1+2x_3)(x_1^2-2x_1x_3\\&+2x_2^2-8x_3^2)\end{array}$\\
\textsc{Partition} & 6 & 4& $P(x_1,\ldots,x_6)=\left(\sum_{i=1}^5x_i\right)^2x_6^2+\sum_{i=1}^5(x_i^2-x_6^2)^2$\\
  Delzell & 4 & 8 & $D(x_1,\ldots,x_4)=x_1^4x_2^2x_4^2+x_2^4x_3^2x_4^2+x_1^2x_3^4x_4^2-3x_1^2x_2^2x_3^2x_4^2+x_3^{8}$\\
  Stengle-$k$ & 3 &$4k+2$  & $S_k(x_1,x_2,x_3)=x_1^{2k+1}x_3^{2k+1}+(x_2^2x_3^{2k-1}-x_1^{2k+1}-x_1x_3^{2k})^2$\\
  \bottomrule
  \end{tabular}
\end{table}

We record the number of subregions generated by Algorithm~\ref{alg:bnb} with $\epsilon_{\text{tol}}=10^{-4}$, using a single projected gradient descent step in~\eqref{eq:pgd_poly} for computing an upper bound at each node of the branch-and-bound tree. We also experiment with the two types of initialization described in Section~\ref{sec:adaptive_grid}. Table~\ref{tab:num_poly} shows that with both initializations, the algorithm terminates with a small number of subregions in most cases. 

% successfully converges for all test cases, requiring only a small number of subregions in most cases. 

\begin{table}
  \centering
  \caption{The number of subregions produced by Algorithm~\ref{alg:bnb} with two types of initialization described in Section~\ref{sec:adaptive_grid}.}
  \label{tab:num_poly}
  \small
  \begin{tabular}{@{}c@{\hspace{1.5em}}c@{}}
  \begin{minipage}[t]{0.45\textwidth}
    \centering
    \begin{tabular}{lll}
      \toprule
      & \multicolumn{2}{c}{\# subregions} \\
      \multirow{-2}{*}{Name} & Init 1 & Init 2 \\\midrule
      Motzkin    & 4   & 7   \\
      Robinson-1 & 4   & 8   \\
      Robinson-2 & 8   & 19  \\
      Choi-Lam-1 & 5   & 15  \\
      Choi-Lam-2 & 4   & 8   \\
      % Horn       & 22  & 102 \\
      Lax        & 98 & 149 \\
      Schm\"udgen & 4 & 5 \\\bottomrule
    \end{tabular}
  \end{minipage}
  &
  \begin{minipage}[t]{0.45\textwidth}
    \centering
    \begin{tabular}{lll}
      \toprule
      & \multicolumn{2}{c}{\# subregions} \\
      \multirow{-2}{*}{Name} & Init 1 & Init 2 \\\midrule
          \textsc{Partition} & 69 & 161 \\
        Delzell & 8 & 5 \\
      Stengle-1   & 4 & 10 \\
      Stengle-2   & 4 & 4 \\
      Stengle-3   & 4 & 4 \\
      Stengle-4   & 4 & 4 \\
      Stengle-5   & 4 & 4 \\
       \bottomrule
    \end{tabular}
  \end{minipage}
  \end{tabular}

    \vspace{.5em}
  \begin{minipage}{0.9\textwidth}
    \small \textit{Remark:} The number of subregions can be lowered, sometimes significantly, by a better implementation of the branch-and-bound algorithm. See \url{https://github.com/YH7422/DisjunctiveSOS} for heuristic improvements to our implementation made by Claude Code, which result, e.g., in 64 (Init 1) and 49 (Init 2) regions for the Lax polynomial, and 32 (Init 1) and 35 (Init 2) regions for the \textsc{Partition} polynomial.      
  \end{minipage}
\end{table}

% {\color{blue}
% For five of these fourteen polynomials (Motzkin, Choi-Lam-1, Choi-Lam-2, Delzell, and Stengle-1), a complementary view is available: a rational disjunctive sos certificate of the form $p = s_1 + h\,s_2 = s_3 - h\,s_4$ with $h(x) = x_i x_j$ is found by the alternating maximization algorithm of \Sec\ref{sec:disos} and verified exactly in $\mathbb{Q}[x]$, with each $s_i$ certified as SOS by a rational $LDL^\tpose$ factorization of its Gram matrix.
% We report each certificate explicitly in \Sec\ref{app:ai}.
% For the remaining nine polynomials, no rational disjunctive sos certificate of the above form is available with the bases tried; the spatial branch-and-bound algorithm of \Sec\ref{sec:bnb} constructs an adaptive partition of $\reals^n$ that succeeds on every test polynomial.
% }

% \begin{table}
%   \centering
%   \caption{The number of iterations to certify nonnegativity of polynomials using Algorithm~\ref{alg:bnb} with two types of initialization described in Section~\ref{sec:adaptive_grid}.}
%       \label{tab:num_poly}
%   \begin{tabular}{lllllll}
%     \toprule
%     & \multicolumn{2}{c}{\# subregions} & & \multicolumn{2}{c}{\# subregions} \\
    
%     \multirow{-2}{*}{Name} & Init 1 & Init 2 &\multirow{-2}{*}{Name} & Init 1 & Init 2 \\\midrule
%     Motzkin       & 4   & 8   & Schm\"udgen & 4 & 4 \\
%     Robinson-1    & 4   & 8   & Stengle-1   & 4 & 5 \\
%     Robinson-2    & 8   & 25  & Stengle-2   & 4 & 4 \\
%     Choi-Lam-1    & 5   & 15  & Stengle-3   & 4 & 4 \\
%     Choi-Lam-2    & 4   & 8   & Stengle-4   & 4 & 4 \\
%     Horn          & 22  & 102 & Stengle-5   & 4 & 5 \\
%     Lax           & 239 & 247 & Stengle-6   & 4 & 7 \\\bottomrule
%   \end{tabular}
% \end{table}

% \begin{table}
%   \centering
%   \begin{tabular}{llllll}
%       \toprule
% Name & Init 1  & Init 2 & Name & Init 1  & Init 2\\\midrule
% Motzkin &  4 & 8&  Schm\"udgen &4 & 4\\
% Robinson-1 & 4 & 8 & Stengle-1 &4&5 \\
% Robinson-2 & 8 & 25 & Stengle-2 &4& 4\\
%   Choi-Lam-1 & 5 & 15 & Stengle-3 &4&4 \\
%   Choi-Lam-2 & 4 & 8 & Stengle-4 &4&4\\
%   Horn & 22& 102& Stengle-5 &4&5\\
%   Lax & 239&247 &Stengle-6&4&7\\\bottomrule

% \end{tabular}
% \end{table}

Figures~\ref{fig:poly_R2_PT_init1} and~\ref{fig:poly_R2_PT_init0} show the lower and upper bounds on the minimum values of the Robinson-2 polynomial and the \textsc{Partition} polynomial versus the number of subregions produced by Algorithm~\ref{alg:bnb}.
% for the Robinson-2 polynomial and the \textsc{Partition} polynomial.

\begin{figure}[H]
\centering
\includegraphics[width=0.45\textwidth]{figs/Robinson2_bnb_sd_init0_new.pdf}\hfill
\includegraphics[width=0.45\textwidth]{figs/Partition_bnb_sd_init0_new.pdf}
    \caption{Lower and upper bounds on the minimum value of a given polynomial versus the number of subregions produced by Algorithm~\ref{alg:bnb} with the first type of initialization.
    Left: Robinson-2 polynomial. Right: \textsc{Partition} polynomial.}
\label{fig:poly_R2_PT_init1}
\end{figure}

\begin{figure}[H]
\centering
\includegraphics[width=0.44\textwidth]{figs/Robinson2_bnb_sd_init1_new.pdf}\hfill
\includegraphics[width=0.45\textwidth]{figs/Partition_bnb_sd_init1_new.pdf}
    \caption{Lower and upper bounds on the minimum value of a given polynomial versus the number of subregions produced by Algorithm~\ref{alg:bnb} with the second type of initialization.
    Left: Robinson-2 polynomial. Right: \textsc{Partition} polynomial.}
\label{fig:poly_R2_PT_init0}
\end{figure}

% \[\begin{aligned}
%     D(x_1^2,\ldots,x_4^2)&=x_1^8x_2^4x_4^4+x_2^8x_3^4x_4^4+x_1^4x_3^8x_4^4-3x_1^4x_2^4x_3^4x_4^4+x_3^{16}\\
%     &=\frac{1}{9}h_1^2+\frac{2}{3}h_2^2+\frac{2}{3}h_3^2+\frac{2}{3}h_4^2+\frac{2}{3}h_5^2+\frac{2}{9}h_6^2+h_7^2,
% \end{aligned}\]
% where
% \[\begin{aligned}
% h_1&=x_1^4x_2^2x_4+x_2^4x_3^2x_4+x_1^2x_3^4x_4-3x_1^2x_2^2x_3^2x_4\\  
% h_2&=x_1x_2^2x_3^3x_4-x_1^3x_2^2x_3x_4\\  
% h_3&=x_1^2x_2x_3^3x_4-x_1^2x_2^3x_3x_4\\  
% h_4&=x_1x_2^3x_3^2x_4-x_1^3x_2x_3^2x_4\\  
% h_5&=x_1^4x_2^2x_4-x_1^2x_3^4x_4\\  
% h_6&=x_1^4x_2^2x_4-2x_2^4x_3^2x_4+x_1^2x_3^4x_4\\  
% h_7&=x_3^8.  
% \end{aligned}\]

% \begin{figure}[H]
% \centering
% % \begin{subfigure}{\textwidth}
% % \includegraphics[width=0.495\textwidth]{figs/Robinson1_bnb2_scatter.pdf}\hfill
% \includegraphics[width=0.5\textwidth]{figs/Robinson2_bnb_sd_init1.pdf}
% % \end{subfigure}
%     \caption{Lower and upper bounds on the optimal value obtained by Algorithm~\ref{alg:bnb} versus the number of iterations for Robinson-2 polynomials, using the first type of initialization. The Robinson-1 plot is omitted since the bounds already match after initialization.}
% \label{fig:adaptive_R1_R2_init1}
% \end{figure}

% \begin{figure}[H]
% \centering
% \includegraphics[width=0.49\textwidth]{figs/Robinson1_bnb_sd_init0.pdf}\hfill
% \includegraphics[width=0.51\textwidth]{figs/Choi-Lam1_bnb_sd_init0.pdf}
%     \caption{Lower and upper bounds on the optimal value obtained by Algorithm~\ref{alg:bnb} versus the number of iterations for the Robinson-1 (left) and Choi-Lam-1 (right) polynomials, using the second type of initialization. The plots using the first type of initialization are omitted since the bounds already match after initialization.}
% \label{fig:adaptive_R1_R2_init0}
% \end{figure}

% \begin{figure}[H]
% \centering
% \includegraphics[width=0.49\textwidth]{figs/Robinson2_bnb_sd_init1.pdf}\hfill
% \includegraphics[width=0.51\textwidth]{figs/Lax_bnb_sd_init1.pdf}
%     \caption{Lower and upper bounds on the optimal value obtained by Algorithm~\ref{alg:bnb} versus the number of iterations for the Robinson-2 (left) and Lax (right) polynomials, using the first type of initialization.}
% \label{fig:adaptive_H_A_init1}
% \end{figure}

% \begin{figure}[H]
% \centering
% \includegraphics[width=0.5\textwidth]{figs/Robinson2_bnb_sd_init0.pdf}\hfill
% \includegraphics[width=0.5\textwidth]{figs/Lax_bnb_sd_init0.pdf}
%     \caption{Lower and upper bounds on the optimal value obtained by Algorithm~\ref{alg:bnb} versus the number of iterations for the Robinson-2 (left) and Lax (right) polynomials, using the second type of initialization.}
% \label{fig:adaptive_H_A_init0}
% \end{figure}

\subsubsection{Applications to Copositive Programming}\label{sec:num_cop}


We now present some numerical experiments with Algorithm~\ref{alg:bnb_simplex}, which can solve copositive programs or the closely related problem of minimizing a quadratic form over the unit simplex.


\paragraph{Solving nonconvex quadratic \textcolor{black}{programs} (QPs).}



% Recall that $\textcolor{black}{\Delta^{n-1}}$ denotes the unit simplex in $\reals^n$ and $\mcC_n$ denotes the cone of $n\times n$ copositive matrices. It is easy to establish that the optimal value of the QP $\min_{x\in\textcolor{black}{\Delta^{n-1}}}\; x^TQx$ is equal to the optimal value of the following copositive program:
% % The QP $\min_{x\in\textcolor{black}{\Delta^{n-1}}}\; x^TQx$
% % is equivalent to the copositive program (see, e.g.,~\cite{Bomze2000}):
% \[            \begin{array}{cl}
%         \max\limits_{t\in\reals} & t \\
%          \st & Q-tJ_n\in \mcC_n.
%     \end{array}\]
We apply Algorithm~\ref{alg:bnb_simplex} to four nonconvex QP instances of the form $\min_{x\in\textcolor{black}{\Delta^{n-1}}}\; x^TQx$ taken from~\cite{Bundfuss2009}.
% We take four matrices from~\cite{Bundfuss2009}.
The four matrices defining the objective functions of these QPs are as follows:
% Recall that $\mcC_n$ denotes the set of positive matrices, and that $\mcC_n(\mcV)$, defined in Definition~\ref{def:copos_disos}, is an inner approximation of $\mcC_n$. Consequently, we apply Algorithm~\ref{alg:bnb_simplex} to solve the standard QP. 
% We test four benchmark instances from~\cite{Bundfuss2009}, with matrices listed below.
\allowdisplaybreaks
\begin{alignat*}{4}
Q_1 &= \begin{bmatrix}
1 & 0 & 1 & 1 & 0 \\
0 & 1 & 0 & 1 & 1 \\
1 & 0 & 1 & 0 & 1 \\
1 & 1 & 0 & 1 & 0 \\
0 & 1 & 1 & 0 & 1
\end{bmatrix},
&\quad
Q_2 &= {\footnotesize \left[\begin{array}{llllllllllll}
1 & 0 & 0 & 0 & 0 & 0 & 1 & 1 & 1 & 1 & 1 & 1 \\
0 & 1 & 0 & 0 & 1 & 1 & 0 & 0 & 1 & 1 & 1 & 1 \\
0 & 0 & 1 & 1 & 0 & 1 & 0 & 1 & 0 & 1 & 1 & 1 \\
0 & 0 & 1 & 1 & 1 & 0 & 1 & 0 & 1 & 0 & 1 & 1 \\
0 & 1 & 0 & 1 & 1 & 0 & 1 & 1 & 0 & 1 & 0 & 1 \\
0 & 1 & 1 & 0 & 0 & 1 & 1 & 1 & 1 & 0 & 0 & 1 \\
1 & 0 & 0 & 1 & 1 & 1 & 1 & 0 & 0 & 1 & 1 & 0 \\
1 & 0 & 1 & 0 & 1 & 1 & 0 & 1 & 1 & 0 & 1 & 0 \\
1 & 1 & 0 & 1 & 0 & 1 & 0 & 1 & 1 & 1 & 0 & 0 \\
1 & 1 & 1 & 0 & 1 & 0 & 1 & 0 & 1 & 1 & 0 & 0 \\
1 & 1 & 1 & 1 & 0 & 0 & 1 & 1 & 0 & 0 & 1 & 0 \\
1 & 1 & 1 & 1 & 1 & 1 & 0 & 0 & 0 & 0 & 0 & 1
\end{array}\right]},
\\[1em]
{Q_3} &= {\footnotesize \begin{bmatrix}
-14 & -15 & -16 & 0 & 0 \\
-15 & -14 & -12.5 & -22.5 & -15 \\
-16 & -12.5 & -10 & -26.5 & -16 \\
0 & -22.5 & -26.5 & 0 & 0 \\
0 & -15 & -16 & 0 & -14
\end{bmatrix},}
&\quad
{ Q_4} &= {\footnotesize \begin{bmatrix}
0.9044 & 0.1054 & 0.5140 & 0.3322 & 0 \\
0.1054 & 0.8715 & 0.7385 & 0.5866 & 0.9751 \\
0.5140 & 0.7385 & 0.6936 & 0.5368 & 0.8086 \\
0.3322 & 0.5866 & 0.5368 & 0.5633 & 0.7478 \\
0 & 0.9751 & 0.8086 & 0.7478 & 1.2932
\end{bmatrix}.}
\end{alignat*}

% \[Q_1=\begin{bmatrix}
% 1 & 0 & 1 & 1 & 0 \\
% 0 & 1 & 0 & 1 & 1 \\
% 1 & 0 & 1 & 0 & 1 \\
% 1 & 1 & 0 & 1 & 0 \\
% 0 & 1 & 1 & 0 & 1
% \end{bmatrix},\quad {\footnotesize Q_2=\left[\begin{array}{llllllllllll}
% 1 & 0 & 0 & 0 & 0 & 0 & 1 & 1 & 1 & 1 & 1 & 1 \\
% 0 & 1 & 0 & 0 & 1 & 1 & 0 & 0 & 1 & 1 & 1 & 1 \\
% 0 & 0 & 1 & 1 & 0 & 1 & 0 & 1 & 0 & 1 & 1 & 1 \\
% 0 & 0 & 1 & 1 & 1 & 0 & 1 & 0 & 1 & 0 & 1 & 1 \\
% 0 & 1 & 0 & 1 & 1 & 0 & 1 & 1 & 0 & 1 & 0 & 1 \\
% 0 & 1 & 1 & 0 & 0 & 1 & 1 & 1 & 1 & 0 & 0 & 1 \\
% 1 & 0 & 0 & 1 & 1 & 1 & 1 & 0 & 0 & 1 & 1 & 0 \\
% 1 & 0 & 1 & 0 & 1 & 1 & 0 & 1 & 1 & 0 & 1 & 0 \\
% 1 & 1 & 0 & 1 & 0 & 1 & 0 & 1 & 1 & 1 & 0 & 0 \\
% 1 & 1 & 1 & 0 & 1 & 0 & 1 & 0 & 1 & 1 & 0 & 0 \\
% 1 & 1 & 1 & 1 & 0 & 0 & 1 & 1 & 0 & 0 & 1 & 0 \\
% 1 & 1 & 1 & 1 & 1 & 1 & 0 & 0 & 0 & 0 & 0 & 1
% \end{array}\right]},\]
% \[{\footnotesize Q_3=\begin{bmatrix}
% -14 & -15 & -16 & 0 & 0 \\
% -15 & -14 & -12.5 & -22.5 & -15 \\
% -16 & -12.5 & -10 & -26.5 & -16 \\
% 0 & -22.5 & -26.5 & 0 & 0 \\
% 0 & -15 & -16 & 0 & -14
% \end{bmatrix},\quad Q_4=\begin{bmatrix}
% 0.9044 & 0.1054 & 0.5140 & 0.3322 & 0 \\
% 0.1054 & 0.8715 & 0.7385 & 0.5866 & 0.9751 \\
% 0.5140 & 0.7385 & 0.6936 & 0.5368 & 0.8086 \\
% 0.3322 & 0.5866 & 0.5368 & 0.5633 & 0.7478 \\
% 0 & 0.9751 & 0.8086 & 0.7478 & 1.2932
% \end{bmatrix}}.\]
We record the number of subregions generated by Algorithm~\ref{alg:bnb_simplex} with $\epsilon_{\text{tol}}=10^{-6}$. 
We perform 5 projected gradient descent steps in~\eqref{eq:pgd_copos} to obtain upper bounds at each node of the branch-and-bound tree. 
% Since the objective is quadratic, each projected gradient descent step is inexpensive and we perform 5 steps to obtain upper bounds.
Table~\ref{tab:num_cop} compares our subregion counts against those reported in~\cite{Bundfuss2009} using an LP-based adaptive approximation algorithm. 
% Our algorithm uses fewer subregions on all four instances. 

% As an illustrative example, Figure~\ref{fig:copositive_Q2} plots the lower and upper bounds on the optimal value versus the number of iterations for $Q_2$, reflecting the tightness of our lower bounds.

\begin{table}
  \centering
    \caption{The number of subregions produced by Algorithm~\ref{alg:bnb_simplex} to solve nonconvex QPs.}
      \label{tab:num_cop}
      \small
  \begin{tabular}{llllll}
    \toprule
 Instance & $Q_1$ & $Q_2$ & $Q_3$ & $Q_4$\\
    \midrule
    Adaptive Linear Approximation~\cite{Bundfuss2009} & 6 & 158& 44 & 27\\
    Algorithm~\ref{alg:bnb_simplex}& 2&42 & 5&17\\\bottomrule
  \end{tabular}
\end{table}

Figure~\ref{fig:num_copositive} shows the lower and upper bounds on the optimal value of the four QPs versus the number of subregions.

\begin{center}
\includegraphics[width=0.45\textwidth]{figs/copositive1_bnb.pdf}\hfill
\includegraphics[width=0.45\textwidth]{figs/copositive2_bnb.pdf}\\
\includegraphics[width=0.45\textwidth]{figs/copositive3_bnb.pdf}\hfill
\includegraphics[width=0.44\textwidth]{figs/copositive4_bnb.pdf}
\captionof{figure}{Lower and upper bounds on the optimal value of the QP $\min_{x\in\textcolor{black}{\Delta^{n-1}}}\; x^TQx$ produced by Algorithm~\ref{alg:bnb_simplex} versus the number of subregions. Top left: $Q=Q_1$. Top right: $Q=Q_2$. Bottom left: $Q=Q_3$. Bottom right: $Q=Q_4$.}
\label{fig:num_copositive}
\end{center}

% \begin{figure}[H]
% \centering
% \includegraphics[width=0.75\textwidth]{figs/copositive2_bnb_new.pdf}
%     \caption{Lower and upper bounds on the optimal value of the standard QP with instance $Q_2$ obtained by Algorithm~\ref{alg:bnb_simplex} versus the number of iterations.}
% \label{fig:copositive_Q2}
% \end{figure}

\paragraph{Computing the clique number of a graph.}
A clique in a graph is a set of pairwise adjacent vertices. The clique number of a graph $G$, denoted by $\omega(G)$, is the size of a maximum
clique in $G$. 
% The clique number of a graph $G$, denoted by $\omega(G)$, is the size of the maximum clique (complete subgraph) in $G$. 
By a theorem of Motzkin and Straus~\cite{Motzkin1965}, we have
% $\omega(G)$ is related to the optimal value of a QP via
% can be obtained by solving a QP of the form
\[\frac{1}{\omega(G)}=\min\limits_{x\in\textcolor{black}{\Delta^{n-1}}} x^T(I_n+\bar{A})x,\]
where $\bar{A}$ is the adjacency matrix of the complement graph $\bar{G}$; i.e., the graph obtained by swapping the edges and non-edges of $G$. Equivalently, the clique number can be written as the optimal value of the copositive program
\begin{equation}\label{eq:clique}
\begin{aligned}
    \omega(G)=\ &\min_{k\in\reals}\ k\\
    &\ \st\ k(I_n+\bar{A}) - J_n\in\mcC_n.
\end{aligned}    
\end{equation}
% to which we apply Algorithm~\ref{alg:bnb_simplex}. 
To estimate the clique number by Algorithm~\ref{alg:bnb_simplex}, at each node of the branch-and-bound tree, we perform 10 projected gradient descent steps in~\eqref{eq:pgd_copos} with $Q=I_n+\bar{A}$ to obtain a lower bound on $\omega(G)$, and apply our inner approximation of $\mcC_n$ to~\eqref{eq:clique} to obtain an upper bound on $\omega(G)$. 
Since the clique number is an integer, the algorithm terminates when the ceiled lower bound and the floored upper bound coincide.

% where, before checking convergence, we round the ceil the lower bound and floor the upper bounds. xxxx to clarify
We present some experiments on Erd\"os-R\'enyi random graphs~\cite{Erdos1959} $G(n,p)$, where $n$ vertices are connected pairwise and independently with probability $p\in(0,1)$. 
We generate four random instances with $n=75$ and $p=0.5$, and compare our bounds against the Schrijver upper bound $\vartheta'(\bar{G})$ on $\omega(G)$~\cite{Schrijver1979}.\footnote{The Schrijver bound $\vartheta'(\bar{G})$ satisfies $\omega(G)\leq \vartheta'(\bar{G}) \leq \vartheta(\bar{G})$, where $\vartheta$ is the more familiar Lov\'asz theta function~\cite{Lovasz1979}.} 
% Since clique numbers are integers, the algorithm terminates when the lower and upper bounds coincide after rounding. 
% Since the problem dimension is higher than in the previous experiments and each projected gradient descent step is inexpensive, we perform 10 steps in Algorithm~\ref{alg:bnb_simplex}.
Table~\ref{tab:clique_random} shows that our algorithm successfully determines the exact clique number in all four instances using a moderate number of subregions.

\begin{table}
\centering
\caption{Lower and upper bounds on the clique number, and the number of subregions produced by Algorithm~\ref{alg:bnb_simplex} for four random instances of $G(75, 0.5)$.}
\label{tab:clique_random}
\small
\begin{tabular}{cccccc}
\toprule
Instance & $\omega(G)$ & $\vartheta'(\bar{G})$ & Upper bound & Lower bound & \# subregions \\
\midrule
1 & 8 & 9.0336 & 8.9014 & 7.9956 & 8 \\
2 & 8 & 9.1274 & 8.9815 & 7.9967 & 48 \\
3 & 8 & 8.9770 & 8.8686 & 7.9950 & 4 \\
4 & 8 & 9.1400 & 8.9952 & 7.9973 & 16 \\
\bottomrule
\end{tabular}
\end{table}

% \begin{table}
% \centering
% \caption{Lower and upper bounds on the clique number, and the number of subregions produced by Algorithm~\ref{alg:bnb_simplex} for four random instances of $G(75, 0.5)$.}
% \label{tab:clique_random}
% \small
% \begin{tabular}{cccccc}
% \toprule
% Instance & $\omega(G)$ & $\vartheta'(\bar{G})$ & Lower bound & Upper bound & \# subregions \\
% \midrule
% 1 &8 &9.0336 &7.9956 & 8.9014&8 \\
% 2 &8 &9.1274 &7.9967& 8.9815& 48 \\
% 3 &8 &8.9770 &7.9950 &8.8686 &4 \\
% 4 & 8&9.1400 &7.9973 &8.9952 & 16\\
% \bottomrule
% \end{tabular}
% \end{table}
Figure~\ref{fig:num_clique} shows the lower and upper bounds, produced by Algorithm~\ref{alg:bnb_simplex}, on the clique number of the four random instances of $G(75,0.5)$ versus the number of subregions.

\begin{center}
\includegraphics[width=0.45\textwidth]{figs/Erdos_n75_p0.5_test1_bnb.pdf}\hfill
\includegraphics[width=0.45\textwidth]{figs/Erdos_n75_p0.5_test2_bnb.pdf}\\[-0.5ex]
\includegraphics[width=0.45\textwidth]{figs/Erdos_n75_p0.5_test3_bnb.pdf}\hfill
\includegraphics[width=0.45\textwidth]{figs/Erdos_n75_p0.5_test4_bnb.pdf}
\captionof{figure}{Lower and upper bounds on the clique numbers of four random instances of $G(75, 0.5)$ produced by Algorithm~\ref{alg:bnb_simplex} versus the number of subregions.}
\label{fig:num_clique}
\end{center}

% \begin{figure}[H]
% \centering
% \includegraphics[width=0.49\textwidth]{figs/Erdos_n75_p0.5_test1_bnb.pdf}\hfill
% \includegraphics[width=0.49\textwidth]{figs/Erdos_n75_p0.5_test2_bnb.pdf}\\
% \includegraphics[width=0.49\textwidth]{figs/Erdos_n75_p0.5_test3_bnb.pdf}\hfill
% \includegraphics[width=0.49\textwidth]{figs/Erdos_n75_p0.5_test4_bnb.pdf}
% \caption{Lower and upper bounds on the clique number $\omega(G)$ obtained by Algorithm~\ref{alg:bnb_simplex} versus the number of iterations for four random instances of $G(75, 0.5)$. Top left: instance 1. Top right: instance 2. Bottom left: instance 3. Bottom right: instance 4.}
% \label{fig:num_clique}
% \end{figure}

% The Schrijver theta number~\cite{Schrijver1979}, denoted as $\vartheta'(G)$, typically provides a good upper bound for $\omega(G)$, by solving the following SDP:
% \[\begin{array}{rll}
% \vartheta'(G):=&\max\limits_{X\in S^{n\times n}} & \operatorname{Tr}(J X) \\
% & \text { s.t. } & X_{i j}=0 \quad \text { if } (i, j) \in E \\
% && \operatorname{Tr}(X)=1 \\
% && X \succeq 0 \\
% && X \geq 0 ,
% \end{array}\]
% where $J\in\reals^{n\times n}$ is the all-one matrix.

% We conduct experiments on Paley graphs $P_q$, constructed on the finite field $\mbF_q$. Two vertices $x,y\in\{0,1,\ldots,q-1\}$ are connected if and only if $x-y=a^2$ for some $a\in \mbF_q$. As a benchmark, the Schrijver theta number $\vartheta'(P_q)$~\cite{Schrijver1979} provides an upper bound on $\omega(P_q)$ and satisfies 
% $\vartheta'(P_q)=\sqrt{q}$ when the prime number $q\equiv 1\pmod{4}$~\cite{godsil2016sabidussi}. 
% However,~\cite{Hanson2021} show that \[\omega(P_q)\le \frac{\sqrt{2q-1}+1}{2},\]
% which implies that the gap between the true clique number and the Schrijver theta number can be substantial for this family of graphs.
% % Therefore, there can be a significant gap between the actual clique number of this family of graphs and the upper bound provided by the Schrijver theta number. 

% We evaluate the Paley graphs with $q=17,29,37,41,61$ and record the lower and upper bounds produced by Algorithm~\ref{alg:bnb_simplex} in Table~\ref{tab:num_clique}. The lower bounds are displayed as integers since the returned values agree with their integer part to at least four decimal places. Since clique numbers are integers, the algorithm terminates when the lower and upper bounds coincide after rounding, or after a prescribed maximum number of iterations. For $q=17,29,41$, the algorithm successfully determines the exact value of $\omega(P_q)$. For $q=37$ and $q=61$, where the maximum iteration limit is reached, the upper bounds obtained by our method still improve considerably upon the Schrijver theta number.

% % demonstrating that the disjunctive copositive approach can yield significantly sharper bounds on this challenging benchmark.
% % The termination criterion is based on the integral nature of the clique number: either the lower and upper bound match after rounding, or the maximum number of iterations is reached. 
% % For the same reason, we also include the rounded result in Table~\ref{tab:clique_disos_paley}. Although our algorithm do not completely solve for $q=37,41$ in 1000 iterations, the upper bound obtained by our algorithm achieves a reasonable improvement compared to the Schrijver theta number. This result demonstrates the strength of our algorithm.

% \begin{table}
%     \centering
%         \caption{Lower and upper bounds on the clique number $\omega(P_q)$ obtained by Algorithm~\ref{alg:bnb_simplex}.}
%     \begin{tabular}{lllllll}
%     \toprule
%     $q$&$\omega(P_q)$&$\vartheta'(P_q)$&$(\sqrt{2q-1}+1)/{2}$&Upper bound& Lower bound& Iterations\\\midrule
% 17  & 3&4.1231&3.3723&	3.9033&3&4\\
% 29  &4&5.3852&4.2749&4.9968&4&88	\\
% 37  &4&6.0828&4.772&5.4007&4&1000(MAX)	\\
% 41  &5&6.4031&5&5.9997&5&462\\
% 61&5&7.8102&6&7.3180&5&1000(MAX)\\\bottomrule
%     \end{tabular}
%     \label{tab:num_clique}
% \end{table}




\section{Conclusions and Future Directions}\label{sec:discussion}
In this paper, we introduced the disjunctive sum of squares method for certifying nonnegativity of polynomials.
Unlike the classical sum of squares approach, which certifies nonnegativity via a single algebraic identity, our method constructs multiple algebraic identities, which together certify nonnegativity. 
We showed that by doing so, the degrees of the sum of squares polynomials involved in the proof (and hence the size of the underlying semidefinite constraints) can be kept low and fixed. We proposed algorithms that can search for disjunctive proofs of nonnegativity, and studied extensions of our method to copositive programs and polynomial optimization problems.


% This keeps the degree of the sum of squares polynomials fixed and low throughout, with the size of the underlying semidefinite constraint remaining fixed across all levels of the hierarchy.
% We prove that one can a priori fix a family of low-degree algebraic disjunctions with respect to which every positive definite form admits a disjunctive sum of squares proof of nonnegativity.
% We also extend the method to copositive programs and general polynomial optimization problems over compact semialgebraic sets.
{
This work leaves several interesting directions for future research. While we proposed a general reduction from constrained polynomial optimization problems to unconstrained ones in Section~\ref{sec:hierarchy}, future work could tailor the design of algebraic disjunctions to the geometry of specific feasible sets (as we did, e.g., in the case of the nonnegative orthant in Section~\ref{sec:copositive}). 
%consider the geometry of specific feasible sets to tailor 
%In future work, we aim to develop disjunctive sum of squares based algorithms that tailor to the geometry of the feasible region of interest, potentially yielding more scalable and problem-specific tools for global optimization.
Another interesting direction is to develop more efficient implementations of the disjunctive sum of squares approach that in addition to utilizing spatial {branch-and-bound} (cf.\ Section~\ref{sec:bnb}), potentially exploit the natural parallelism across subregions. 
% A natural next step is to conduct more extensive numerical experiments on both constrained and unconstrained polynomial optimization problems, supported by a more efficient implementation, potentially with parallelism across subregions.
Along these lines, it would be interesting to combine the disjunctive sum of squares method with techniques for exploiting sparsity~\cite{Waki2006, Zheng2022, Zheng2021} and symmetry~\cite{Gatermann2004, Riener2013} in sum of squares programs. Another avenue to study is the dual side of our framework through the theory of moments~\cite{lasserre2001global, Lasserre2019}, which has proven useful for extracting solutions in polynomial optimization problems~\cite{Henrion2005}.
Finally, a promising direction is the design of instance-specific algebraic disjunctions that require fewer subregions than the universal constructions provided in this paper (cf.\ Theorems~\ref{cor:strictly_pos} and~\ref{cor:strictly_pos_new}). The branch-and-bound framework of Section~\ref{sec:bnb} constructs such disjunctions adaptively, but there could be techniques that in addition exploit algebraic properties of the given problem instance.}



{We end with an open problem of theoretical interest: does there exist a nonnegative form that admits no disjunctive sum of squares proof of nonnegativity with any number of identities? This question has two natural variants: one where the degree of the sum of squares proof is required to be the same as {that of} the given form, and one where it is allowed to be larger. By Theorem~\ref{cor:strictly_pos}, any such form must have a zero.
A natural candidate would be a polynomial with a {``bad point''}; see, e.g.,~\cite{Reznick2005, Benoist2021} and references therein. 
% : given a polynomial $p\in P_{n,d}$, a point $x^*$ is said to be a bad point if for every polynomial $q$ with $q(x^*)\neq 0$, $qf$ is not sos. 
An example of such a polynomial is Delzell's polynomial $D$ (see~\cite{Delzell1980})
$$D(x_1,\ldots,x_4)=x_1^4x_2^2x_4^2+x_2^4x_3^2x_4^2+x_1^2x_3^4x_4^2-3x_1^2x_2^2x_3^2x_4^2+x_3^{8},$$
which was also a test case for us in Section~\ref{sec:num_poly}. It is known that while $D$ is nonnegative, $\left(\sum_{i=1}^4 x_i^2\right)^k D(x_1,\ldots,x_4)$ is not sos for any integer $k \ge 0$, which implies that the hierarchy of SDPs based on Reznick's Positivstellensatz~\cite{reznick1995uniform} cannot certify nonnegativity of~$D$. 
It turns out, however, that the disjunctive sum of squares approach can produce a low-degree proof of nonnegativity of~$D$ only with two subregions: 
%One might hope that this obstruction carries over to disjunctive sum of squares certificates, but this turns out not to be the case:
\ifpreprint
\[
\begin{aligned}
D(x_1,\ldots,x_4) &= (x_1^2x_2x_4 - x_1x_3^2x_4)^2 + (x_2^2x_3x_4 - x_1x_2x_3x_4)^2 + (x_3^4)^2 + 2x_1x_2(x_1x_3x_4-x_2x_3x_4)^2\\
    & = (x_1^2x_2x_4 + x_1x_3^2x_4)^2 + (x_2^2x_3x_4 + x_1x_2x_3x_4)^2 + (x_3^4)^2 - 2x_1x_2(x_1x_3x_4+x_2x_3x_4)^2.
    \end{aligned}
\]
\else
\[
\begin{aligned}
D(x_1,\ldots,x_4) ={}& (x_1^2x_2x_4 - x_1x_3^2x_4)^2 + (x_2^2x_3x_4 - x_1x_2x_3x_4)^2 + (x_3^4)^2\\
		     &\hspace{3em} + 2x_1x_2(x_1x_3x_4-x_2x_3x_4)^2\\
    ={}& (x_1^2x_2x_4 + x_1x_3^2x_4)^2 + (x_2^2x_3x_4 + x_1x_2x_3x_4)^2 + (x_3^4)^2\\
       &\hspace{3em} - 2x_1x_2(x_1x_3x_4+x_2x_3x_4)^2.
    \end{aligned}
\]
\fi
% Whether a form exists that truly escapes all disjunctive sum of squares certificates remains an open question.
}

% In this paper, we introduce the disjunctive sum of squares method for certifying nonnegativity of polynomials. The key idea is to introduce algebraic disjunctions, which use low-degree polynomials to cover $\reals^n$, and to construct algebraic certificates on each subregion. 
% Unlike classical sum of squares hierarchies, which increase the degree of the sos polynomials at each level, our method keeps the degree fixed and low throughout, with the size of the underlying semidefinite constraint remaining fixed across all levels of the hierarchy. 
% We prove that one can a priori fix a family of low-degree algebraic disjunctions with respect to which every positive definite form admits a disjunctive sum of squares proof of nonnegativity.
% We also extend the method to copositive programs and general polynomial optimization problems over compact semialgebraic sets.

% We end by highlighting two open questions for future research.
% Despite these advances, important questions remain open. We highlight two below. 
% \begin{itemize}
%     \item {\bf Optimal algebraic disjunctions.} 
% We end with a direction for future research. 
% The family of algebraic disjunctions we construct for Theorem~\ref{cor:strictly_pos} is universal: it is defined independently of the problem and is guaranteed to work for any positive definite form. In practice, however, instance-specific disjunctions may make it easier to find certificates. 
% Formally, given fixed integers $r, n_1,\ldots,n_r$ as in Definition~\ref{def:disjunction}, a polynomial $p\in H_{n,d}$, and an integer $\bar{d}\ge d$, the optimal algebraic disjunction maximizes
% \[
%     \begin{array}{rl}
%         \gamma^*(\mcD)\defeq \max\limits_{\gamma\in{\reals}} & \gamma\\
%         \st & p(x)-\gamma\|x\|_2^d\in \DiSOS_{n,d}^{\bar{d}}(\mcD),
%     \end{array}
% \]
% over all algebraic disjunctions $\mcD$ whose elements have degree at most $\bar{d}$.
% Finding such a disjunction is a non-convex design problem.
% Section~\ref{sec:disos} introduces an alternating maximization algorithm as a heuristic approach, and developing more principled strategies for constructing an optimal or near-optimal algebraic disjunction is a promising research direction.
    
    % Section~\ref{sec:disos} and Section~\ref{sec:disos_psatz} present several strategies for constructing algebraic disjunctions, but these are not necessarily optimal. For a polynomial $p\in H_{n,d}$ and integer $\bar{d}\ge d$, different algebraic disjunctions of $\reals^n$, even with the same number of subsets, can produce different quality lower bounds for the problem
    % \[
    %     \begin{array}{rl}
    %         \gamma^*(\mcD)\defeq \max\limits_{\gamma\in{\reals}} & \gamma\\
    %         \st & p(x)-\gamma\|x\|_2^d\in \DiSOS_{n,d}^{\bar{d}}(\mcD).
    %     \end{array}
    % \]
    % For a fixed integer $r$, an optimal algebraic disjunction maximizes $\gamma^*(\mcD)$ over all algebraic disjunctions with $r$ subsets. Finding such a disjunction is a non-convex design problem. Developing principled approaches or approximation schemes for constructing near-optimal disjunctions would be a promising research direction.

%     \item 

% {\bf Polynomials without any disjunctive certificate.} For a polynomial $p\in H_{n,d}$, given an algebraic disjunction $\mcD$ and degree $\bar{d}\ge d$, membership in $\DiSOS_{n,d}^{\bar{d}}(\mcD)$ can be checked efficiently via a semidefinite program. In contrast, determining whether $p$ does not belong to $\DiSOS_{n,d}^{\bar{d}}(\mcD)$ for any algebraic disjunction $\mcD$ is substantially harder.
%     We attempted to construct a polynomial with no disjunctive sum of squares proof of nonnegativity, even with $\bar{d}=d$, but were unable to find such an example. Identifying polynomials that escape all disjunctive certificates would provide valuable insights into the limitations of our method.
% \end{itemize}
% In future work, we aim to develop disjunctive sum of squares based algorithms tailored to the geometry of the feasible region, potentially yielding more scalable and problem-specific tools for global optimization. 
% We hope that our results stimulate further investigation into both the capabilities and the limitations of the disjunctive sum of squares method, and inspire new algorithmic developments for challenging polynomial optimization problems.

\section*{Acknowledgments}
We are grateful to Abraar Chaudhry, Georgina Hall, and Jeffrey Zhang for insightful discussions.

\ifpreprint\else
% Springer "Statements and Declarations" (in the preprint, funding appears in the author footnotes).
\section*{Statements and Declarations}
\bmhead{Funding}
A.\,A.\,Ahmadi was partially supported by the Sloan Fellowship, the Princeton AI Lab Seed Grant, the Princeton SEAS Innovation Grant, and a Research Gift in Mathematical Optimization.
Y.\,Hua was partially supported by the Princeton AI Lab Seed Grant, the Princeton SEAS Innovation Grant, and the IBM PhD Fellowship.

\bmhead{Competing Interests}
% TODO: confirm the competing-interests statement with all co-authors before submission.
The authors have no competing interests to declare that are relevant to the content of this article.

\bmhead{Data Availability}
The code and data that support the findings of this study are openly available in the DisjunctiveSOS repository at \url{https://github.com/YH7422/DisjunctiveSOS}.
\fi





%Bibliography
\ifpreprint
\printbibliography
\else
\bibliography{bibliography}
\fi


\end{document}
